\chapter{控制导论}
\label{ch:control}

% ════════════════════════════════════════════════════════════════
%  控制导论：反馈与稳定性(Lyapunov/LaSalle) + PID + 状态空间/可控可观
%  + LQR(Riccati 完整推导) + MPC(滚动时域 + 稳定性证明) + 机器人控制
%  (逆动力学/计算力矩/操作空间/WBC/足式凸 MPC 概览)。
%  个人控制笔记未同步，全章据权威教材/讲义/原论文重建，整体打 \rebuilt。
%  教学层遵循 v5.0 理论教学规范。
% ════════════════════════════════════════════════════════════════

\rebuilt{本章对应的个人控制笔记尚未同步到本仓库，故全章主要据权威教材、课程讲义与原始论文重建（Åström--Murray《Feedback Systems》、Borrelli--Bemporad--Morari 与 Rawlings--Mayne--Diehl 的《Model Predictive Control》、Tedrake《Underactuated Robotics》、Khalil《Nonlinear Systems》、Spong 等《Robot Modeling and Control》，以及 Di Carlo 等 2018、Wensing 等 2024 两篇足式控制论文）。逐处 \texttt{\textbackslash cite} 标源；待个人笔记回填后二次核对。}

\begin{practice}[前置自测]
开始之前，先试答下列问题。若已能流畅作答，可只速读本章表格与速查卡；若卡住，正是本章要为你解决的。
\begin{enumerate}
  \item 一个仅含比例（P）反馈的恒温器为什么总会留下一点稳态温差？加上“积分”项后这点温差为何会消失？你能不假设系统线性就说清楚吗？
  \item 旋转矩阵 $\mathbf{R}\in\SO(3)$ 的“稳定”在数学上如何定义？“轨迹始终留在附近”和“轨迹最终回到平衡点”是同一回事吗？
  \item 给定线性系统 $\dot{\mathbf{x}}=\mathbf{A}\mathbf{x}+\mathbf{B}\mathbf{u}$ 与二次代价 $\int(\mathbf{x}^\top\mathbf{Q}\mathbf{x}+\mathbf{u}^\top\mathbf{R}\mathbf{u})\,\mathrm dt$，你打算如何求“最优”的反馈增益 $\mathbf{K}$？它和一个叫 Riccati 的方程有什么关系？
  \item 把“最优控制律”用在系统上，闭环就一定稳定吗？（这是一个陷阱问题。）
  \item 机器人手臂的动力学 $\mathbf{M}(\mathbf{q})\ddot{\mathbf{q}}+\mathbf{C}\dot{\mathbf{q}}+\mathbf{g}=\boldsymbol{\tau}$ 是强非线性的。有没有办法选一个 $\boldsymbol{\tau}$，让每个关节的跟踪误差服从一个简单的线性二阶方程 $\ddot{\mathbf{e}}+\mathbf{K}_d\dot{\mathbf{e}}+\mathbf{K}_p\mathbf{e}=\mathbf{0}$？
\end{enumerate}
\noindent\emph{答案线索}：1$\to$\cref{sec:ctrl-feedback}（积分作用的“魔力”证明）；2$\to$\cref{sec:ctrl-lyapunov}（Lyapunov/渐近/指数三层定义）；3$\to$\cref{sec:ctrl-lqr}（HJB 与离散 DP 两条推导）；4$\to$\cref{sec:ctrl-lqr} 的 Kalman 反例与 \cref{sec:ctrl-mpc}；5$\to$\cref{sec:ctrl-robot}（计算力矩/反馈线性化）。
\end{practice}

\section*{本章目标}
学完本章，你应当能够：
\begin{enumerate}
  \item \textbf{说清}开环与闭环的根本差异、用时域/积分性能指标量化控制好坏，并独立\textbf{证明}积分作用消除稳态误差（不假设线性）；
  \item \textbf{读出}增益裕度/相位裕度与根轨迹，并\textbf{陈述}并\textbf{运用} Lyapunov 稳定/渐近稳定/指数稳定三层定义、Lyapunov 直接法与 LaSalle 不变原理，判定一个平衡点的稳定性；
  \item \textbf{推导} PID 三项的作用、滤波微分、抗积分饱和与数字实现，并用 Routh--Hurwitz 判 PID 闭环稳定；
  \item \textbf{判别}线性系统的可控性/可观性（秩判据、Gramian、PBH），\textbf{设计}状态反馈与 Luenberger 观测器，并用\textbf{分离原理}组合输出反馈（理解 LQG）；
  \item \textbf{独立推导} LQR：经动态规划（离散）与 HJB（连续）两条路得到 Riccati 方程与最优增益 $\mathbf{K}=\mathbf{R}^{-1}\mathbf{B}^\top\mathbf{P}$，解释“最优 $\ne$ 稳定”为何需要无限时域，并\textbf{推广}到非线性的 iLQR/DDP；
  \item \textbf{表述} MPC 的滚动时域问题、凝聚为 QP，并\textbf{证明}其闭环渐近稳定性（终端代价 + 终端集 = 闭环 Lyapunov 函数）；
  \item \textbf{连接}机器人动力学与控制：计算力矩（反馈线性化）、PD+重力补偿（无源性 + LaSalle）、操作空间与力/阻抗控制、全身控制 QP 与足式凸 MPC 的分层架构。
\end{enumerate}

\subsection*{本章知识导航}
本章是一条“\emph{从经典单回路到现代最优、再到机器人全身控制}”的主线。把六大主题排成一张依赖图：

\begin{center}
\small
\begin{tikzpicture}[
  node distance=4mm and 7mm, >=Stealth,
  blk/.style={draw, rounded corners, align=center, font=\small, inner sep=3pt, fill=main!6, draw=main!60!black},
  grp/.style={draw, rounded corners, align=center, font=\small, inner sep=3pt, fill=second!6, draw=second!60!black},
  adv/.style={draw, rounded corners, align=center, font=\small, inner sep=3pt, fill=third!8, draw=third!60!black},
  ln/.style={->, thick, gray!60!black}]
\node[blk] (fb) {反馈 + 稳定性\\ Lyapunov/LaSalle};
\node[blk, right=of fb] (pid) {PID\\（经典控制）};
\node[grp, below=8mm of fb] (ss) {状态空间\\ 可控/可观};
\node[grp, right=of ss] (lqr) {LQR\\ Riccati};
\node[adv, right=of lqr] (mpc) {MPC\\ 滚动时域 + QP};
\node[adv, below=8mm of ss] (robot) {机器人控制\\ 逆动力学/WBC};
\node[adv, right=of robot] (legged) {足式控制\\ 凸 MPC + 分层};
\draw[ln] (fb)--(pid);
\draw[ln] (fb.south)-- (ss.north);
\draw[ln] (ss)--(lqr); \draw[ln] (lqr)--(mpc);
\draw[ln] (pid.south) -- ++(0,-3mm) -| (lqr.north);
\draw[ln] (ss.south)-- (robot.north);
\draw[ln] (robot)--(legged);
\draw[ln] (lqr.south) -- ++(0,-12mm) -| (legged.north);
\draw[ln] (mpc.south) |- (legged.east);
\end{tikzpicture}
\end{center}

\noindent\textbf{推荐路径（脉络层）}：第 1--2 节（反馈/稳定性、PID）是\emph{骨干}，建立“误差—反馈—稳定”的基本语言，任何读者必读；第 3--4 节（状态空间、LQR）是\emph{骨干}，是现代控制的核心，SLAM 后端的 Riccati、Kalman 增益都与此对偶；第 5 节（MPC）是\emph{次优核心}，给出处理约束的现代框架及其稳定性证明，初读可只看问题表述与主定理结论；第 6 节（机器人控制）是\emph{次优}，把前面工具用到机械臂/足式机器人；第 7 节（足式概览）是\emph{前沿}，以 MIT Cheetah 凸 MPC 为旗舰例，选读。

\subsection*{前置知识桥接}
本章用到三处前置：
\begin{itemize}
  \item \cref{ch:lie} 的旋转 $\mathbf{R}\in\SO(3)$、反对称算子 $(\cdot)^\wedge$、右扰动：第 6 节任务空间姿态误差与第 7 节单刚体动力学的 $\dot{\mathbf{R}}=\boldsymbol{\omega}^\wedge\mathbf{R}$ 都直接复用。回顾要点：$\boldsymbol{\phi}^\wedge$ 把三维向量映成反对称矩阵，$\boldsymbol{\phi}^\wedge\mathbf{y}=\boldsymbol{\phi}\times\mathbf{y}$；$\dot{\mathbf{R}}\mathbf{R}^\top$ 恒为反对称。
  \item \cref{ch:nlopt} 的二次型最小化与正规方程：LQR 的“对 $\mathbf{u}$ 求导置零”与 MPC 的 QP 是同一套最小二乘技术，$\mathbf{H}=\mathbf{J}^\top\mathbf{W}^{-1}\mathbf{J}$ 的 Hessian 思想在此重现。
  \item \cref{ch:eskf} 的 Riccati 递推与协方差传播：LQR 的 Riccati 与 Kalman 滤波的 Riccati 互为\emph{对偶}（一个后向、一个前向），本章 \cref{sec:ctrl-lqe-dual} 会专门点明，并提醒\textbf{字母 $\mathbf{P}$ 的双关}（此处是值函数权重，那里是协方差）。
\end{itemize}

\subsection*{如果跳过本章会怎样}
\begin{itemize}
  \item 在做机器人轨迹跟踪时，你会把 PID 增益盲调到“看起来能动”，却不知道积分饱和为何让系统大幅超调、为何微分项会放大传感器噪声——这些都是本章 \cref{sec:ctrl-pid} 用具体机制解释的；
  \item 在读 VIO/SLAM 后端代码时，你会看到 Riccati 递推（\cref{ch:eskf}）却不知它与最优控制的 Riccati 是同一个数学对象，错失“估计—控制对偶”这条统一主线；
  \item 在用 OSQP/qpOASES 写一个足式机器人的全身控制器时，你会照抄 QP 而不理解“为什么动力学约束对 $(\dot{\boldsymbol{\nu}},\boldsymbol{\tau},\boldsymbol{\lambda})$ 是线性的、为什么摩擦锥要近似成金字塔”——这正是 \cref{sec:ctrl-legged} 要讲清的。
\end{itemize}

\begin{table}[H]\centering\small
\caption{本章阅读方式建议}
\label{tab:ctrl-reading}
\begin{tabular}{lll}
\toprule
阅读方式 & 时间 & 适合谁 \\
\midrule
精读（含全部推导与练习） & 6--8 小时 & 首次系统学控制、要做后端/机器人控制推导 \\
速读（跳过冗长推导细节） & 2--3 小时 & 有控制基础、想对齐本书记号与右扰动衔接 \\
速查（只看小结的符号表 + 定理速查表） & 20 分钟 & 写代码时回来查公式 \\
\bottomrule
\end{tabular}
\end{table}

\begin{note}
\textbf{本章记号与约定（务必先读）.}\;
状态 $\mathbf{x}\in\mathbb{R}^n$、控制输入 $\mathbf{u}\in\mathbb{R}^m$、输出 $\mathbf{y}\in\mathbb{R}^p$。连续 LTI 系统 $\dot{\mathbf{x}}=\mathbf{A}\mathbf{x}+\mathbf{B}\mathbf{u},\ \mathbf{y}=\mathbf{C}\mathbf{x}+\mathbf{D}\mathbf{u}$；离散 $\mathbf{x}_{k+1}=\mathbf{A}\mathbf{x}_k+\mathbf{B}\mathbf{u}_k$。
\begin{itemize}
  \item \textbf{$\mathbf{A},\mathbf{B},\mathbf{C}$ 与估计章的对应}：本章沿用控制界惯例 $\mathbf{A},\mathbf{B},\mathbf{C},\mathbf{D}$。其中 $\mathbf{A}$ 对应 \cref{ch:eskf} 的离散状态转移矩阵 $\mathbf{F}$（连续时为雅可比），$\mathbf{C}$ 对应观测雅可比 $\mathbf{H}$。
  \item \textbf{Riccati 解矩阵记 $\mathbf{P}$，且\emph{并非}协方差}：LQR/MPC 的 Riccati 解 $\mathbf{P}\succ0$ 是\emph{值函数（代价-to-go）的权重}（即 cost-to-go 的 Hessian），与 \cref{ch:eskf} 的状态协方差 $\mathbf{P}$ 同字母但\textbf{语义相反}（一个是代价，一个是不确定度）。本章凡出现 $\mathbf{P}$ 均指前者；引用 Tedrake 记 $\mathbf{S}$ 处已统一改写为 $\mathbf{P}$。
  \item \textbf{代价权重 $\mathbf{Q},\mathbf{R}$ 也\emph{并非}噪声协方差}：LQR 的 $\mathbf{Q}\succeq0$（罚状态偏差）、$\mathbf{R}\succ0$（罚控制能量）是\emph{代价权重}，与 \cref{ch:eskf} 的过程/观测噪声协方差 $\boldsymbol{\Sigma}_{\mathbf{w}},\boldsymbol{\Sigma}_{\mathbf{v}}$ 不同（控制文献也常用 $\mathbf{Q},\mathbf{R}$ 记噪声，这是“控制—估计对偶”的巧合，但物理意义不同）。
  \item \textbf{反馈律} $\mathbf{u}=-\mathbf{K}\mathbf{x}$，增益 $\mathbf{K}$ \emph{不含}负号；\textbf{代价不带} $\tfrac12$（与 Tedrake/Bansal 一致；个别引 Rawlings 处其代价带 $\tfrac12$，仅整体缩放，不改最优解与 $\mathbf{K}$）。
  \item \textbf{误差}取“期望减实际” $\mathbf{e}=\mathbf{x}_d-\mathbf{x}$（控制惯例）；时域长度记 $N$，与“交叉权重”区分（本章不用 $N$ 记交叉权重，需要时记 $\mathbf{S}_{xu}$）。
  \item \textbf{机器人动力学}沿用广义坐标 $\mathbf{q}$、惯性阵 $\mathbf{M}(\mathbf{q})\succ0$、科氏阵 $\mathbf{C}(\mathbf{q},\dot{\mathbf{q}})$、重力 $\mathbf{g}(\mathbf{q})$；任务雅可比 $\mathbf{J}$（$\dot{\mathbf{x}}=\mathbf{J}\dot{\mathbf{q}}$）；接触力 $\boldsymbol{\lambda}$、接触雅可比 $\mathbf{J}_c$、驱动选择矩阵 $\mathbf{S}=[\mathbf{0}\ \mathbf{I}]$。旋转沿用本书 $\mathbf{R}\in\SO(3)$、$(\cdot)^\wedge$。
\end{itemize}
\end{note}

% ════════════════════════════════════════════════════════════════
\section{反馈控制与稳定性}
\label{sec:ctrl-feedback}

\paragraph{动机：为什么要反馈。}
机器人要把温度、速度、关节角、位姿稳定在某个期望值，并在扰动（坡道、负载、风）下维持。最朴素的想法是\emph{开环}：算好需要的输入 $\mathbf{u}(t)$ 直接施加。可一旦模型不准或有扰动，开环就失准且无从纠正。\textbf{反馈}（闭环）控制把测得的输出 $y$ 与参考 $r$ 比较，用误差 $e=r-y$ 实时修正输入。这一节先讲开环与闭环的根本差异、衡量控制好坏的性能指标，再讲反馈的基本结构与它为何能消除稳态误差，最后建立判断闭环是否“稳定”的数学工具——稳定性是后面 LQR、MPC、机器人控制一切收敛性论证的公共地基。

\subsection{开环 vs 闭环：为什么必须反馈}
\label{sec:ctrl-open-vs-closed}

把同一个对象 $P$ 用两种方式控制，对比最能说清反馈的价值\cite{astrom2008feedback}。设对象零频增益为 $P(0)=g_0$。

\paragraph{开环（前馈）。}
不测输出，直接令 $u=F(s)\,r$。要在常参考下令 $y=r$，须取 $F(0)=1/g_0$。则输出对增益误差的\emph{相对灵敏度}为 $1$：若真实增益变成 $g_0(1+\delta)$，输出就变成 $r(1+\delta)$，误差与模型误差\emph{同阶}。开环对扰动 $d$ 毫无抑制（$y=g_0(u+d)$ 里 $d$ 原样进入输出），但\emph{不会引入不稳定}——开环稳定的对象，开环控制必稳定。

\paragraph{闭环（反馈）。}
取 $u=C(s)(r-y)$，环路增益 $L=CP$。由 \cref{eq:ctrl-comp-sens}，输出对对象增益的相对灵敏度降为 $S=1/(1+L)$：高环路增益 $|L|\gg1$ 时 $|S|\approx1/|L|\ll1$，模型误差被\emph{压缩 $1/|L|$ 倍}；扰动同样被压 $1/|L|$（\cref{sec:ctrl-feedback-struct}）。代价是：高增益可能\emph{失稳}（\cref{eq:ctrl-ess-P} 后的数值例），且测量噪声经 $-T$ 进入输出。

\begin{table}[H]\centering\small
\caption{开环 vs 闭环控制对比}
\label{tab:ctrl-open-closed}
\begin{tabular}{p{0.20\textwidth} p{0.34\textwidth} p{0.34\textwidth}}
\toprule
维度 & 开环（前馈） & 闭环（反馈） \\
\midrule
是否用输出 & 否，靠模型预先算 $u$ & 是，用误差 $e=r-y$ 实时修正 \\
对模型误差 & 灵敏度 $=1$（误差同阶传递） & 灵敏度 $=S=1/(1+L)$，高增益时大幅压缩 \\
对扰动 & 无抑制 & 抑制 $1/(1+L)$ 倍 \\
对测量噪声 & 不受影响（不测量） & 经 $-T$ 进入输出（高频折中） \\
稳定性风险 & 不引入新的不稳定 & 高增益可致失稳，须设计 \\
典型用途 & 已知、可重复、无扰的标称动作 & 一切需镇定/抗扰/纠差的场景 \\
\bottomrule
\end{tabular}
\end{table}

\begin{insight}
反馈的本质是“\emph{用测量换鲁棒}”（\cref{tab:ctrl-open-closed} 的对比）：它把对模型精度的依赖（开环灵敏度 $1$）换成对环路增益的依赖（闭环灵敏度 $1/(1+L)$），从而在模型不准、有扰动时仍能逼近目标。工程上几乎总是\textbf{前馈 + 反馈并用}：前馈负责标称轨迹的快速跟随（不引入稳定性风险），反馈负责纠正模型误差与扰动（\cref{eq:ctrl-setpoint-weight} 的两自由度 PID、\cref{sec:ctrl-robot-ctc} 的计算力矩“前馈逆动力学 + 反馈 PD”都是这一范式）。
\end{insight}

\subsection{性能指标：时域与积分判据}
\label{sec:ctrl-performance}

“稳定”只是底线；工程上还要量化\emph{跟得多快、多准、多稳}。给定单位阶跃参考，闭环输出 $y(t)$ 的标准\textbf{时域指标}\cite{astrom2008feedback}（\cref{tab:ctrl-pid-effect} 的趋势即就这些指标而言）：

\begin{itemize}
  \item \textbf{上升时间} $t_r$：$y$ 从终值 $10\%$ 升到 $90\%$ 所需时间（衡量快慢）；
  \item \textbf{峰值时间} $t_p$：$y$ 首次到达最大值的时刻；
  \item \textbf{超调} $M_p=\dfrac{y_{\max}-y_\infty}{y_\infty}\times100\%$：超过终值的最大相对幅度（衡量振荡）；
  \item \textbf{调节时间} $t_s$：$y$ 进入并保持在终值 $\pm2\%$（或 $\pm5\%$）带内所需时间（衡量稳定快慢）；
  \item \textbf{稳态误差} $e_{ss}=\lim_{t\to\infty}(r-y)$（衡量准，\cref{thm:ctrl-integral} 论其消除）。
\end{itemize}

对标准二阶系统 $\dfrac{\omega_n^2}{s^2+2\zeta\omega_n s+\omega_n^2}$（阻尼比 $\zeta$、自然频率 $\omega_n$），这些指标有闭式\cite{astrom2008feedback}：
\begin{equation}
  M_p=e^{-\pi\zeta/\sqrt{1-\zeta^2}}\ (\zeta<1),\qquad t_s\approx\frac{4}{\zeta\omega_n}\ (\pm2\%),\qquad t_p=\frac{\pi}{\omega_n\sqrt{1-\zeta^2}}.
  \label{eq:ctrl-2nd-order-spec}
\end{equation}
$\zeta$ 越大超调越小（$\zeta\ge1$ 无超调），$\omega_n$ 越大越快——这给了“极点放哪”的物理含义（\cref{sec:ctrl-pole-placement}），也是 \cref{eq:ctrl-ctc-err} 取临界阻尼 $\mathbf{K}_d=2\sqrt{\mathbf{K}_p}$ 的由来。

\paragraph{积分（标量）判据。}
单一数衡量整条误差曲线，便于自动整定与最优控制对比\cite{astrom2008feedback}：
\begin{equation}
  \mathrm{IAE}=\int_0^\infty|e|\,\mathrm dt,\quad
  \mathrm{ISE}=\int_0^\infty e^2\,\mathrm dt,\quad
  \mathrm{ITAE}=\int_0^\infty t\,|e|\,\mathrm dt,\quad
  \mathrm{IE}=\int_0^\infty e\,\mathrm dt.
  \label{eq:ctrl-perf-integral}
\end{equation}
ISE 重罚大误差（与 LQR 的二次代价 \cref{eq:ctrl-lqr-cost} 同源，惟 LQR 还罚控制能量）；ITAE 用时间加权、压制长拖尾，整定出的响应往往最“干净”；IAE 折中；IE（带符号）即 \cref{sec:ctrl-integral-magic} 用到的积分误差，对纯衰减响应等于 IAE。

\begin{insight}
时域指标（$t_r,M_p,t_s$）直观但彼此\emph{冲突}——快（小 $t_r$、大 $\omega_n$）常伴随大超调，是同一组极点的不同侧面。积分判据把整条曲线压成一个数，正好把“多目标手调”变成“单目标优化”：LQR/MPC 正是沿这条路把控制设计\emph{自动化}（最小化 ISE 型代价），这是从经典 PID 整定走向现代最优控制的认知桥。
\end{insight}

\subsection{反馈的基本结构与灵敏度折中}
\label{sec:ctrl-feedback-struct}

标准单回路（SISO、LTI）结构：参考 $r\to$ 误差 $e=r-y\to$ 控制器 $C(s)\to$ 控制量 $u\to$ 对象 $P(s)\to$ 输出 $y$，$y$ 经传感器反馈回求和点\cite{astrom2008feedback}。记环路传递函数 $L(s)=C(s)P(s)$，则

\begin{align}
  \text{参考到输出}\;&: & T(s)&=\frac{Y(s)}{R(s)}=\frac{L(s)}{1+L(s)},\label{eq:ctrl-comp-sens}\\
  \text{灵敏度函数}\;&: & S(s)&=\frac{E(s)}{R(s)}=\frac{1}{1+L(s)},\label{eq:ctrl-sens}
\end{align}
二者满足恒等式
\begin{equation}
  S(s)+T(s)=1.
  \label{eq:ctrl-S-plus-T}
\end{equation}
对象输入处扰动 $d$ 到输出为 $Y/D=P/(1+CP)$；测量噪声 $n$ 到输出为 $Y/N=-T(s)$。

\begin{insight}
\cref{eq:ctrl-S-plus-T} 揭示了反馈设计的根本折中：低频要 $|S|$ 小（跟踪准、抑扰好），高频要 $|T|$ 小（抑制测量噪声），但 $S+T\equiv1$ 使二者\emph{不能同时任意小}。这就是为什么没有“万能”控制器——必须按频段分配性能，这一约束（Bode 灵敏度积分）是经典控制的核心智慧。
\end{insight}

\begin{figure}[ht]
\centering
\begin{tikzpicture}[>=Stealth, node distance=10mm, font=\small,
  blk/.style={draw, rounded corners, minimum height=8mm, minimum width=14mm, fill=main!6, draw=main!60!black},
  sum/.style={draw, circle, inner sep=1pt, minimum size=5mm}]
  \node (r) {$r$};
  \node[sum, right=of r] (s1) {};
  \node[blk, right=of s1] (C) {$C(s)$};
  \node[sum, right=of C] (s2) {};
  \node[blk, right=of s2] (P) {$P(s)$};
  \node[right=14mm of P] (y) {$y$};
  \node[above=5mm of s2] (d) {$d$};
  \draw[->] (r)--(s1) node[pos=0.4,above]{$+$};
  \draw[->] (s1)--(C) node[midway,above]{$e$};
  \draw[->] (C)--(s2) node[midway,above]{$u$};
  \draw[->] (d)--(s2) node[pos=0.3,right]{$+$};
  \draw[->] (s2)--(P);
  \draw[->] (P)--(y);
  \coordinate (fb) at ($(P)!0.5!(y)$);
  \draw[->] (fb) |- ($(s1)+(0,-9mm)$) -| (s1) node[pos=0.93,left]{$-$};
\end{tikzpicture}
\caption{单回路反馈结构：误差 $e=r-y$ 驱动控制器 $C$，对象 $P$ 受控制量 $u$ 与输入扰动 $d$ 作用产生输出 $y$，经单位负反馈回到求和点。（据 \cite{astrom2008feedback} 重绘）}
\label{fig:ctrl-loop}
\end{figure}

\paragraph{纯比例反馈留下稳态误差。}
取纯比例控制 $C(s)=k_p$，参考到误差的传递函数（\cref{eq:ctrl-sens}）为 $G_{er}(s)=1/(1+k_pP(s))$。单位阶跃下（设闭环稳定）稳态误差由终值定理
\begin{equation}
  e_{ss}=G_{er}(0)=\frac{1}{1+k_pP(0)}.
  \label{eq:ctrl-ess-P}
\end{equation}
以对象 $P(s)=1/(s+1)^3$ 为数值例\cite{astrom2008feedback}：$k_p=1,2,5$ 时稳态误差分别为 $0.5,\ 0.33,\ 0.17$；增益越大误差越小但系统越振荡，$k_p=8$ 时\emph{失稳}。可见单靠增大 $k_p$ 无法既消差又稳定——这正是引入积分作用的动机。

\subsection{频域稳定裕度：增益裕度与相位裕度}
\label{sec:ctrl-margins}

“$k_p=8$ 失稳”自然引出一个定量问题：闭环离失稳还有\emph{多远}？频域用环路传递函数 $L(j\omega)=C(j\omega)P(j\omega)$ 的 Nyquist 图（或 Bode 图）回答。闭环极点是 $1+L(s)=0$ 的根，故失稳临界点是 $L(j\omega)=-1$（幅值 $1$、相位 $-180^\circ$）。两个裕度衡量到该点的距离\cite{astrom2008feedback}：

\begin{definition}[增益裕度与相位裕度]
\label{def:ctrl-margins}
设开环稳定、最小相位系统的环路 $L(j\omega)$。
\begin{itemize}
  \item \textbf{相位穿越频率} $\omega_{pc}$：相位 $\angle L(j\omega_{pc})=-180^\circ$ 处；\textbf{增益裕度} $\mathrm{GM}=1/|L(j\omega_{pc})|$（常以 dB 计 $-20\log_{10}|L(j\omega_{pc})|$）——环路增益还能放大多少倍才到 $-1$。
  \item \textbf{增益穿越频率} $\omega_{gc}$：幅值 $|L(j\omega_{gc})|=1$ 处；\textbf{相位裕度} $\mathrm{PM}=180^\circ+\angle L(j\omega_{gc})$——相位还能滞后多少度才到 $-180^\circ$。
\end{itemize}
开环稳定时，$\mathrm{GM}>1$（$>0$\,dB）且 $\mathrm{PM}>0^\circ$ $\Rightarrow$ 闭环稳定（Nyquist 判据）。
\end{definition}

\begin{insight}
裕度是“离失稳还有多少安全余量”的工程语言：$\mathrm{PM}$ 还直接关联时域阻尼——经验关系 $\zeta\approx\mathrm{PM}^\circ/100$，故 $\mathrm{PM}\approx60^\circ$ 对应 $\zeta\approx0.6$（\cref{eq:ctrl-2nd-order-spec} 下超调约 $10\%$），是常用的设计靶值。$\mathrm{PM}$ 也限定可容忍的\emph{纯时延}：时延 $\tau_d$ 在 $\omega_{gc}$ 处贡献相位 $-\omega_{gc}\tau_d$，故系统能容忍的最大时延约 $\mathrm{PM}_{\rm rad}/\omega_{gc}$——这解释了为什么高带宽（大 $\omega_{gc}$）系统对时延/采样延迟格外敏感（呼应 \cref{sec:ctrl-pid-digital} 数字实现的延迟最小化）。\textbf{边界}：GM/PM 是 SISO、最小相位下的可靠判据；MIMO 或非最小相位需回到 Nyquist 围线或 \cref{sec:ctrl-feedback-struct} 的灵敏度峰值 $\|S\|_\infty$。
\end{insight}

\begin{note}
\textbf{LQR 的“天然好裕度”.}\;
\cref{sec:ctrl-lqr} 的（全状态反馈）LQR 有一条著名的鲁棒性保证：单输入情形其环路传递在每个通道上至少有 $\mathrm{GM}\in[\tfrac12,\infty)$（即 $\ge6$\,dB 增益裕度、可承受增益减半到无穷大）与 $\mathrm{PM}\ge60^\circ$（Safonov--Athans 1977；亦见 Anderson--Moore\cite{tedrake_underactuated,rawlings2020mpc}）。这是 LQR 相对手工极点配置的一大优势——最优反馈\emph{自带}良好裕度。\textbf{警告}：该保证属于\emph{状态反馈}；一旦把状态换成 Kalman 估计（LQG，\cref{sec:ctrl-lqe-dual}），裕度保证\emph{不再成立}（Doyle 1978“LQG 无保证裕度”），需用 LQG/LTR 等手段找回。
\end{note}

\paragraph{根轨迹：增益变化时极点怎么走。}
另一经典视角是\textbf{根轨迹}：把闭环特征方程写成 $1+kL_0(s)=0$，画增益 $k$ 从 $0\to\infty$ 时闭环极点在复平面的轨迹。它直观显示“增大增益如何把极点推向虚轴乃至右半平面”。以上面的 $P(s)=1/(s+1)^3$、$C=k_p$ 为例（三重极点在 $s=-1$），根轨迹三条分支从 $s=-1$ 出发，沿 $\pm60^\circ,180^\circ$ 渐近线张开，其中两条随 $k_p$ 增大\emph{穿过虚轴}进入右半平面——穿越点对应的 $k_p$ 正是临界增益（此例 $k_p=8$，呼应 \cref{eq:ctrl-ess-P} 后的数值）。

\begin{figure}[ht]
\centering
\begin{tikzpicture}[>=Stealth, scale=1.0, font=\small]
  % axes
  \draw[->, gray!60!black] (-3.2,0)--(1.6,0) node[right]{$\operatorname{Re}$};
  \draw[->, gray!60!black] (0,-2.2)--(0,2.2) node[above]{$\operatorname{Im}$};
  % open-loop triple pole at s=-1 (cross marker)
  \draw[main, thick] (-1.12,-0.12)--(-0.88,0.12) (-1.12,0.12)--(-0.88,-0.12);
  \node[main, below right, font=\scriptsize] at (-0.95,-0.05) {$s=-1$ (三重极)};
  % asymptote center at -1, angles +-60,180
  % branch 1: stays real, goes left along negative real axis
  \draw[third, thick, ->] (-1,0) -- (-3.0,0);
  % branch 2 & 3: leave at +-60deg, curve toward and cross imag axis
  \draw[third, thick, ->] (-1,0) .. controls (-0.9,0.6) and (-0.3,1.1) .. (0.25,1.7);
  \draw[third, thick, ->] (-1,0) .. controls (-0.9,-0.6) and (-0.3,-1.1) .. (0.25,-1.7);
  % crossing points on imag axis (k = ku)
  \fill[red!70!black] (0,1.06) circle (1.3pt) node[right, red!70!black, font=\scriptsize] {$j\omega_u$ ($k_p=k_u$)};
  \fill[red!70!black] (0,-1.06) circle (1.3pt);
  \node[third, font=\scriptsize] at (0.95,0.95) {$k_p\uparrow$};
  \node[gray!50!black, font=\scriptsize, align=center] at (-2.2,1.4) {左半平面\\（稳定）};
  \node[gray!50!black, font=\scriptsize, align=center] at (1.0,1.7) {右半平面\\（不稳定）};
\end{tikzpicture}
\caption{$P(s)=1/(s+1)^3$ 在比例反馈下的根轨迹示意：三条分支自三重极点 $s=-1$ 出发，一条沿负实轴远离，另两条以 $\pm60^\circ$ 渐近线张开并随 $k_p$ 增大穿过虚轴进入右半平面；穿越点 $j\omega_u$ 处的增益即临界增益 $k_u$（此例 $k_p=8$）。（据 \cite{astrom2008feedback} 原理重绘）}
\label{fig:ctrl-root-locus}
\end{figure}

\begin{insight}
根轨迹与裕度是同一件事的两种画法：根轨迹在 $s$ 平面看\emph{极点随增益的移动}，Nyquist/Bode 在频域看 $L(j\omega)$ \emph{离 $-1$ 的距离}；二者都把“增益—稳定性”的折中可视化。它们也是 \cref{sec:ctrl-pole-placement} 极点配置与 \cref{sec:ctrl-pid-tuning} Ziegler--Nichols 整定（临界增益 $k_u$ 正是根轨迹的虚轴穿越点）的共同几何背景。现代最优控制（LQR/MPC）不再手画这些图，但其“极点放哪、裕度多大”的判读仍是验证设计的必备工具。
\end{insight}

\subsection{积分作用消除稳态误差：一般性证明}
\label{sec:ctrl-integral-magic}

\paragraph{反面：前馈消差要精确模型。}
若改用前馈 $u(t)=k_pe(t)+u_{\rm ff}$，取 $u_{\rm ff}=r/P(0)$，则常参考下稳态输出恰等于参考\cite{astrom2008feedback}。但这要求\emph{精确已知}零频增益 $P(0)$，实际中模型总有误差，前馈消差因此脆弱。积分作用则不需要这种先验知识。

\begin{theorem}[积分作用消除稳态误差]
\label{thm:ctrl-integral}
考虑含积分项（$k_i\neq0$）的控制器
\begin{equation}
  u(t)=k_pe(t)+k_i\int_0^t e(\tau)\,\mathrm d\tau+k_d\frac{\mathrm de}{\mathrm dt}.
  \label{eq:ctrl-pid-time}
\end{equation}
若闭环存在稳态（即 $u(t)\to u_0$、$e(t)\to e_0$ 均收敛到有限常值），则必有 $e_0=0$。此结论\textbf{不假设系统线性或时不变}，只假设存在一个平衡点。
\end{theorem}
\begin{proof}\rebuilt{（据 \cite{astrom2008feedback}）}
对 \cref{eq:ctrl-pid-time} 取 $t\to\infty$ 极限。微分项 $k_d\,\mathrm de/\mathrm dt\to0$（因 $e\to e_0$ 常值），比例项 $\to k_pe_0$，于是
\begin{equation}
  u_0=k_pe_0+k_i\lim_{t\to\infty}\int_0^t e(\tau)\,\mathrm d\tau.
  \label{eq:ctrl-integral-limit}
\end{equation}
关键观察：右端积分项 $\int_0^t e\,\mathrm d\tau$，\emph{除非} $e(\tau)\to0$，否则当 $\tau$ 充分大后被积函数趋于非零常值 $e_0$，积分将发散到 $\pm\infty$。而左端 $u_0$ 有限、$k_i\neq0$，矛盾。故必有 $e_0=0$。
\end{proof}

\begin{insight}
这就是“积分作用的魔力”：积分器是误差的\emph{记忆}——只要还有残差，积分项就持续累积、不断推动控制量，直到误差归零才停下。它的威力\emph{不依赖模型精度}（对比前馈需精确 $P(0)$），这是反馈相对前馈的根本优越性。频域看同一件事：PID 传递函数 $C(s)=k_p+k_i/s+k_ds$ 在零频 $C(0)=\infty$，由 \cref{eq:ctrl-sens} 得 $G_{er}(0)=1/(1+\infty)=0$，阶跃无稳态误差。
\end{insight}

\paragraph{积分增益与扰动抑制。}
稳定闭环初始静止，在对象输入加单位阶跃负载扰动，瞬态后 $e(t)\to0$，由控制律取稳态\cite{astrom2008feedback}：$u(\infty)=k_i\int_0^\infty e(t)\,\mathrm dt$。故单位阶跃扰动的\textbf{积分误差} $\mathrm{IE}=\int_0^\infty e\,\mathrm dt$ 与 $k_i$ 成反比——$k_i$ 越大扰动抑制越好，但过大致振荡、鲁棒性差乃至失稳。这给出 $k_i$ 整定的物理含义。

\begin{pitfall}[概念误区：把“稳态误差为零”当成“任何时候误差都为零”]
\cref{thm:ctrl-integral} 只保证\emph{若存在稳态则稳态误差为零}，它\textbf{不}保证系统一定收敛到稳态，也不说瞬态多好。积分增益过大会使闭环失稳，根本到不了稳态（无稳态时定理前提不成立）。正确理解：积分消除的是\emph{稳态}误差，瞬态性能（超调、调节时间）由全部三项与对象共同决定，仍需整定。
\end{pitfall}

\begin{practice}
\begin{enumerate}
  \item 对 $P(s)=1/(s+1)^3$，用 \cref{eq:ctrl-ess-P} 验证 $k_p=2$ 时 $e_{ss}=1/3$。再说明为何增大 $k_p$ 同时减小 $e_{ss}$ 却恶化稳定性。
  \item （证明）把 \cref{thm:ctrl-integral} 的证明改写为离散时间版本 $u_k=k_pe_k+k_iT_s\sum_{j\le k}e_j$，论证若 $e_k\to e_0$、$u_k\to u_0$ 则 $e_0=0$。\emph{线索}：离散求和发散同理。
  \item （开放）若对象本身含一个积分器（$P(s)$ 在 $s=0$ 有极点，即“1 型系统”），纯比例反馈对阶跃参考还有稳态误差吗？对阶跃\emph{扰动}呢？
\end{enumerate}
\end{practice}

% ────────────────────────────────────────────────
\subsection{Lyapunov 稳定性：三层定义}
\label{sec:ctrl-lyapunov}

\paragraph{动机：什么叫“稳定”。}
直觉上“稳定”是“受扰后不跑飞”，但这含糊。考虑自治系统 $\dot{\mathbf{x}}=f(\mathbf{x})$，平衡点 $f(\mathbf{x}_e)=\mathbf{0}$。我们需要把“留在附近”“最终回到平衡”“以多快回到”分清楚——这正是 Lyapunov 的三层定义\cite{khalil2002nonlinear}。它们是后面一切收敛性论证（LQR 闭环、PD+重力补偿、MPC）的统一语言。

\begin{definition}[Lyapunov 稳定]
\label{def:ctrl-lyap-stable}
平衡点 $\mathbf{x}_e$ \textbf{Lyapunov 稳定（stable i.s.L.）}，若
\begin{equation}
  \forall\,\epsilon>0,\ \exists\,\delta>0\ \text{s.t.}\ \|\mathbf{x}(0)-\mathbf{x}_e\|<\delta\ \Rightarrow\ \forall t\ge0,\ \|\mathbf{x}(t)-\mathbf{x}_e\|<\epsilon.
  \label{eq:ctrl-lyap-stable}
\end{equation}
\end{definition}

\begin{definition}[渐近稳定]
\label{def:ctrl-asymp-stable}
$\mathbf{x}_e$ \textbf{渐近稳定（asymptotically stable）}，若它 Lyapunov 稳定，且
\begin{equation}
  \exists\,\delta>0\ \text{s.t.}\ \|\mathbf{x}(0)-\mathbf{x}_e\|<\delta\ \Rightarrow\ \lim_{t\to\infty}\|\mathbf{x}(t)-\mathbf{x}_e\|=0.
\end{equation}
\end{definition}

\begin{definition}[指数稳定]
\label{def:ctrl-exp-stable}
$\mathbf{x}_e$ \textbf{指数稳定（exponentially stable）}，若它渐近稳定，且存在 $\alpha,\beta,\delta>0$ 使
\begin{equation}
  \|\mathbf{x}(0)-\mathbf{x}_e\|<\delta\ \Rightarrow\ \|\mathbf{x}(t)-\mathbf{x}_e\|\le\alpha\,\|\mathbf{x}(0)-\mathbf{x}_e\|\,e^{-\beta t},\quad\forall t\ge0.
\end{equation}
\end{definition}

\begin{insight}
三者的差别是\emph{逐级加强}：指数稳定 $\subset$ 渐近稳定 $\subset$ Lyapunov 稳定。
\begin{itemize}
  \item Lyapunov 稳定 = “只要起点够近，全程就不会跑远”，但\emph{未必回到}平衡（如无阻尼单摆绕平衡点永远摆动，是 Lyapunov 稳定但非渐近稳定）；
  \item 渐近稳定 = 既不跑远\emph{又最终回到}平衡；
  \item 指数稳定 = 回到平衡且收敛速度\emph{至少几何级快}（给出明确衰减率 $\beta$）。
\end{itemize}
\textbf{边界提醒}：渐近稳定\emph{不能}省去“Lyapunov 稳定”这一前提——存在“最终回到平衡但中途先跑很远”的反例，那样不算渐近稳定。前缀“全局（global）”表示对\emph{所有}初始状态成立（$\delta=\infty$）。
\end{insight}

\begin{figure}[ht]
\centering
\begin{tikzpicture}[>=Stealth, font=\small]
% (a) undamped: centers -> Lyapunov stable but not asymptotic
\begin{scope}
  \draw[->, gray!60!black] (-1.9,0)--(1.9,0) node[right]{$x$};
  \draw[->, gray!60!black] (0,-1.9)--(0,1.9) node[above]{$\dot x$};
  \draw[main, thick] (0,0) circle (0.55);
  \draw[main, thick] (0,0) circle (1.05);
  \draw[main, thick] (0,0) circle (1.5);
  \draw[main, ->] (1.05,0.02)--(1.05,-0.02);
  \draw[main, ->] (-1.05,-0.02)--(-1.05,0.02);
  \fill (0,0) circle (1.3pt);
  \node[align=center, font=\scriptsize] at (0,-2.4) {(a) 无阻尼 $\ddot x+x=0$\\闭轨（中心）：Lyapunov 稳定\\但\emph{非}渐近稳定};
\end{scope}
% (b) damped: stable spiral -> asymptotically stable
\begin{scope}[xshift=5.2cm]
  \draw[->, gray!60!black] (-1.9,0)--(1.9,0) node[right]{$x$};
  \draw[->, gray!60!black] (0,-1.9)--(0,1.9) node[above]{$\dot x$};
  \draw[third, thick, ->] (1.5,0) .. controls (1.4,1.0) and (0.4,1.35) .. (-0.55,0.75)
        .. controls (-1.1,0.35) and (-0.9,-0.55) .. (-0.35,-0.7)
        .. controls (0.25,-0.85) and (0.55,-0.35) .. (0.35,0.1)
        .. controls (0.2,0.35) and (-0.1,0.25) .. (0,0);
  \fill (0,0) circle (1.3pt);
  \node[align=center, font=\scriptsize] at (0,-2.4) {(b) 有阻尼 $\ddot x+b\dot x+x=0$\\收缩螺旋：\emph{渐近}稳定\\（$V$ 沿轨迹严格耗散）};
\end{scope}
\end{tikzpicture}
\caption{相平面（$x$ 对 $\dot x$）直观区分两层稳定：(a) 无阻尼振子的轨迹是绕原点的\emph{闭轨}，能量 $V=\tfrac12\dot x^2+\tfrac12x^2$ 守恒（$\dot V=0$），故 Lyapunov 稳定但不回到平衡；(b) 加阻尼后轨迹\emph{螺旋}收敛到原点，$\dot V=-b\dot x^2<0$ 严格耗散，故渐近稳定。这正是 \cref{thm:ctrl-lyap-direct} 中 $\dot V\le0$（半负定，对应 a）与 $\dot V<0$（负定，对应 b）之别。}
\label{fig:ctrl-phase-portrait}
\end{figure}

\paragraph{离散时间版。}
对 $\mathbf{x}_{k+1}=f(\mathbf{x}_k)$，把 $t$ 换成 $k$、$\lim_{t\to\infty}$ 换成 $\lim_{k\to\infty}$ 即得对应定义\cite{khalil2002nonlinear}。MPC 章用的就是离散版。

% ────────────────────────────────────────────────
\subsection{Lyapunov 直接法（第二法）}
\label{sec:ctrl-lyap-direct}

\paragraph{动机与几何直觉。}
要判稳定，最直接是求解 $\dot{\mathbf{x}}=f(\mathbf{x})$ 看轨迹（如 \cref{fig:ctrl-phase-portrait} 的相平面）——但非线性系统通常无闭式解。Lyapunov 的天才在于：\emph{不解方程}，只找一个“广义能量”函数 $V(\mathbf{x})$，若沿轨迹这个能量不增（甚至递减），系统就不会远离平衡。物理直觉：一个有阻尼的机械系统，总能量随时间耗散，故最终停在最低能量点（平衡，即 \cref{fig:ctrl-phase-portrait}b 的收缩螺旋）。

\begin{theorem}[Lyapunov 直接法]
\label{thm:ctrl-lyap-direct}
对 $\dot{\mathbf{x}}=f(\mathbf{x})$，平衡点取在 $\mathbf{x}=\mathbf{0}$。若存在连续可微 $V:\mathbb{R}^n\to\mathbb{R}$ 满足
\begin{enumerate}
  \item $V(\mathbf{0})=0$，且 $V(\mathbf{x})>0,\ \forall\mathbf{x}\neq\mathbf{0}$（\textbf{正定}）；
  \item $\dot V(\mathbf{x})=\nabla V(\mathbf{x})\cdot f(\mathbf{x})=\sum_{i}\frac{\partial V}{\partial x_i}f_i(\mathbf{x})\le0,\ \forall\mathbf{x}\neq\mathbf{0}$（沿轨迹\textbf{半负定}），
\end{enumerate}
则原点 Lyapunov 稳定。若进一步 $\dot V(\mathbf{x})<0,\ \forall\mathbf{x}\neq\mathbf{0}$（\textbf{负定}），则原点\textbf{渐近稳定}。若再加 $V$ \textbf{径向无界}（$V(\mathbf{x})\to\infty$ 当 $\|\mathbf{x}\|\to\infty$），则\textbf{全局渐近稳定}。
\end{theorem}

\begin{insight}
把 $V$ 想成“到平衡点的广义距离/能量”。$\dot V\le0$ 说能量不增，故轨迹被困在 $V$ 的某个水平集内（不跑远）$\Rightarrow$ Lyapunov 稳定；$\dot V<0$ 说能量严格耗散到零，故状态被逼向平衡 $\Rightarrow$ 渐近稳定；径向无界保证水平集 $\{V\le c\}$ 有界，避免轨迹沿“能量谷”逃逸到无穷 $\Rightarrow$ 全局。$V$ 是\emph{纯代数对象}（不必有物理能量意义），但取能量函数往往最自然——这正是 \cref{sec:ctrl-robot} PD+重力补偿证明的思路。
\end{insight}

\paragraph{承重假设三问（径向无界）。}
为什么“全局”要求径向无界？\emph{(i) 用在哪步}：用于断言水平集 $\{V\le c\}$ 有界，从而轨迹不能从任意远处出发逃逸；\emph{(ii) 删了崩在哪}：若 $V$ 不径向无界，可能存在 $V$ 沿某方向趋于有限上界的“能量谷”，轨迹沿谷跑向无穷而 $V$ 仍递减，结论崩在“轨迹有界”这一步；\emph{(iii) 可否弱化}：可换成“某个水平集 $\{V\le c_0\}$ 有界”，则得\emph{局部}（吸引域至少含该水平集）而非全局结论。

% ────────────────────────────────────────────────
\subsection{LaSalle 不变原理}
\label{sec:ctrl-lasalle}

\paragraph{动机：$\dot V$ 只半负定怎么办。}
机械系统最常见的情形是 $V=$ 总能量、$\dot V=-$ 阻尼耗散，而阻尼\emph{仅在速度非零时}严格负——即 $\dot V\le0$ 但 $\dot V=0$ 在整条“速度为零”的子流形上成立，\emph{不是}负定。此时直接法只给“稳定”，给不出“渐近稳定”。LaSalle 不变原理正是为补这一步而生，它是 PD+重力补偿、Cheetah 落脚控制等机器人稳定性证明的关键工具\cite{khalil2002nonlinear}。

\begin{theorem}[LaSalle 不变原理]
\label{thm:ctrl-lasalle}
自治系统 $\dot{\mathbf{x}}=f(\mathbf{x})$，$f(\mathbf{0})=\mathbf{0}$。设存在 $C^1$ 函数 $V(\mathbf{x})$ 使 $\dot V(\mathbf{x})\le0$（半负定）。令
\begin{equation}
  \mathcal{I}=\text{集合 }\{\mathbf{x}:\dot V(\mathbf{x})=0\}\text{ 内的\emph{最大不变集}（完整轨迹之并）.}
\end{equation}
则任意有界轨迹的极限点集都含于 $\mathcal{I}$，即轨迹趋于 $\mathcal{I}$。
\emph{渐近稳定推论}：若再有 (a) $V$ 正定、$V(\mathbf{0})=0$；(b) $\mathcal{I}$ 中除 $\mathbf{x}(t)\equiv\mathbf{0}$ 外不含任何完整轨迹；(c) $V$ 径向无界，则原点\textbf{全局渐近稳定}。
\end{theorem}
\begin{proof}\rebuilt{（标准证明骨架，据 \cite{khalil2002nonlinear}）}
紧正不变集 $\Omega$（如某个有界水平集 $\{V\le c\}$）上 $\dot V\le0$，故 $V$ 沿任一轨迹单调不增且有下界，从而 $V(\mathbf{x}(t))\to L$（某常数）。设 $\boldsymbol{\omega}(\mathbf{x}_0)$ 为该轨迹的正极限集，它是闭、不变的，且因 $V$ 连续，$V\equiv L$ 于 $\boldsymbol{\omega}(\mathbf{x}_0)$ 上。在 $\boldsymbol{\omega}(\mathbf{x}_0)$ 上 $V$ 恒为常数 $\Rightarrow\dot V=0$，故 $\boldsymbol{\omega}(\mathbf{x}_0)\subseteq\{\dot V=0\}$；又 $\boldsymbol{\omega}(\mathbf{x}_0)$ 不变，故 $\boldsymbol{\omega}(\mathbf{x}_0)\subseteq\mathcal{I}$（$\{\dot V=0\}$ 中最大不变集）。轨迹趋于其正极限集，即趋于 $\mathcal{I}$。在推论条件 (b) 下 $\mathcal{I}=\{\mathbf{0}\}$，故轨迹趋于原点；配 (a)(c) 即全局渐近稳定。
\end{proof}

\begin{pitfall}[思维陷阱：把 LaSalle 用于非自治系统]
LaSalle 只对\emph{自治}（时不变）系统 $\dot{\mathbf{x}}=f(\mathbf{x})$ 成立——证明依赖正极限集的不变性，时变系统没有这个结构。自适应控制、轨迹跟踪等系统是\emph{非自治}的（$f$ 显含 $t$），此时不能用 LaSalle，须改用 \textbf{Barbalat 引理}（见下）。误用 LaSalle 于非自治系统是常见错误，会给出错误的渐近稳定结论。
\end{pitfall}

% ────────────────────────────────────────────────
\subsection{线性系统的稳定性与 Lyapunov 方程}
\label{sec:ctrl-lyap-linear}

线性系统是 Lyapunov 理论最干净的特例，也是 LQR 闭环分析的工具。

\begin{theorem}[线性系统稳定性与 Lyapunov 方程]
\label{thm:ctrl-lyap-eq}
线性系统 $\dot{\mathbf{x}}=\mathbf{A}\mathbf{x}$ 渐近稳定（对线性系统等价于指数稳定、全局）当且仅当 $\mathbf{A}$ 的所有特征值实部为负（\textbf{Hurwitz}）。等价地：对任意给定 $\mathbf{Q}=\mathbf{Q}^\top\succ0$，Lyapunov 方程
\begin{equation}
  \boxed{\mathbf{A}^\top\mathbf{M}+\mathbf{M}\mathbf{A}=-\mathbf{Q}}
  \label{eq:ctrl-lyap-eq}
\end{equation}
有唯一正定解 $\mathbf{M}=\mathbf{M}^\top\succ0$，对应 Lyapunov 函数 $V(\mathbf{x})=\mathbf{x}^\top\mathbf{M}\mathbf{x}$。离散版：$\mathbf{x}_{k+1}=\mathbf{A}\mathbf{x}_k$ 稳定 $\iff$ $\mathbf{A}$ 谱半径 $<1$（\textbf{Schur}），对应离散 Lyapunov 方程 $\mathbf{A}^\top\mathbf{M}\mathbf{A}-\mathbf{M}=-\mathbf{Q}$。
\end{theorem}
\begin{proof}\rebuilt{（$\mathbf{A}$ Hurwitz $\Rightarrow$ 解存在且正定，据 \cite{khalil2002nonlinear}）}
$\mathbf{A}$ Hurwitz 时，定义
\begin{equation}
  \mathbf{M}=\int_0^\infty e^{\mathbf{A}^\top t}\,\mathbf{Q}\,e^{\mathbf{A}t}\,\mathrm dt.
  \label{eq:ctrl-lyap-int}
\end{equation}
因 $\mathbf{A}$ Hurwitz，$\|e^{\mathbf{A}t}\|$ 指数衰减，积分收敛；被积矩阵正定（$\mathbf{Q}\succ0$）故 $\mathbf{M}\succ0$。代入验证：
\begin{align}
  \mathbf{A}^\top\mathbf{M}+\mathbf{M}\mathbf{A}
  &=\int_0^\infty\!\Big(\mathbf{A}^\top e^{\mathbf{A}^\top t}\mathbf{Q}e^{\mathbf{A}t}+e^{\mathbf{A}^\top t}\mathbf{Q}e^{\mathbf{A}t}\mathbf{A}\Big)\mathrm dt
   =\int_0^\infty\frac{\mathrm d}{\mathrm dt}\Big(e^{\mathbf{A}^\top t}\mathbf{Q}e^{\mathbf{A}t}\Big)\mathrm dt\notag\\
  &=\Big[e^{\mathbf{A}^\top t}\mathbf{Q}e^{\mathbf{A}t}\Big]_0^\infty=\mathbf{0}-\mathbf{Q}=-\mathbf{Q}.
\end{align}
中间用 $\frac{\mathrm d}{\mathrm dt}e^{\mathbf{A}t}=\mathbf{A}e^{\mathbf{A}t}$ 与上限处 $e^{\mathbf{A}t}\to0$。唯一性由方程对 $\mathbf{M}$ 是线性、且 $\mathbf{A}\oplus\mathbf{A}$（Kronecker 和）无零特征值（因 $\mathbf{A}$ Hurwitz $\Rightarrow\lambda_i+\lambda_j\neq0$）保证。
\end{proof}

\begin{pitfall}[证明陷阱：把“正交/对称矩阵”性质误用在 $\mathbf{A}^\top\mathbf{M}+\mathbf{M}\mathbf{A}$ 上]
$\mathbf{A}^\top\mathbf{M}+\mathbf{M}\mathbf{A}$ 是 $\dot V$ 的核（$V=\mathbf{x}^\top\mathbf{M}\mathbf{x}\Rightarrow\dot V=\mathbf{x}^\top(\mathbf{A}^\top\mathbf{M}+\mathbf{M}\mathbf{A})\mathbf{x}$），它\emph{对称}（即便 $\mathbf{A}$ 不对称）。新手易误以为“$\mathbf{A}$ 稳定就有 $\mathbf{A}+\mathbf{A}^\top\prec0$”——\textbf{错}：$\mathbf{A}+\mathbf{A}^\top$ 是 $\mathbf{M}=\mathbf{I}$ 这一\emph{特例}的 $\dot V$ 核，稳定矩阵未必满足（如 $\mathbf{A}=\big[\begin{smallmatrix}-1&10\\0&-1\end{smallmatrix}\big]$ 稳定但 $\mathbf{A}+\mathbf{A}^\top$ 不负定）。正确做法是\emph{对给定 $\mathbf{Q}$ 解出 $\mathbf{M}$}（\cref{eq:ctrl-lyap-int}），不能直接取 $\mathbf{M}=\mathbf{I}$。
\end{pitfall}

% ────────────────────────────────────────────────
\subsection{Lyapunov 间接法与 Barbalat 引理}
\label{sec:ctrl-indirect}

\begin{theorem}[Lyapunov 间接法（线性化）]
\label{thm:ctrl-indirect}
平衡点 $\mathbf{x}_e$ 处取雅可比 $\mathbf{J}=\nabla f(\mathbf{x}_e)$。则\cite{khalil2002nonlinear}：
\begin{itemize}
  \item 若 $\mathbf{J}$ 所有特征值实部 $<0$，则 $\mathbf{x}_e$ 渐近稳定；
  \item 若 $\mathbf{J}$ 有特征值实部 $>0$，则 $\mathbf{x}_e$ 不稳定；
  \item 若 $\mathbf{J}$ 有特征值实部 $=0$（临界），间接法\textbf{失效}，须用直接法或 LaSalle。
\end{itemize}
\end{theorem}

间接法把非线性稳定性归结为线性化的特征值，简单但只给局部结论且在临界情形失效。临界情形的典型例是无阻尼振子，其线性化有纯虚特征值，需直接法判定。

\begin{lemma}[Barbalat 引理]
\label{lem:ctrl-barbalat}
若 $f(t)$ 当 $t\to\infty$ 有有限极限，且 $\dot f$ 一致连续（充分条件：$\ddot f$ 有界），则 $\dot f(t)\to0$（$t\to\infty$）\cite{khalil2002nonlinear}。
\end{lemma}

\begin{insight}
Barbalat 引理是“非自治版的 LaSalle”：当 $\dot V\le-W(t)\le0$ 而系统时变（不能用 LaSalle）时，先由 $V$ 有下界且单调不增得 $V\to$ 常数，故 $\int_0^\infty W\,\mathrm dt<\infty$；若 $W$ 一致连续，Barbalat 给 $W(t)\to0$。这是自适应控制收敛性证明的核心。\textbf{边界}：Barbalat 不要求自治，但要求 $\dot f$（或 $W$）一致连续——少了这个条件结论不成立（存在 $\int f<\infty$ 但 $f\not\to0$ 的尖峰反例）。
\end{insight}

\begin{practice}
\begin{enumerate}
  \item （证明）对 $\dot x=-x^3$，取 $V=\tfrac12x^2$，证 $\dot V=-x^4<0$（$x\neq0$），故全局渐近稳定。再用间接法（\cref{thm:ctrl-indirect}）看线性化 $\mathbf{J}=0$ 为何失效——这正说明直接法更强。
  \item （推导，在草稿纸上完成）对 $\dot{\mathbf{x}}=\mathbf{A}\mathbf{x}$，$\mathbf{A}=\big[\begin{smallmatrix}0&1\\-2&-3\end{smallmatrix}\big]$，取 $\mathbf{Q}=\mathbf{I}$ 解 Lyapunov 方程 \cref{eq:ctrl-lyap-eq} 得 $\mathbf{M}$，验证 $\mathbf{M}\succ0$，从而确认 $\mathbf{A}$ Hurwitz。
  \item （开放）单摆 $\ddot\theta+\sin\theta=0$（无阻尼）取 $V=\tfrac12\dot\theta^2+(1-\cos\theta)$。算 $\dot V$，说明只能判“Lyapunov 稳定”而非渐近稳定。加阻尼项 $-b\dot\theta$ 后用 LaSalle 论证渐近稳定。
\end{enumerate}
\end{practice}

% ════════════════════════════════════════════════════════════════
\section{PID 控制}
\label{sec:ctrl-pid}

\paragraph{动机：最朴素也最普及的控制器。}
PID（比例-积分-微分）是工业界最广泛使用的控制器：约九成的工业回路是 PID 或其变体。它只需对象的极少信息（甚至不需模型），却能稳定一大类一阶/二阶系统。本节系统讲三项的作用、整定方法、两个工程必须处理的实际问题（微分放大噪声、积分饱和），以及数字实现——这些是机器人关节伺服、姿态环、温控等的基础。

\subsection{三项与两种参数化}
\label{sec:ctrl-pid-three}

\paragraph{时域控制律。}
误差 $e(t)=r(t)-y(t)$（参考减输出），理想 PID\cite{astrom2008feedback}：
\begin{equation}
  u(t)=k_pe(t)+k_i\int_0^t e(\tau)\,\mathrm d\tau+k_d\frac{\mathrm de}{\mathrm dt}
       =k_p\!\left(e+\frac{1}{T_i}\int_0^t e(\tau)\,\mathrm d\tau+T_d\frac{\mathrm de}{\mathrm dt}\right).
  \label{eq:ctrl-pid-two}
\end{equation}
历史上称“三项控制器”：比例项看\emph{现在}的误差、积分项看\emph{过去}的累积、微分项\emph{预测}未来。传递函数
\begin{equation}
  C(s)=k_p+\frac{k_i}{s}+k_ds=\frac{k_ds^2+k_ps+k_i}{s}.
  \label{eq:ctrl-pid-tf}
\end{equation}

\paragraph{两种等价参数化。}
\begin{itemize}
  \item \textbf{增益式}：比例增益 $k_p$、积分增益 $k_i$、微分增益 $k_d$；
  \item \textbf{时间式（ISA）}：\emph{积分时间} $T_i=k_p/k_i$、\emph{微分时间} $T_d=k_d/k_p$（量纲为时间，可与对象时间常数对比）；\emph{比例带} $\mathrm{PB}=100/k_p$（\%），$\mathrm{PB}=10\%$ 表示控制器在测量量程 $10\%$ 内线性工作。
\end{itemize}

\begin{table}[H]\centering\small
\caption{PID 三项的定性影响（趋势表；各参数耦合，非独立，单调改一项的趋势）}
\label{tab:ctrl-pid-effect}
\begin{tabular}{lccccc}
\toprule
参数增大 & 上升时间 & 超调 & 调节时间 & 稳态误差 & 稳定性 \\
\midrule
$k_p\uparrow$ & 减小 & 增大 & 小变化 & 减小 & 变差 \\
$k_i\uparrow$ & 减小 & 增大 & 增大 & 消除 & 变差 \\
$k_d\uparrow$ & 微变 & 减小 & 减小 & 无影响 & 改善（小增益时） \\
\bottomrule
\end{tabular}
\end{table}

\begin{insight}
三项的本质各异：\textbf{P} 看现在（对当前误差的刚度）；\textbf{I} 看过去（记忆，消稳态误差，\cref{thm:ctrl-integral}）；\textbf{D} 看未来。微分是\emph{预测}：把比例+微分合写 $u=k_p(e+T_d\,\mathrm de/\mathrm dt)=:k_pe_p$，其中 $e_p$ 恰是误差在 $t+T_d$ 的\emph{线性外推预测}，预测时间正是微分时间 $T_d$。这解释了 D 为何提供阻尼——它对“误差正在变大”提前反应。
\end{insight}

\subsection{微分滤波与设定点加权}
\label{sec:ctrl-pid-filter}

\paragraph{反面：理想微分放大噪声。}
理想微分 $k_ds$ 的幅频随频率线性增长，对高频测量噪声增益无穷大——实际中会把传感器噪声放大成剧烈的控制抖振。对策是“信号减其低通版本”实现\emph{滤波微分}\cite{astrom2008feedback}：
\begin{equation}
  \frac{k_ds}{1+sT_f},\qquad T_f=\frac{T_d}{N},\ N\in[5,20].
  \label{eq:ctrl-filt-deriv}
\end{equation}
低频（$|s|\ll1/T_f$）时近似理想微分 $k_ds$，高频时增益饱和为常值 $k_d/T_f=k_pN$（不再无穷），既保留预测又不爆噪声。

\paragraph{设定点加权（避免微分踢与阶跃冲击）。}
误差反馈型三项都作用于 $e=r-y$，参考阶跃时 $\mathrm de/\mathrm dt$ 出现冲激（“derivative kick”），微分项产生巨大控制峰值。两自由度（2-DOF）PID 让比例/微分只作用于\emph{测量} $y$、积分仍作用于误差\cite{astrom2008feedback}：
\begin{equation}
  u=k_p(\beta r-y)+k_i\int_0^t\!\big(r(\tau)-y(\tau)\big)\mathrm d\tau+k_d\!\left(\gamma\frac{\mathrm dr}{\mathrm dt}-\frac{\mathrm dy}{\mathrm dt}\right),
  \label{eq:ctrl-setpoint-weight}
\end{equation}
$\beta,\gamma\in[0,1]$ 为\emph{设定点权}。\textbf{积分必须作用于误差}以保证稳态误差为零。$\gamma$ 通常取 $0$（杜绝微分踢），$\beta<1$ 减小阶跃超调；$\beta=\gamma=0$ 称 I-PD 控制器。不同 $\beta,\gamma$ 对扰动/噪声响应相同，只改变对参考的响应。

\begin{pitfall}[编程陷阱：数字 PID 直接对参考做微分]
代码里若写 \texttt{D = kd*(e - e\_prev)/h} 而 \texttt{e=r-y}，则参考 $r$ 一旦阶跃，$D$ 项瞬间产生巨大尖峰（微分踢），可能让执行器饱和甚至损坏。现象：每次改设定点，控制量先猛冲一下。根本原因：阶跃参考的导数是冲激。正确做法：微分只对测量 $y$（\cref{eq:ctrl-setpoint-weight} 取 $\gamma=0$）——\texttt{D = ad*D - bd*(y - y\_prev)}，即对 $y$ 用滤波后向差分（见 \cref{sec:ctrl-pid-digital}）。
\end{pitfall}

\subsection{PID 整定：Ziegler--Nichols 与改进规则}
\label{sec:ctrl-pid-tuning}

Ziegler--Nichols（Z-N，1942\cite{ziegler1942optimum}）是首套系统整定规则：做简单实验提取对象的时/频域特征，再查表定参数。

\paragraph{(a) 阶跃响应法（时域）。}
对开环对象做单位阶跃，由响应曲线读两参数：$a$（最陡切线在纵轴的截距）与 $\tau$（切线与时间轴交点，近似时延）。$a/\tau$ 是最陡斜率。

\paragraph{(b) 频率响应法（临界比例度法）。}
接上控制器，置 $k_i=k_d=0$，逐步增大 $k_p$ 至闭环出现\emph{等幅振荡}，记此临界增益 $k_u$（亦记 $K_c$）与振荡周期 $T_u$（亦记 $T_c$）。由 Nyquist 判据，此时 $L=k_uP$ 在 $\omega_u=2\pi/T_u$ 恰过临界点 $-1$，即捕获了对象相位滞后 $180^\circ$ 的点——几何上正是根轨迹（\cref{fig:ctrl-root-locus}）穿过虚轴的那一点，对应增益裕度恰好用尽（$\mathrm{GM}=1$，\cref{def:ctrl-margins}）。

\begin{table}[H]\centering\small
\caption{Ziegler--Nichols 整定规则（频率响应法/临界比例度法），以临界增益 $k_u$、临界周期 $T_u$ 表}
\label{tab:ctrl-zn-freq}
\begin{tabular}{lccc}
\toprule
类型 & $k_p$ & $T_i$ & $T_d$ \\
\midrule
P   & $0.5\,k_u$  & --        & --        \\
PI  & $0.45\,k_u$ & $T_u/1.2$ & --        \\
PID & $0.6\,k_u$  & $T_u/2$   & $T_u/8$   \\
\bottomrule
\end{tabular}
\end{table}

\begin{table}[H]\centering\small
\caption{Ziegler--Nichols 整定规则（阶跃响应法），以 $a$、$\tau$ 表}
\label{tab:ctrl-zn-step}
\begin{tabular}{lccc}
\toprule
类型 & $k_p$ & $T_i$ & $T_d$ \\
\midrule
P   & $1/a$   & --          & --        \\
PI  & $0.9/a$ & $\tau/0.3$  & --        \\
PID & $1.2/a$ & $\tau/0.5$  & $0.5\tau$ \\
\bottomrule
\end{tabular}
\end{table}

\rebuilt{\cref{tab:ctrl-zn-freq} 与教材另一常见写法（以 $k_i,k_d$ 直接给）一致：PID 取 $k_p=0.6k_u,\ k_i=1.2k_u/T_u,\ k_d=3k_uT_u/40$，换算回 ISA 得 $T_i=k_p/k_i=0.5T_u$、$T_d=k_d/k_p=T_u/8$，与本表吻合。}

\paragraph{继电反馈自整定。}
把对象接入与\emph{继电（relay）非线性}的反馈环，多数系统会起振：继电输出为方波（幅值 $d$，基波幅值 $4d/\pi$），过程输出近正弦（幅值 $a$），二者基波相差 $180^\circ$，即以临界周期 $T_u$ 振荡\cite{astrom2008feedback}。在 $\omega_u$ 处过程增益 $|P(i\omega_u)|=\pi a/(4d)$，故
\begin{equation}
  k_u=\frac{4d}{\pi a}.
  \label{eq:ctrl-relay}
\end{equation}
得 $k_u,T_u$ 后查 Z-N 表。此法可自动化（一键触发、自动保持小振荡、整定毕切回 PID）。Z-N 的缺陷是用过程信息太少，所得闭环鲁棒性偏弱；基于 FOTD（一阶+时延）模型的改进规则通常给出更低增益、更好鲁棒性\cite{astrom2008feedback}。

\subsection{积分饱和与抗饱和}
\label{sec:ctrl-pid-windup}

\paragraph{现象：integrator windup。}
执行器有限（电机限速、阀门全开/全闭）。控制量达限后，反馈环\emph{断开}、系统开环运行；但误差非零，积分项仍持续累积、变得很大；即便误差后来变号，巨大的积分项仍把控制信号顶在饱和限，需很久才退出饱和，造成大超调与长瞬态。这就是\textbf{积分饱和}\cite{astrom2008feedback}。

\paragraph{对策：回算（back-calculation）。}
设控制器输出 $u_a$、经饱和后实际输出 $u=\mathrm{sat}(u_a)$，差 $e_s=u-u_a$ 经增益 $k_{aw}$ 馈回积分器输入\cite{astrom2008feedback,khalil2002nonlinear}：
\begin{equation}
  \dot I=k_ie+\frac{1}{T_{aw}}\big(\mathrm{sat}(u_a)-u_a\big),\qquad T_{aw}=1/k_{aw}.
  \label{eq:ctrl-antiwindup}
\end{equation}
无饱和时 $e_s=0$，额外环无作用；饱和时 $e_s$ 反馈使积分器输出迅速复位，把控制器输出\emph{保持在饱和限附近}，误差一变号即可退出饱和。复位速率由 $k_{aw}$ 控制（常取为积分增益 $k_i$ 的若干倍，如一个数量级）；过大则测量噪声会引起不良复位，常取跟踪时间 $T_{aw}\approx\sqrt{T_iT_d}$ 或 $T_{aw}=T_i$。\textbf{增量式数字 PID}（见下）天然抗饱和，因为它不显式累加积分。

\begin{pitfall}[思维陷阱：以为限幅输出就等于抗饱和]
新手常在控制律最后加一句 \texttt{u = clamp(u, umin, umax)} 就以为解决了饱和。这只限制了\emph{输出}，但积分器内部仍在疯狂累积（它看不到饱和），windup 照样发生，退饱和照样慢。根本原因：抗饱和必须让\emph{积分器本身}感知饱和并停止/回退累积。正确做法：用回算 \cref{eq:ctrl-antiwindup}（把饱和差反馈进积分器），或条件积分（饱和且误差加剧饱和时停止累加），或改用增量式。
\end{pitfall}

\subsection{数字（离散）实现}
\label{sec:ctrl-pid-digital}

\paragraph{位置式与增量式。}
采样周期 $T_s$，$e_k=e(kT_s)$。\textbf{位置式}（直接算绝对控制量）：
\begin{equation}
  u_k=k_pe_k+k_iT_s\sum_{j=0}^{k}e_j+\frac{k_d}{T_s}(e_k-e_{k-1}).
  \label{eq:ctrl-pid-position}
\end{equation}
\textbf{增量式}（算控制量增量，对抗饱和、无扰切换更友好）：
\begin{equation}
  \Delta u_k=k_p(e_k-e_{k-1})+k_iT_s\,e_k+\frac{k_d}{T_s}(e_k-2e_{k-1}+e_{k-2}),\quad u_k=u_{k-1}+\Delta u_k.
  \label{eq:ctrl-pid-velocity}
\end{equation}

\paragraph{滤波微分的后向差分。}
连续滤波微分 $T_f\dot D+D=-k_d\dot y$（对测量 $y$ 微分），用后向差分离散\cite{astrom2008feedback}：
\begin{equation}
  D(t_k)=\frac{T_f}{T_f+T_s}D(t_{k-1})-\frac{k_d}{T_f+T_s}\big(y(t_k)-y(t_{k-1})\big).
  \label{eq:ctrl-pid-Ddiff}
\end{equation}
后向差分的优点：系数 $T_f/(T_f+T_s)\in[0,1)$ 对所有 $T_s>0$ 成立，保证差分方程稳定（前向差分则不然）。积分用前向近似加回算项（\cref{eq:ctrl-antiwindup} 离散化）。

\begin{codebox}[Python]{带滤波微分 + 抗饱和 + 设定点加权的数字 PID（一个采样周期）}
# 预计算系数（h = Ts 采样周期）
bi = ki * h                 # 积分项系数
ad = Tf / (Tf + h)          # 滤波微分的状态保持系数, in [0,1)
bd = kd / (Tf + h)          # 滤波微分的测量增益
br = h / Taw                # 抗饱和(回算)系数, Taw = 1/kaw

# 状态(跨周期保持): I 积分, D 滤波微分, y_old 上一次测量
def pid_step(r, y):
    global I, D, y_old
    P  = kp * (b * r - y)            # 比例 + 设定点加权 (b = beta)
    D  = ad * D - bd * (y - y_old)   # 滤波微分, 只对测量 y (避免 derivative kick)
    u_a = P + I + D                  # 临时输出(未饱和)
    u  = max(u_low, min(u_high, u_a))# 执行器饱和 sat()
    I  = I + bi * (r - y) + br * (u - u_a)  # 更新积分(含抗饱和回算)
    y_old = y                        # 更新微分状态
    return u
# 注: 先算输出再更新状态, 最小化输入到输出的延迟; 全部为定点/浮点可实现
\end{codebox}

\subsection{PID 闭环稳定性：Routh--Hurwitz}
\label{sec:ctrl-pid-routh}

\rebuilt{PID 接在对象上构成闭环特征多项式 $1+C(s)P(s)=0$，其稳定性由该多项式根的位置决定，可用 Routh--Hurwitz 判据判定（据 \cite{astrom2008feedback} 通用结论与控制教材）。}
以二阶对象 $P(s)=1/(s^2+a_1s+a_0)$ 为例。PD 控制 $C(s)=k_p+k_ds$：
\begin{equation}
  1+\frac{k_p+k_ds}{s^2+a_1s+a_0}=0\;\Rightarrow\;s^2+(a_1+k_d)s+(a_0+k_p)=0,
\end{equation}
二阶 Routh--Hurwitz 稳定 $\iff$ 系数全正：$a_1+k_d>0$ 且 $a_0+k_p>0$——可见 $k_d$ 增阻尼、$k_p$ 提刚度。加积分变三阶 $C=k_p+k_i/s+k_ds$：
\begin{equation}
  s^3+(a_1+k_d)s^2+(a_0+k_p)s+k_i=0,
\end{equation}
三阶 Routh--Hurwitz 稳定条件：所有系数 $>0$ 且
\begin{equation}
  (a_1+k_d)(a_0+k_p)>k_i.
  \label{eq:ctrl-routh3}
\end{equation}
\cref{eq:ctrl-routh3} 给出 $k_i$ 的\textbf{上界}：积分增益过大则失稳——这正呼应 Z-N 整定为何不能把 $k_i$ 取得太大，也印证了 \cref{sec:ctrl-integral-magic} 关于“$k_i$ 大扰动抑制好但鲁棒性差”的论断。

\begin{practice}
\begin{enumerate}
  \item （推导）对一阶对象 $P(s)=b/(s+a)$ 用 PI 控制，写出闭环特征多项式，令其等于 $s^2+a_1s+a_2$，解出 $k_p=(a_1-a)/b,\ k_i=a_2/b$（极点配置）。
  \item （整定）对 $P(s)=e^{-s}/(s+1)$ 用继电法（\cref{eq:ctrl-relay}）估 $k_u$，再查 \cref{tab:ctrl-zn-freq} 设计 PI，说明为何 Z-N 结果偏激进。
  \item （编程/分析）解释 \cref{eq:ctrl-pid-Ddiff} 中系数 $T_f/(T_f+T_s)$ 为何必落在 $[0,1)$，从而差分方程对任意 $T_s>0$ 稳定；前向差分 $D_k=(1-T_s/T_f)D_{k-1}+\cdots$ 何时会失稳？
\end{enumerate}
\end{practice}

% ════════════════════════════════════════════════════════════════
\section{状态空间：可控性与可观性}
\label{sec:ctrl-statespace}

\paragraph{动机：从单输入单输出到多变量。}
PID 是 SISO（单输入单输出）+ 固定结构的控制器。机器人是\emph{多输入多输出（MIMO）}、强耦合的：多个关节力矩同时影响多个关节角。要系统地设计 MIMO 控制器（如下一节的 LQR、状态反馈），先要回答两个根本问题：能否用输入把状态\emph{驱动}到任意目标（可控性）？能否由输出\emph{推断}内部状态（可观性）？这两个概念也是 LQR/MPC 解存在唯一性的前提，并与 \cref{ch:eskf} 的滤波可观性直接对偶。

\subsection{可控性}
\label{sec:ctrl-controllability}

\begin{definition}[完全状态可控]
\label{def:ctrl-controllable}
系统 $\mathbf{x}_{k+1}=\mathbf{A}\mathbf{x}_k+\mathbf{B}\mathbf{u}_k$（或连续 $\dot{\mathbf{x}}=\mathbf{A}\mathbf{x}+\mathbf{B}\mathbf{u}$）\textbf{完全状态可控}，若存在有限时间与输入序列，能把状态从\emph{任意}初态 $\mathbf{x}_0$ 转移到\emph{任意}终态 $\mathbf{z}$。
\end{definition}

\paragraph{Kalman 秩判据（离散代数推导）。}
由 $\mathbf{x}_{k+1}=\mathbf{A}\mathbf{x}_k+\mathbf{B}\mathbf{u}_k$ 迭代 $n$ 步\cite{rawlings2020mpc}：
\begin{equation}
  \mathbf{x}_n=\mathbf{A}^n\mathbf{x}_0+\sum_{i=0}^{n-1}\mathbf{A}^{n-1-i}\mathbf{B}\mathbf{u}_i
  =\mathbf{A}^n\mathbf{x}_0+\begin{bmatrix}\mathbf{B}&\mathbf{A}\mathbf{B}&\cdots&\mathbf{A}^{n-1}\mathbf{B}\end{bmatrix}
  \begin{bmatrix}\mathbf{u}_{n-1}\\\mathbf{u}_{n-2}\\\vdots\\\mathbf{u}_0\end{bmatrix}.
  \label{eq:ctrl-reach}
\end{equation}
要对任意 $\mathbf{z}$ 解出输入，需 $\mathbf{z}-\mathbf{A}^n\mathbf{x}_0$ 落在右端矩阵的列空间，即对任意右端有解 $\iff$ 该矩阵行满秩。由 Cayley--Hamilton，$\mathbf{A}^k\ (k\ge n)$ 可由 $\{\mathbf{A}^0,\dots,\mathbf{A}^{n-1}\}$ 线性表出，故 $n$ 步已捕获所有可达方向——\emph{若 $n$ 步内不可达，任意步数皆不可达}。定义
\begin{equation}
  \boxed{\mathcal{C}=\begin{bmatrix}\mathbf{B}&\mathbf{A}\mathbf{B}&\mathbf{A}^2\mathbf{B}&\cdots&\mathbf{A}^{n-1}\mathbf{B}\end{bmatrix}\in\mathbb{R}^{n\times nm},}
  \label{eq:ctrl-cmat}
\end{equation}
则 $(\mathbf{A},\mathbf{B})$ 可控 $\iff\operatorname{rank}\mathcal{C}=n$（\textbf{可控性矩阵满行秩}）。

\paragraph{可控性 Gramian。}
连续时不变情形，无限时域可控性 Gramian $\mathbf{W}_c\succeq0$ 满足 Lyapunov 方程\cite{rawlings2020mpc}
\begin{equation}
  \mathbf{A}\mathbf{W}_c+\mathbf{W}_c\mathbf{A}^\top=-\mathbf{B}\mathbf{B}^\top,
  \label{eq:ctrl-gramian}
\end{equation}
可控 $\iff\mathbf{W}_c\succ0$。有限时域积分形式 $\mathbf{W}_c(t_0,t_1)=\int_{t_0}^{t_1}\boldsymbol{\phi}(t_0,t)\mathbf{B}\mathbf{B}^\top\boldsymbol{\phi}(t_0,t)^\top\,\mathrm dt$（$\boldsymbol{\phi}$ 为状态转移矩阵）。

\begin{theorem}[Gramian 正定 $\iff$ 秩满]
\label{thm:ctrl-gramian-rank}
系统可控 $\iff$ 可控性 Gramian 正定 $\iff$ 可控性矩阵 \cref{eq:ctrl-cmat} 满秩。
\end{theorem}
\begin{proof}\rebuilt{（标准双向证明，据 \cite{rawlings2020mpc} 与可控性理论）}
($\Rightarrow$ 秩满 $\Rightarrow\mathbf{W}_c\succ0$，逆否法) 设 $\mathbf{W}_c$ 非正定，则存在 $\mathbf{v}\neq\mathbf{0}$ 使
\[
  \mathbf{v}^\top\mathbf{W}_c\mathbf{v}=\int_{t_0}^{t_1}\big\|\mathbf{B}^\top\boldsymbol{\phi}(t_0,t)^\top\mathbf{v}\big\|^2\,\mathrm dt=0\;\Rightarrow\;\mathbf{B}^\top\boldsymbol{\phi}(t_0,t)^\top\mathbf{v}\equiv\mathbf{0}.
\]
在 $t=t_0$ 处对 $t$ 反复求导（用 $\dot{\boldsymbol{\phi}}=\mathbf{A}\boldsymbol{\phi}$，故每求导一次多出一个 $\mathbf{A}^\top$）并令 $t=t_0$，得 $\mathbf{v}^\top\mathbf{B}=\mathbf{v}^\top\mathbf{A}\mathbf{B}=\cdots=\mathbf{v}^\top\mathbf{A}^{n-1}\mathbf{B}=\mathbf{0}$，即 $\mathbf{v}^\top\mathcal{C}=\mathbf{0}$，$\mathcal{C}$ 不满秩。逆否即得满秩 $\Rightarrow\mathbf{W}_c\succ0$。

($\Leftarrow$ $\mathbf{W}_c\succ0\Rightarrow$ 可控，构造性) 取输入
\[
  \mathbf{u}(t)=\mathbf{B}^\top\boldsymbol{\phi}(t_0,t)^\top\mathbf{W}_c^{-1}\big(\mathbf{x}_1-\boldsymbol{\phi}(t_0,t_1)\mathbf{x}_0\big),
\]
代入解的积分表达式可直接验证它把 $\mathbf{x}_0$ 在 $[t_0,t_1]$ 驱到 $\mathbf{x}_1$，故任意态可达，系统可控。
\end{proof}

\paragraph{PBH 判据与可镇定。}
Popov--Belevitch--Hautus（PBH）判据\cite{rawlings2020mpc}：可控 $\iff$
\begin{equation}
  \operatorname{rank}\begin{bmatrix}\lambda\mathbf{I}-\mathbf{A}&\mathbf{B}\end{bmatrix}=n,\quad\forall\lambda\in\mathbb{C}
  \label{eq:ctrl-pbh}
\end{equation}
（实际只需对 $\mathbf{A}$ 的特征值 $\lambda$ 检验，因非特征值处前 $n$ 列已满秩）。\textbf{可镇定（stabilizable）}是松弛版：只要求\emph{不稳定}模态可控，即 \cref{eq:ctrl-pbh} 仅对 $\operatorname{Re}\lambda\ge0$（连续）或 $|\lambda|\ge1$（离散）的 $\lambda$ 成立。LQR 的解存在性只需可镇定（见下节）。

\subsection{可观性与对偶}
\label{sec:ctrl-observability}

\begin{definition}[可观]
\label{def:ctrl-observable}
$(\mathbf{A},\mathbf{C})$ \textbf{可观}，若能由有限段输出（与输入）唯一推断内部状态。
\end{definition}

可观性矩阵与秩判据\cite{rawlings2020mpc}：
\begin{equation}
  \boxed{\mathcal{O}=\begin{bmatrix}\mathbf{C}\\\mathbf{C}\mathbf{A}\\\vdots\\\mathbf{C}\mathbf{A}^{n-1}\end{bmatrix}\in\mathbb{R}^{pn\times n},\qquad\text{可观}\iff\operatorname{rank}\mathcal{O}=n.}
  \label{eq:ctrl-omat}
\end{equation}
不可观子空间 $\mathcal{N}=\bigcap_{k=0}^{n-1}\ker(\mathbf{C}\mathbf{A}^k)=\ker(\mathcal{O})$，可观 $\iff\mathcal{N}=\{\mathbf{0}\}$。

\begin{theorem}[可控--可观对偶]
\label{thm:ctrl-duality}
$(\mathbf{A},\mathbf{B})$ 可控 $\iff(\mathbf{A}^\top,\mathbf{B}^\top)$ 可观；等价地 $(\mathbf{A},\mathbf{C})$ 可观 $\iff(\mathbf{A}^\top,\mathbf{C}^\top)$ 可控\cite{rawlings2020mpc}。
\end{theorem}

\begin{insight}
对偶是把可控性结论“翻译”成可观性的省力工具：把 $\mathbf{A}\to\mathbf{A}^\top,\ \mathbf{B}\to\mathbf{C}^\top$ 一替换，可控性矩阵 \cref{eq:ctrl-cmat} 就变成可观性矩阵 \cref{eq:ctrl-omat}$^\top$、可控 Gramian 变成可观 Gramian、PBH 上下叠放。这对偶在 \cref{ch:eskf} 也现身：Kalman 滤波的可观性正是这里 $(\mathbf{A},\mathbf{C})$ 的可观性，而 LQR 增益与 Kalman 增益由\emph{对偶}的两个 Riccati 给出（\cref{sec:ctrl-lqe-dual}）。\textbf{边界}：对偶是\emph{数学}上的镜像，物理意义相反（一个是“能否驱动”，一个是“能否观测”），别把两者混为一谈。
\end{insight}

可观 Gramian 满足 $\mathbf{A}^\top\mathbf{W}_o+\mathbf{W}_o\mathbf{A}=-\mathbf{C}^\top\mathbf{C}$；PBH 可观判据 $\operatorname{rank}\big[\begin{smallmatrix}\lambda\mathbf{I}-\mathbf{A}\\\mathbf{C}\end{smallmatrix}\big]=n,\ \forall\lambda$。\textbf{可检测（detectable）}：只要求不稳定模态可观（对偶于可镇定）。任意 LTI 系统可经相似变换做 \textbf{Kalman 分解}为四块（可控可观/可控不可观/不可控可观/不可控不可观），传递函数只反映可控且可观部分。

\subsection{状态反馈与极点配置}
\label{sec:ctrl-pole-placement}

\rebuilt{若 $(\mathbf{A},\mathbf{B})$ 可控，则状态反馈 $\mathbf{u}=-\mathbf{K}\mathbf{x}$ 可把闭环 $\dot{\mathbf{x}}=(\mathbf{A}-\mathbf{B}\mathbf{K})\mathbf{x}$ 的极点\emph{任意}配置（极点配置定理；可控性是充要条件）。} 单输入情形 $\mathbf{K}$ 由 \textbf{Ackermann 公式} $\mathbf{K}=[0\,\cdots\,0\ 1]\,\mathcal{C}^{-1}\,p_d(\mathbf{A})$ 给出（$p_d$ 为期望特征多项式，$\mathcal{C}$ 可控性矩阵 \cref{eq:ctrl-cmat}），可控性保证 $\mathcal{C}$ 可逆。

\paragraph{状态观测器（Luenberger）。}
状态反馈 $\mathbf{u}=-\mathbf{K}\mathbf{x}$ 需要\emph{全状态} $\mathbf{x}$，但实际只测得输出 $\mathbf{y}=\mathbf{C}\mathbf{x}$。\textbf{Luenberger 观测器}用一个“模型 + 输出纠错”的副本重构状态\cite{astrom2008feedback}：
\begin{equation}
  \dot{\hat{\mathbf{x}}}=\underbrace{\mathbf{A}\hat{\mathbf{x}}+\mathbf{B}\mathbf{u}}_{\text{模型预测}}+\underbrace{\mathbf{L}(\mathbf{y}-\mathbf{C}\hat{\mathbf{x}})}_{\text{输出残差纠正}},
  \label{eq:ctrl-luenberger}
\end{equation}
$\mathbf{L}$ 为观测器增益。定义估计误差 $\tilde{\mathbf{x}}=\mathbf{x}-\hat{\mathbf{x}}$，由 \cref{eq:ctrl-luenberger} 减真实动力学 $\dot{\mathbf{x}}=\mathbf{A}\mathbf{x}+\mathbf{B}\mathbf{u}$（注意 $\mathbf{B}\mathbf{u}$ 项与输入相消）：
\begin{equation}
  \boxed{\dot{\tilde{\mathbf{x}}}=(\mathbf{A}-\mathbf{L}\mathbf{C})\tilde{\mathbf{x}}.}
  \label{eq:ctrl-obs-error}
\end{equation}
\emph{误差动态不依赖输入}，只由 $(\mathbf{A}-\mathbf{L}\mathbf{C})$ 决定。若 $(\mathbf{A},\mathbf{C})$ 可观，由对偶（\cref{thm:ctrl-duality}）$(\mathbf{A}^\top,\mathbf{C}^\top)$ 可控，故可经对 $\mathbf{A}^\top-\mathbf{C}^\top\mathbf{L}^\top$ 的极点配置任意配 $(\mathbf{A}-\mathbf{L}\mathbf{C})$ 的极点——把误差衰减率调到任意快。

\begin{note}
\textbf{Luenberger 观测器与 Kalman 滤波同形.}\;
\cref{eq:ctrl-luenberger} 与 \cref{ch:eskf} 的 Kalman 滤波\emph{结构完全一致}：“预测（用模型）+ 校正（用观测残差 $\mathbf{y}-\mathbf{C}\hat{\mathbf{x}}$ 乘增益）”。差别仅在增益的\emph{来路}——Luenberger 由\emph{极点配置}手选 $\mathbf{L}$，Kalman 由\emph{噪声协方差经 Riccati} 最优给出 $\mathbf{L}=\mathbf{K}_{\rm Kalman}$（\cref{sec:ctrl-lqe-dual}）。换言之：\textbf{Kalman 滤波 = 最优的 Luenberger 观测器}，正如 LQR = 最优的极点配置（状态反馈）。这是“估计--控制对偶”在结构层面的最直接体现。
\end{note}

\paragraph{分离原理。}
把估计的状态喂给状态反馈：$\mathbf{u}=-\mathbf{K}\hat{\mathbf{x}}=-\mathbf{K}(\mathbf{x}-\tilde{\mathbf{x}})$。这就是\emph{基于观测器的输出反馈}。它的闭环稳定性由下面的优美结果保证。

\begin{proposition}[分离原理]
\label{prop:ctrl-separation}
对 LTI 系统，用观测器估计的状态做反馈 $\mathbf{u}=-\mathbf{K}\hat{\mathbf{x}}$，则闭环（合并 $\mathbf{x}$ 与 $\tilde{\mathbf{x}}$）的系统矩阵分块上三角，其特征值是\emph{控制器极点} $\mathrm{eig}(\mathbf{A}-\mathbf{B}\mathbf{K})$ 与\emph{观测器极点} $\mathrm{eig}(\mathbf{A}-\mathbf{L}\mathbf{C})$ 之\textbf{并}。故 $\mathbf{K}$ 与 $\mathbf{L}$ 可\emph{独立设计}\cite{astrom2008feedback}。
\end{proposition}
\begin{proof}\rebuilt{（标准分块对角化，据 \cite{astrom2008feedback}）}
取状态 $(\mathbf{x},\tilde{\mathbf{x}})$。由 $\mathbf{u}=-\mathbf{K}(\mathbf{x}-\tilde{\mathbf{x}})$ 代入 $\dot{\mathbf{x}}=\mathbf{A}\mathbf{x}+\mathbf{B}\mathbf{u}$ 得 $\dot{\mathbf{x}}=(\mathbf{A}-\mathbf{B}\mathbf{K})\mathbf{x}+\mathbf{B}\mathbf{K}\tilde{\mathbf{x}}$；误差动态由 \cref{eq:ctrl-obs-error} 为 $\dot{\tilde{\mathbf{x}}}=(\mathbf{A}-\mathbf{L}\mathbf{C})\tilde{\mathbf{x}}$。合并：
\[
  \begin{bmatrix}\dot{\mathbf{x}}\\\dot{\tilde{\mathbf{x}}}\end{bmatrix}
  =\begin{bmatrix}\mathbf{A}-\mathbf{B}\mathbf{K}&\mathbf{B}\mathbf{K}\\\mathbf{0}&\mathbf{A}-\mathbf{L}\mathbf{C}\end{bmatrix}
  \begin{bmatrix}\mathbf{x}\\\tilde{\mathbf{x}}\end{bmatrix}.
\]
块上三角矩阵的特征值等于两对角块特征值之并，故闭环极点 $=\mathrm{eig}(\mathbf{A}-\mathbf{B}\mathbf{K})\cup\mathrm{eig}(\mathbf{A}-\mathbf{L}\mathbf{C})$，二者互不影响。
\end{proof}

\begin{insight}
分离原理是“先估计、再控制”这一工程范式的理论许可证：它说\emph{可以}把状态估计器（\cref{ch:eskf} 的 Kalman/观测器）和状态反馈控制器（本节/LQR）\emph{分开设计}，拼起来仍稳定。把“最优观测器（Kalman/LQE）”与“最优状态反馈（LQR）”按此原理组合，即得 \textbf{LQG（线性二次高斯）}最优输出反馈控制器（\cref{sec:ctrl-lqe-dual}）。\textbf{边界（重要）}：分离原理的“极点可分别配”对 LTI 严格成立，但 LQG 的\emph{鲁棒裕度不可分离}——见 \cref{sec:ctrl-margins} 的 note（Doyle 反例）；非线性/受约束系统也只在局部近似成立。
\end{insight}

对偶地，可观系统可设计观测器（增益 $\mathbf{L}$）任意配置误差动态 $(\mathbf{A}-\mathbf{L}\mathbf{C})$ 的极点。极点配置要人工选极点，LQR 则给出 $\mathbf{K}$ 的\emph{最优}选择，是下一节的主题。

\begin{pitfall}[证明陷阱：把“可控”当成“能任意快地驱动”]
可控性只保证能在\emph{有限}时间到达任意状态，\textbf{不}保证能任意快或用有限大的输入到达。极点配置理论上能把极点配到任意远（任意快），但实际受执行器饱和与模型有效频率上限约束——配得越快，所需控制量越大、对模型误差越敏感。正确理解：可控性是\emph{定性}的“能否”，速度/代价是\emph{定量}问题，由 LQR 的 $\mathbf{Q},\mathbf{R}$ 权衡（\cref{sec:ctrl-lqr}）。\textbf{观测器同理}：误差极点配得太快会放大测量噪声（高 $\mathbf{L}$ 把噪声灌进 $\hat{\mathbf{x}}$）——这正是 Kalman 滤波用噪声协方差\emph{最优折中}速度与噪声、而非一味求快的原因。
\end{pitfall}

\begin{practice}
\begin{enumerate}
  \item （判别）对 $\mathbf{A}=\big[\begin{smallmatrix}0&1\\0&0\end{smallmatrix}\big],\ \mathbf{B}=\big[\begin{smallmatrix}0\\1\end{smallmatrix}\big]$（双积分器），写出 $\mathcal{C}$ 并验证满秩、系统可控。这正是机器人“位置+速度”的最简模型。
  \item （对偶）对上题取 $\mathbf{C}=[1\ 0]$（只测位置），写出 $\mathcal{O}$，验证可观；并用 \cref{thm:ctrl-duality} 说明它等价于 $(\mathbf{A}^\top,\mathbf{C}^\top)$ 可控。
  \item （跨章综合）回顾 \cref{ch:eskf} 中“可观性决定哪些状态能被估计”，与本节“可控性决定哪些状态能被驱动”对照，举一个机器人例子说明“可控但不可观”会带来什么后果。
\end{enumerate}
\end{practice}

% ════════════════════════════════════════════════════════════════
\section{线性二次调节器 LQR 与 Riccati 方程}
\label{sec:ctrl-lqr}

\paragraph{动机：让“最优”自动给出反馈增益。}
极点配置要人工选极点，但“极点放哪最好”往往不清楚。LQR 换个问法：给定一个\emph{二次代价}（同时罚状态偏差与控制能量，与 \cref{eq:ctrl-perf-integral} 的 ISE 判据同源——惟 LQR 还罚控制能量），求使代价最小的反馈律。答案是一个优雅的线性反馈 $\mathbf{u}=-\mathbf{K}\mathbf{x}$，其中 $\mathbf{K}$ 由一个矩阵方程——\textbf{Riccati 方程}——给出。LQR 是现代最优控制的基石，也是 MPC 的无约束无限时域特例，其 Riccati 与 Kalman 滤波的 Riccati 对偶。本节给\emph{两条}完整推导（离散动态规划、连续 HJB），并用 Kalman 的反例揭示“最优 $\ne$ 稳定”这一深刻教训。

\subsection{离散有限时域 LQR：动态规划完整推导}
\label{sec:ctrl-lqr-dp}

\paragraph{问题设置。}
离散 LTI 系统 $\mathbf{x}_{t+1}=\mathbf{A}\mathbf{x}_t+\mathbf{B}\mathbf{u}_t$，二次代价\cite{bansal2017lqr}
\begin{equation}
  J(U,\mathbf{x}_0)=\sum_{\tau=0}^{N-1}\big(\mathbf{x}_\tau^\top\mathbf{Q}\mathbf{x}_\tau+\mathbf{u}_\tau^\top\mathbf{R}\mathbf{u}_\tau\big)+\mathbf{x}_N^\top\mathbf{Q}_f\mathbf{x}_N,
  \label{eq:ctrl-lqr-cost}
\end{equation}
$\mathbf{Q},\mathbf{Q}_f\succeq0$（罚状态偏差/终端态），$\mathbf{R}\succ0$（罚控制，且保唯一），$U=(\mathbf{u}_0,\dots,\mathbf{u}_{N-1})$。目标 $\min_U J$。

\paragraph{动态规划与 Bellman 递推。}
定义从 $t$ 时刻、状态 $\mathbf{x}_t$ 起的最优\emph{cost-to-go}（待付代价）$J_t^*(\mathbf{x}_t)=\min_{U_t}\sum_{\tau=t}^{N-1}(\cdots)+\mathbf{x}_N^\top\mathbf{Q}_f\mathbf{x}_N$。由 Bellman 最优性原理，$\mathbf{u}_t$ 同时影响当前代价与下一状态的 cost-to-go\cite{bansal2017lqr}：
\begin{equation}
  J_t^*(\mathbf{x}_t)=\min_{\mathbf{u}_t}\Big[\mathbf{x}_t^\top\mathbf{Q}\mathbf{x}_t+\mathbf{u}_t^\top\mathbf{R}\mathbf{u}_t+J_{t+1}^*(\mathbf{A}\mathbf{x}_t+\mathbf{B}\mathbf{u}_t)\Big].
  \label{eq:ctrl-bellman}
\end{equation}

\begin{theorem}[离散 LQR 的最优 cost-to-go 与反馈律]
\label{thm:ctrl-lqr-dp}
\cref{eq:ctrl-lqr-cost} 的最优 cost-to-go 与最优控制为
\begin{equation}
  J_t^*(\mathbf{z})=\mathbf{z}^\top\mathbf{P}_t\mathbf{z},\qquad \mathbf{u}_t^*=-\mathbf{K}_t\mathbf{z},
\end{equation}
其中 $\mathbf{P}_N=\mathbf{Q}_f$，且对 $t=N-1,\dots,0$ 后向递推（\textbf{离散 Riccati 递推}）：
\begin{align}
  \mathbf{K}_t&=(\mathbf{R}+\mathbf{B}^\top\mathbf{P}_{t+1}\mathbf{B})^{-1}\mathbf{B}^\top\mathbf{P}_{t+1}\mathbf{A},\label{eq:ctrl-dlqr-K}\\
  \mathbf{P}_t&=\mathbf{Q}+\mathbf{A}^\top\mathbf{P}_{t+1}\mathbf{A}-\mathbf{A}^\top\mathbf{P}_{t+1}\mathbf{B}(\mathbf{R}+\mathbf{B}^\top\mathbf{P}_{t+1}\mathbf{B})^{-1}\mathbf{B}^\top\mathbf{P}_{t+1}\mathbf{A}.\label{eq:ctrl-dlqr-P}
\end{align}
\end{theorem}
\begin{proof}\rebuilt{（归纳，逐步，据 \cite{bansal2017lqr}）}
\emph{基例 $t=N$}：代价 \cref{eq:ctrl-lqr-cost} 中 $t=N$ 项 $\mathbf{x}_N^\top\mathbf{Q}_f\mathbf{x}_N$ 与控制无关，故 $J_N^*(\mathbf{z})=\mathbf{z}^\top\mathbf{Q}_f\mathbf{z}$，即 $\mathbf{P}_N=\mathbf{Q}_f$。

\emph{归纳步}：设 $J_{k}^*(\mathbf{x})=\mathbf{x}^\top\mathbf{P}_k\mathbf{x}$，证 $J_{k-1}^*$ 同形。由 \cref{eq:ctrl-bellman}：
\begin{equation}
  J_{k-1}^*(\mathbf{z})=\min_{\mathbf{u}}\Big[\mathbf{z}^\top\mathbf{Q}\mathbf{z}+\mathbf{u}^\top\mathbf{R}\mathbf{u}+(\mathbf{A}\mathbf{z}+\mathbf{B}\mathbf{u})^\top\mathbf{P}_k(\mathbf{A}\mathbf{z}+\mathbf{B}\mathbf{u})\Big].
  \label{eq:ctrl-dp-step}
\end{equation}
对 $\mathbf{u}$ 求梯度置零（用 $\mathbf{P}_k=\mathbf{P}_k^\top$）：
\begin{equation}
  2\mathbf{R}\mathbf{u}+2\mathbf{B}^\top\mathbf{P}_k(\mathbf{A}\mathbf{z}+\mathbf{B}\mathbf{u})=\mathbf{0}\;\Rightarrow\;(\mathbf{R}+\mathbf{B}^\top\mathbf{P}_k\mathbf{B})\mathbf{u}=-\mathbf{B}^\top\mathbf{P}_k\mathbf{A}\mathbf{z},
\end{equation}
因 $\mathbf{R}\succ0\Rightarrow\mathbf{R}+\mathbf{B}^\top\mathbf{P}_k\mathbf{B}\succ0$ 可逆，解出
\begin{equation}
  \mathbf{u}^*=-(\mathbf{R}+\mathbf{B}^\top\mathbf{P}_k\mathbf{B})^{-1}\mathbf{B}^\top\mathbf{P}_k\mathbf{A}\,\mathbf{z}=-\mathbf{K}_{k-1}\mathbf{z}.
\end{equation}
二阶条件：Hessian $2(\mathbf{R}+\mathbf{B}^\top\mathbf{P}_k\mathbf{B})\succ0$，确为极小。回代 \cref{eq:ctrl-dp-step}，记 $\mathbf{A}\mathbf{z}+\mathbf{B}\mathbf{u}^*=(\mathbf{A}-\mathbf{B}\mathbf{K}_{k-1})\mathbf{z}$：
\begin{align}
  J_{k-1}^*(\mathbf{z})
  &=\mathbf{z}^\top\mathbf{Q}\mathbf{z}+\mathbf{z}^\top\mathbf{K}_{k-1}^\top\mathbf{R}\mathbf{K}_{k-1}\mathbf{z}+\mathbf{z}^\top(\mathbf{A}-\mathbf{B}\mathbf{K}_{k-1})^\top\mathbf{P}_k(\mathbf{A}-\mathbf{B}\mathbf{K}_{k-1})\mathbf{z}\notag\\
  &=\mathbf{z}^\top\underbrace{\Big[\mathbf{Q}+\mathbf{K}_{k-1}^\top\mathbf{R}\mathbf{K}_{k-1}+(\mathbf{A}-\mathbf{B}\mathbf{K}_{k-1})^\top\mathbf{P}_k(\mathbf{A}-\mathbf{B}\mathbf{K}_{k-1})\Big]}_{=:\,\mathbf{P}_{k-1}}\mathbf{z}.
  \label{eq:ctrl-dlqr-joseph}
\end{align}
这就是 $J_{k-1}^*(\mathbf{z})=\mathbf{z}^\top\mathbf{P}_{k-1}\mathbf{z}$，完成归纳。\cref{eq:ctrl-dlqr-joseph} 称 Riccati 递推的 \textbf{Joseph（稳定化）形式}，把 \cref{eq:ctrl-dlqr-K} 代入展开、消去 $\mathbf{K}_{k-1}$ 即得紧凑形式 \cref{eq:ctrl-dlqr-P}（二者代数恒等；Joseph 形式数值更稳，自动保对称半正定，紧凑式更省算）。
\end{proof}

\begin{insight}
两点本质：(i) 最优控制是状态的\emph{线性}函数 $\mathbf{u}_t^*=-\mathbf{K}_t\mathbf{x}_t$（线性状态反馈），最优 cost-to-go 是\emph{二次}型——这正是“线性二次”名称的来源；(ii) $\mathbf{P}_t,\mathbf{K}_t$ 从终端 $t=N$ \emph{后向}递推，与系统状态\emph{前向}演化方向相反。这种“后向算增益、前向跑状态”的结构在 MPC 的 Riccati 求解（\cref{sec:ctrl-mpc-numerical}）里原样重现。
\end{insight}

\paragraph{仿射系统的增广技巧。}
若动力学含常数项 $\mathbf{x}_{t+1}=\mathbf{A}\mathbf{x}_t+\mathbf{B}\mathbf{u}_t+\mathbf{c}$（如重力、参考前馈），增广状态 $\mathbf{z}_t=[\mathbf{x}_t;1]$\cite{bansal2017lqr}：
\begin{equation}
  \mathbf{z}_{t+1}=\begin{bmatrix}\mathbf{A}&\mathbf{c}\\\mathbf{0}&1\end{bmatrix}\mathbf{z}_t+\begin{bmatrix}\mathbf{B}\\\mathbf{0}\end{bmatrix}\mathbf{u}_t,
  \label{eq:ctrl-affine}
\end{equation}
即化回标准 LQR 形式。这一技巧在 \cref{sec:ctrl-legged} 的 Cheetah 凸 MPC 里用来把重力并入状态。

\subsection{无限时域 LQR、DARE 与“最优 $\ne$ 稳定”}
\label{sec:ctrl-lqr-dare}

$N\to\infty$ 时 Riccati 递推 \cref{eq:ctrl-dlqr-P} 收敛到\emph{稳态解} $\mathbf{P}$，满足\textbf{离散代数 Riccati 方程（DARE）}\cite{bansal2017lqr,rawlings2020mpc}：
\begin{equation}
  \boxed{\mathbf{P}=\mathbf{Q}+\mathbf{A}^\top\mathbf{P}\mathbf{A}-\mathbf{A}^\top\mathbf{P}\mathbf{B}(\mathbf{R}+\mathbf{B}^\top\mathbf{P}\mathbf{B})^{-1}\mathbf{B}^\top\mathbf{P}\mathbf{A},}
  \label{eq:ctrl-dare}
\end{equation}
最优为\emph{常增益}状态反馈 $\mathbf{u}_t=-\mathbf{K}\mathbf{x}_t,\ \mathbf{K}=(\mathbf{R}+\mathbf{B}^\top\mathbf{P}\mathbf{B})^{-1}\mathbf{B}^\top\mathbf{P}\mathbf{A}$。

\begin{theorem}[LQR 收敛性，Rawlings--Mayne--Diehl Lemma 1.3]
\label{thm:ctrl-lqr-converge}
若 $(\mathbf{A},\mathbf{B})$ 可控、$\mathbf{Q},\mathbf{R}\succ0$，则无限时域 LQR 给出收敛的闭环系统 $\mathbf{x}_{k+1}=(\mathbf{A}-\mathbf{B}\mathbf{K})\mathbf{x}_k$（即 $\mathbf{A}-\mathbf{B}\mathbf{K}$ 渐近稳定）\cite{rawlings2020mpc}。
\end{theorem}
\begin{proof}（据 \cite{rawlings2020mpc}）
\emph{代价有界}：$(\mathbf{A},\mathbf{B})$ 可控 $\Rightarrow$ 存在 $n$ 步输入把任意 $\mathbf{x}_0$ 驱到 $\mathbf{x}_n=\mathbf{0}$，其后取零控制使代价各项为零，故最优代价有限。代价对 $U$ 严格凸（$\mathbf{R}\succ0$），最优解唯一。
\emph{收敛}：沿最优闭环考察 cost-to-go 序列 $V_k=V^0(\mathbf{x}_k)$，由 DP 有 $V_k=V_{k+1}+\big(\mathbf{x}_k^\top\mathbf{Q}\mathbf{x}_k+\mathbf{u}_k^\top\mathbf{R}\mathbf{u}_k\big)$。故 $\{V_k\}$ 非增且有下界 0，收敛，从而增量 $\to0$：$\mathbf{x}_k^\top\mathbf{Q}\mathbf{x}_k\to0$、$\mathbf{u}_k^\top\mathbf{R}\mathbf{u}_k\to0$。因 $\mathbf{Q},\mathbf{R}\succ0$，得 $\mathbf{x}_k\to\mathbf{0},\ \mathbf{u}_k\to\mathbf{0}$，闭环收敛（对线性系统即渐近稳定）。
\end{proof}

\rebuilt{条件可放宽：可控 $\to$ 可镇定；$\mathbf{Q}\succ0\to\mathbf{Q}\succeq0$ 且 $(\mathbf{A},\mathbf{Q}^{1/2})$ 可检测；$\mathbf{R}\succ0$ 保唯一。此时 DARE 仍有唯一正定镇定解（据 \cite{rawlings2020mpc,khalil2002nonlinear}）。}

\paragraph{Kalman 的警示：有限时域最优未必稳定。}
\begin{example}[Kalman 反例：最优 $\ne$ 稳定]
\label{ex:ctrl-kalman}
\rebuilt{（据 \cite{rawlings2020mpc} 逐字数值）} Kalman（1960）指出：工程界常\emph{默认且错误}地以为“带最优控制律的系统必稳定”。反例：
\[
  \mathbf{A}=\begin{bmatrix}4/3&-2/3\\1&0\end{bmatrix},\ \mathbf{B}=\begin{bmatrix}1\\0\end{bmatrix},\ \mathbf{C}=[-2/3\ \ 1],
\]
其传递函数在 $z=3/2$ 有不稳定零点。取 $\mathbf{Q}=\mathbf{C}^\top\mathbf{C}+0.001\mathbf{I}$（罚输出）、$\mathbf{R}=0.001$（很小，鼓励激进控制），用\emph{有限}时域 $N$ 的首个增益 $\mathbf{K}_0$ 作控制律，则闭环稳定性由 $\mathbf{A}-\mathbf{B}\mathbf{K}_0$ 的特征值决定：
\begin{itemize}
  \item $N=5$：$\operatorname{eig}(\mathbf{A}-\mathbf{B}\mathbf{K}_0)=\{1.307,\ 0.001\}$ —— 有特征值 $>1$，\textbf{闭环不稳定}！
  \item $N=7$：$\{0.989,\ 0.001\}$ —— 勉强稳定；
  \item $N\to\infty$（DARE 解）：$\{0.664,\ 0.001\}$ —— 稳定且有合理裕度。
\end{itemize}
\end{example}

\begin{insight}
\cref{ex:ctrl-kalman} 是本章最重要的教训之一：\textbf{“最优”不蕴含“稳定”}——短时域的最优控制器可能让闭环发散。\emph{无限}时域 LQR 才有稳定性保证（\cref{thm:ctrl-lqr-converge}），因为它把“无穷远的未来代价”也算进去，不会为短期收益牺牲长期稳定。这正是 \cref{sec:ctrl-mpc} 中 MPC 必须用\emph{终端代价 + 终端约束}来“伪造”无限时域、修复稳定性的根本动机。读 MPC 稳定性证明前，请记住这个反例。
\end{insight}

\subsection{连续时间 LQR：HJB 完整推导}
\label{sec:ctrl-lqr-hjb}

\paragraph{HJB 方程。}
连续无限时域：$\dot{\mathbf{x}}=\mathbf{A}\mathbf{x}+\mathbf{B}\mathbf{u}$，$J=\int_0^\infty(\mathbf{x}^\top\mathbf{Q}\mathbf{x}+\mathbf{u}^\top\mathbf{R}\mathbf{u})\,\mathrm dt$。设最优值函数 $J^*(\mathbf{x})$。由 Bellman 原理 $J^*(\mathbf{x}(t))=\min_{\mathbf{u}}[\ell\,\mathrm dt+J^*(\mathbf{x}(t+\mathrm dt))]$，Taylor 展开 $J^*(\mathbf{x}(t+\mathrm dt))\approx J^*+\frac{\partial J^*}{\partial\mathbf{x}}\dot{\mathbf{x}}\,\mathrm dt$，令 $\mathrm dt\to0$（定常无限时域 $\partial J^*/\partial t=0$），得 \textbf{Hamilton--Jacobi--Bellman (HJB)} 方程\cite{tedrake_underactuated}：
\begin{equation}
  0=\min_{\mathbf{u}}\Big[\mathbf{x}^\top\mathbf{Q}\mathbf{x}+\mathbf{u}^\top\mathbf{R}\mathbf{u}+\frac{\partial J^*}{\partial\mathbf{x}}(\mathbf{A}\mathbf{x}+\mathbf{B}\mathbf{u})\Big].
  \label{eq:ctrl-hjb}
\end{equation}

\paragraph{二次型拟设与逐步求解。}
设 $J^*(\mathbf{x})=\mathbf{x}^\top\mathbf{P}\mathbf{x},\ \mathbf{P}=\mathbf{P}^\top\succeq0$，则 $\partial J^*/\partial\mathbf{x}=2\mathbf{x}^\top\mathbf{P}$。代入 \cref{eq:ctrl-hjb} 的方括号，对 $\mathbf{u}$ 求梯度置零：
\begin{equation}
  \frac{\partial}{\partial\mathbf{u}}\big[\mathbf{u}^\top\mathbf{R}\mathbf{u}+2\mathbf{x}^\top\mathbf{P}\mathbf{B}\mathbf{u}\big]=2\mathbf{R}\mathbf{u}+2\mathbf{B}^\top\mathbf{P}\mathbf{x}=\mathbf{0}
  \;\Rightarrow\;\boxed{\mathbf{u}^*=-\mathbf{R}^{-1}\mathbf{B}^\top\mathbf{P}\mathbf{x}=-\mathbf{K}\mathbf{x},\ \mathbf{K}=\mathbf{R}^{-1}\mathbf{B}^\top\mathbf{P}.}
  \label{eq:ctrl-clqr-K}
\end{equation}
二阶条件 $2\mathbf{R}\succ0$ 确为极小。把 $\mathbf{u}^*$ 代回 \cref{eq:ctrl-hjb}，逐项展开（不跳步）：
\begin{itemize}
  \item $\mathbf{u}^{*\top}\mathbf{R}\mathbf{u}^*=(\mathbf{R}^{-1}\mathbf{B}^\top\mathbf{P}\mathbf{x})^\top\mathbf{R}(\mathbf{R}^{-1}\mathbf{B}^\top\mathbf{P}\mathbf{x})=\mathbf{x}^\top\mathbf{P}\mathbf{B}\mathbf{R}^{-1}\mathbf{B}^\top\mathbf{P}\mathbf{x}$（用 $\mathbf{R}=\mathbf{R}^\top$）；
  \item $2\mathbf{x}^\top\mathbf{P}\mathbf{B}\mathbf{u}^*=-2\mathbf{x}^\top\mathbf{P}\mathbf{B}\mathbf{R}^{-1}\mathbf{B}^\top\mathbf{P}\mathbf{x}$；
  \item $2\mathbf{x}^\top\mathbf{P}\mathbf{A}\mathbf{x}=\mathbf{x}^\top(\mathbf{P}\mathbf{A}+\mathbf{A}^\top\mathbf{P})\mathbf{x}$（标量等于其转置，对称化）。
\end{itemize}
合并三项（$\mathbf{P}\mathbf{B}\mathbf{R}^{-1}\mathbf{B}^\top\mathbf{P}$ 的系数为 $1-2=-1$）：
\begin{equation}
  0=\mathbf{x}^\top\big[\mathbf{Q}+\mathbf{P}\mathbf{A}+\mathbf{A}^\top\mathbf{P}-\mathbf{P}\mathbf{B}\mathbf{R}^{-1}\mathbf{B}^\top\mathbf{P}\big]\mathbf{x},\quad\forall\mathbf{x},
\end{equation}
故方括号为零，得\textbf{连续代数 Riccati 方程（CARE）}：
\begin{equation}
  \boxed{\mathbf{P}\mathbf{A}+\mathbf{A}^\top\mathbf{P}-\mathbf{P}\mathbf{B}\mathbf{R}^{-1}\mathbf{B}^\top\mathbf{P}+\mathbf{Q}=\mathbf{0}.}
  \label{eq:ctrl-care}
\end{equation}
系统可镇定（且 $(\mathbf{A},\mathbf{Q}^{1/2})$ 可检测）时存在唯一正定镇定解 $\mathbf{P}\succ0$\cite{tedrake_underactuated}。

\begin{theorem}[LQR 闭环稳定（值函数即 Lyapunov 函数）]
\label{thm:ctrl-lqr-lyap}
连续 LQR 闭环 $\dot{\mathbf{x}}=(\mathbf{A}-\mathbf{B}\mathbf{K})\mathbf{x}$ 渐近稳定，$V(\mathbf{x})=\mathbf{x}^\top\mathbf{P}\mathbf{x}$ 即其 Lyapunov 函数。
\end{theorem}
\begin{proof}\rebuilt{（据 \cite{tedrake_underactuated} 与 Lyapunov 理论）}
取 $V=\mathbf{x}^\top\mathbf{P}\mathbf{x}\succ0$。沿闭环求导：
\[
  \dot V=\mathbf{x}^\top\big[(\mathbf{A}-\mathbf{B}\mathbf{K})^\top\mathbf{P}+\mathbf{P}(\mathbf{A}-\mathbf{B}\mathbf{K})\big]\mathbf{x}.
\]
代入 $\mathbf{K}=\mathbf{R}^{-1}\mathbf{B}^\top\mathbf{P}$ 并用 CARE（$\mathbf{P}\mathbf{A}+\mathbf{A}^\top\mathbf{P}=\mathbf{P}\mathbf{B}\mathbf{R}^{-1}\mathbf{B}^\top\mathbf{P}-\mathbf{Q}$）：
\begin{align*}
  (\mathbf{A}-\mathbf{B}\mathbf{K})^\top\mathbf{P}+\mathbf{P}(\mathbf{A}-\mathbf{B}\mathbf{K})
  &=\mathbf{A}^\top\mathbf{P}+\mathbf{P}\mathbf{A}-2\mathbf{P}\mathbf{B}\mathbf{R}^{-1}\mathbf{B}^\top\mathbf{P}\\
  &=(\mathbf{P}\mathbf{B}\mathbf{R}^{-1}\mathbf{B}^\top\mathbf{P}-\mathbf{Q})-2\mathbf{P}\mathbf{B}\mathbf{R}^{-1}\mathbf{B}^\top\mathbf{P}=-\mathbf{Q}-\mathbf{P}\mathbf{B}\mathbf{R}^{-1}\mathbf{B}^\top\mathbf{P}.
\end{align*}
故 $\dot V=-\mathbf{x}^\top(\mathbf{Q}+\mathbf{P}\mathbf{B}\mathbf{R}^{-1}\mathbf{B}^\top\mathbf{P})\mathbf{x}\le0$。$\mathbf{Q}\succ0$ 时 $\dot V<0$ 负定，直接渐近稳定；$\mathbf{Q}\succeq0$ 且 $(\mathbf{A},\mathbf{Q}^{1/2})$ 可检测时 $\dot V$ 沿非平凡轨迹不恒零，由 LaSalle（\cref{thm:ctrl-lasalle}）得渐近稳定。
\end{proof}

\subsection{有限时域连续 LQR 与 Riccati 微分方程}
\label{sec:ctrl-lqr-rde}

有限时域 $J=\mathbf{x}(t_f)^\top\mathbf{Q}_f\mathbf{x}(t_f)+\int_0^{t_f}(\mathbf{x}^\top\mathbf{Q}\mathbf{x}+\mathbf{u}^\top\mathbf{R}\mathbf{u})\,\mathrm dt$，设\emph{含时}值函数 $J^*(\mathbf{x},t)=\mathbf{x}^\top\mathbf{P}(t)\mathbf{x}$。时变 HJB 含 $\partial J^*/\partial t=\mathbf{x}^\top\dot{\mathbf{P}}(t)\mathbf{x}$ 项，重复 \cref{sec:ctrl-lqr-hjb} 的求解（此时不令 $\partial_tJ^*=0$）得\textbf{连续 Riccati 微分方程（RDE）}\cite{tedrake_underactuated}：
\begin{equation}
  \boxed{-\dot{\mathbf{P}}(t)=\mathbf{P}(t)\mathbf{A}+\mathbf{A}^\top\mathbf{P}(t)-\mathbf{P}(t)\mathbf{B}\mathbf{R}^{-1}\mathbf{B}^\top\mathbf{P}(t)+\mathbf{Q},\qquad\mathbf{P}(t_f)=\mathbf{Q}_f,}
  \label{eq:ctrl-rde}
\end{equation}
从终端 $t_f$ \emph{后向}积分；最优时变控制 $\mathbf{u}^*=-\mathbf{R}^{-1}\mathbf{B}^\top\mathbf{P}(t)\mathbf{x}$。当 $t_f\to\infty$ 且系统可镇定时 $\mathbf{P}(t)\to\mathbf{P}_\infty$（CARE 的稳态解）。

\paragraph{协态（Pontryagin）视角。}
\rebuilt{第二条等价推导用变分/极小值原理（据 \cite{tedrake_underactuated} 与 Pontryagin 框架）。} 引入协态 $\boldsymbol{\lambda}(t)$，Hamilton 函数 $H=\tfrac12(\mathbf{x}^\top\mathbf{Q}\mathbf{x}+\mathbf{u}^\top\mathbf{R}\mathbf{u})+\boldsymbol{\lambda}^\top(\mathbf{A}\mathbf{x}+\mathbf{B}\mathbf{u})$。一阶必要条件：状态方程 $\dot{\mathbf{x}}=\partial H/\partial\boldsymbol{\lambda}$；协态方程 $\dot{\boldsymbol{\lambda}}=-\partial H/\partial\mathbf{x}=-\mathbf{Q}\mathbf{x}-\mathbf{A}^\top\boldsymbol{\lambda}$；平稳条件 $\partial H/\partial\mathbf{u}=\mathbf{R}\mathbf{u}+\mathbf{B}^\top\boldsymbol{\lambda}=\mathbf{0}\Rightarrow\mathbf{u}=-\mathbf{R}^{-1}\mathbf{B}^\top\boldsymbol{\lambda}$；终端 $\boldsymbol{\lambda}(t_f)=\mathbf{Q}_f\mathbf{x}(t_f)$。设“扫描” ansatz $\boldsymbol{\lambda}(t)=\mathbf{P}(t)\mathbf{x}(t)$，代入协态方程并对所有 $\mathbf{x}$ 成立，即重得 RDE \cref{eq:ctrl-rde}——两条路殊途同归。

\subsection{Riccati 方程的求解：Hamilton 矩阵法}
\label{sec:ctrl-riccati-solve}

CARE \cref{eq:ctrl-care} 是 $\mathbf{P}$ 的二次矩阵方程，闭式解经 $2n\times2n$ \textbf{Hamilton 矩阵}\cite{rawlings2020mpc}：
\begin{equation}
  \mathbf{Z}=\begin{bmatrix}\mathbf{A}&-\mathbf{B}\mathbf{R}^{-1}\mathbf{B}^\top\\-\mathbf{Q}&-\mathbf{A}^\top\end{bmatrix}.
  \label{eq:ctrl-hamilton}
\end{equation}
若 $\mathbf{Z}$ 无虚轴特征值，则恰有一半特征值实部为负。取其\emph{稳定不变子空间}（负实部特征值对应的特征/广义特征向量张成的 $n$ 维子空间）的一组基 $\big[\begin{smallmatrix}\mathbf{U}_{11}\\\mathbf{U}_{21}\end{smallmatrix}\big]\in\mathbb{R}^{2n\times n}$，则\textbf{镇定解}
\begin{equation}
  \mathbf{P}=\mathbf{U}_{21}\mathbf{U}_{11}^{-1},
  \label{eq:ctrl-hamilton-sol}
\end{equation}
且闭环 $\mathbf{A}-\mathbf{B}\mathbf{R}^{-1}\mathbf{B}^\top\mathbf{P}$ 的特征值恰为 $\mathbf{Z}$ 的负实部特征值。DARE 类似，用\emph{辛矩阵}并取单位圆内特征值。工程上更常用迭代法（直接迭代 Riccati 递推至收敛）或调库（MATLAB \texttt{care/dare}、SciPy \texttt{solve\_continuous\_are}）。

\begin{algo}[无限时域 LQR 求解]
\label{alg:ctrl-lqr}
输入 $\mathbf{A},\mathbf{B},\mathbf{Q}\succeq0,\mathbf{R}\succ0$；输出常增益 $\mathbf{K}$。
\begin{enumerate}
  \item 检查 $(\mathbf{A},\mathbf{B})$ 可镇定、$(\mathbf{A},\mathbf{Q}^{1/2})$ 可检测（否则无镇定解）；
  \item 解 CARE \cref{eq:ctrl-care}（Hamilton 矩阵 \cref{eq:ctrl-hamilton-sol} 或迭代法）得 $\mathbf{P}\succ0$；
  \item $\mathbf{K}=\mathbf{R}^{-1}\mathbf{B}^\top\mathbf{P}$，控制律 $\mathbf{u}=-\mathbf{K}\mathbf{x}$；
  \item 验证 $\mathbf{A}-\mathbf{B}\mathbf{K}$ 特征值实部均 $<0$。
\end{enumerate}
\end{algo}

\subsection{含交叉项的统一公式与 LQR--Kalman 对偶}
\label{sec:ctrl-lqe-dual}

\paragraph{含交叉项 $\mathbf{S}_{xu}$ 的统一式。}
代价含状态-输入交叉项 $2\mathbf{x}^\top\mathbf{S}_{xu}\mathbf{u}$ 时\cite{tedrake_underactuated}，CARE 推广为 $\mathbf{A}^\top\mathbf{P}+\mathbf{P}\mathbf{A}-(\mathbf{P}\mathbf{B}+\mathbf{S}_{xu})\mathbf{R}^{-1}(\mathbf{B}^\top\mathbf{P}+\mathbf{S}_{xu}^\top)+\mathbf{Q}=\mathbf{0}$，增益 $\mathbf{K}=\mathbf{R}^{-1}(\mathbf{B}^\top\mathbf{P}+\mathbf{S}_{xu}^\top)$。可经 $\boldsymbol{\mathcal{A}}=\mathbf{A}-\mathbf{B}\mathbf{R}^{-1}\mathbf{S}_{xu}^\top,\ \boldsymbol{\mathcal{Q}}=\mathbf{Q}-\mathbf{S}_{xu}\mathbf{R}^{-1}\mathbf{S}_{xu}^\top$ 把交叉项吸收，化回无交叉项标准式。

\begin{note}
\textbf{LQR--Kalman（LQE）对偶与 LQG.}\;
LQR 的 Riccati（\cref{eq:ctrl-dare}）与 \cref{ch:eskf} 的 Kalman 滤波 Riccati（协方差预测-更新）结构相同：一个\emph{后向}（LQR，控制增益），一个\emph{前向}（LQE，估计增益），此即\textbf{估计--控制对偶}。把“最优状态估计（Kalman/LQE）”与“最优状态反馈（LQR）”合并即 \textbf{LQG（线性二次高斯）}控制，\emph{分离原理}保证二者可独立设计。\textbf{再次提醒}：LQR 的 $\mathbf{P}$ 是值函数权重（cost-to-go Hessian），Kalman 的 $\mathbf{P}$ 是状态协方差——同字母、对偶结构、物理意义相反，本书在本章用 $\mathbf{P}$ 专指前者。
\end{note}

\begin{pitfall}[证明陷阱：用有限时域首增益当无限时域增益却不验稳定]
看到 Riccati 递推收敛快，有人直接用某个 $N$ 步的首增益 $\mathbf{K}_0$ 当常增益用，省去解 DARE。\cref{ex:ctrl-kalman} 已证：小 $N$ 的 $\mathbf{K}_0$ 可能让闭环\emph{发散}（特征值出单位圆）。现象：仿真里状态缓慢甚至迅速放大。根本原因：有限时域最优不含“无穷远未来”的代价，不保证稳定。正确做法：要么把 $N$ 取到 Riccati 充分收敛（$\mathbf{P}_t$ 不再变），要么直接解 DARE/CARE 得稳态 $\mathbf{P}$，并\emph{务必}验 $\mathbf{A}-\mathbf{B}\mathbf{K}$ 的特征值。
\end{pitfall}

\begin{practice}
\begin{enumerate}
  \item （推导，在草稿纸上完成）对标量系统 $\dot x=ax+bu$、$J=\int_0^\infty(qx^2+ru^2)\mathrm dt$，由 CARE \cref{eq:ctrl-care} 解出 $P=\frac{r}{b^2}\big(a+\sqrt{a^2+b^2q/r}\big)$（取正根），并算 $K=bP/r$。验证闭环 $a-bK<0$。
  \item （复现 Kalman 反例）对 \cref{ex:ctrl-kalman} 的 $\mathbf{A},\mathbf{B},\mathbf{Q},\mathbf{R}$，手工或用代码迭代 Riccati \cref{eq:ctrl-dlqr-P} 从 $\mathbf{P}=\mathbf{Q}$ 起，分别取 $N=5,7,\infty$ 算 $\mathbf{K}_0$ 与 $\operatorname{eig}(\mathbf{A}-\mathbf{B}\mathbf{K}_0)$，重现表中数值。
  \item （开放）解释为什么增大 $\mathbf{R}$（更罚控制）会使 LQR 增益 $\mathbf{K}$ 变小、闭环极点更靠近开环极点；增大 $\mathbf{Q}$ 又如何？把这与 \cref{tab:ctrl-pid-effect} 的 PID 直觉类比，并说明类比的\emph{边界}（LQR 是多变量最优，PID 是单回路手调）。
\end{enumerate}
\end{practice}

\subsection{非线性推广：iLQR 与 DDP}
\label{sec:ctrl-ilqr-ddp}

\paragraph{动机：把 LQR 用到非线性轨迹优化。}
LQR 的 Riccati 后向递推（\cref{thm:ctrl-lqr-dp}）只对\emph{线性}动力学 + \emph{二次}代价成立。但机器人是强非线性的（\cref{eq:ctrl-manip}）。\textbf{迭代 LQR（iLQR）}与\textbf{微分动态规划（DDP）}的思路：围绕一条标称轨迹反复做“\emph{二次近似 $\to$ 解一次 LQR $\to$ 前向滚动改进轨迹}”，把非线性最优控制化为一串 LQR 子问题\cite{tedrake_underactuated}。它是当代轨迹优化（操作、足式、MPC 的非线性内核）的主力算法，也把本节的 Riccati 与 \cref{ch:planning} 的轨迹优化、\cref{ch:nlopt} 的高斯--牛顿连成一线。

\paragraph{问题与值函数展开。}
非线性离散系统 $\mathbf{x}_{k+1}=f(\mathbf{x}_k,\mathbf{u}_k)$，一般阶段代价 $\ell(\mathbf{x}_k,\mathbf{u}_k)$、终端代价 $\ell_f(\mathbf{x}_N)$。给定当前标称轨迹 $(\bar{\mathbf{x}}_k,\bar{\mathbf{u}}_k)$，记偏差 $\delta\mathbf{x}_k,\delta\mathbf{u}_k$。定义\emph{状态-动作值函数（Q 函数）} $Q_k(\delta\mathbf{x},\delta\mathbf{u})=\ell(\cdot)+V_{k+1}(f(\cdot))$ 围绕标称点的二阶展开\cite{tedrake_underactuated}（下标记偏导，$V_{k+1}',V_{k+1}''$ 为下一步值函数的一/二阶导）：
\begin{align}
  Q_{\mathbf{x}}&=\ell_{\mathbf{x}}+f_{\mathbf{x}}^\top V_{k+1}', &
  Q_{\mathbf{u}}&=\ell_{\mathbf{u}}+f_{\mathbf{u}}^\top V_{k+1}', \notag\\
  Q_{\mathbf{x}\mathbf{x}}&=\ell_{\mathbf{x}\mathbf{x}}+f_{\mathbf{x}}^\top V_{k+1}''f_{\mathbf{x}}\,(+\,V_{k+1}'\!\cdot f_{\mathbf{x}\mathbf{x}}), &
  Q_{\mathbf{u}\mathbf{u}}&=\ell_{\mathbf{u}\mathbf{u}}+f_{\mathbf{u}}^\top V_{k+1}''f_{\mathbf{u}}\,(+\,V_{k+1}'\!\cdot f_{\mathbf{u}\mathbf{u}}), \notag\\
  Q_{\mathbf{u}\mathbf{x}}&=\ell_{\mathbf{u}\mathbf{x}}+f_{\mathbf{u}}^\top V_{k+1}''f_{\mathbf{x}}\,(+\,V_{k+1}'\!\cdot f_{\mathbf{u}\mathbf{x}}). &&
  \label{eq:ctrl-ilqr-Q}
\end{align}

\begin{insight}
\textbf{iLQR vs DDP 的唯一区别}就在 \cref{eq:ctrl-ilqr-Q} 括号里的项：\textbf{DDP} 保留动力学的二阶导 $f_{\mathbf{x}\mathbf{x}},f_{\mathbf{u}\mathbf{u}},f_{\mathbf{u}\mathbf{x}}$（张量，二阶方法），\textbf{iLQR} 丢掉它们（只留代价二阶导，是“$1.5$ 阶”高斯--牛顿型方法）。丢掉张量项省去昂贵的二阶动力学微分、且让 $Q_{\mathbf{u}\mathbf{u}}$ 更易保持正定，故 iLQR 在机器人上更常用；DDP 收敛更快（二次）但每步更贵。这与 \cref{ch:nlopt} 中“高斯--牛顿丢掉残差二阶项 $\approx$ 牛顿法”\emph{完全是同一种近似哲学}。
\end{insight}

\paragraph{后向传播（局部 LQR）。}
对 $\delta\mathbf{u}$ 极小化二次 $Q$（$Q_{\mathbf{u}\mathbf{u}}\succ0$）得\emph{前馈 + 反馈}更新律\cite{tedrake_underactuated}：
\begin{equation}
  \delta\mathbf{u}^*=\mathbf{k}+\mathbf{K}\,\delta\mathbf{x},\qquad
  \mathbf{k}=-Q_{\mathbf{u}\mathbf{u}}^{-1}Q_{\mathbf{u}},\quad \mathbf{K}=-Q_{\mathbf{u}\mathbf{u}}^{-1}Q_{\mathbf{u}\mathbf{x}},
  \label{eq:ctrl-ilqr-gain}
\end{equation}
其中 $\mathbf{k}$ 是前馈项（标称轨迹非最优时非零）、$\mathbf{K}$ 是反馈增益。回代得值函数后向递推（与 LQR Riccati 同构，但围绕轨迹、且多了前馈）：
\begin{equation}
  V_k'=Q_{\mathbf{x}}-\mathbf{K}^\top Q_{\mathbf{u}\mathbf{u}}\mathbf{k},\qquad
  V_k''=Q_{\mathbf{x}\mathbf{x}}-\mathbf{K}^\top Q_{\mathbf{u}\mathbf{u}}\mathbf{K},
  \label{eq:ctrl-ilqr-V}
\end{equation}
从终端 $V_N'=\ell_{f,\mathbf{x}},\ V_N''=\ell_{f,\mathbf{x}\mathbf{x}}$ 起后向算到 $k=0$。

\paragraph{前向滚动（line search）。}
用新律\emph{非线性}前向滚动（这是 DDP/iLQR 与朴素 LQR 的关键差别——不是线性外推，而是真动力学积分）：$\mathbf{u}_k=\bar{\mathbf{u}}_k+\alpha\mathbf{k}_k+\mathbf{K}_k(\mathbf{x}_k-\bar{\mathbf{x}}_k)$，$\mathbf{x}_{k+1}=f(\mathbf{x}_k,\mathbf{u}_k)$，步长 $\alpha\in(0,1]$ 经线搜索保证代价下降；接受后更新标称轨迹，迭代至收敛。

\begin{algo}[iLQR / DDP 一次迭代]
\label{alg:ctrl-ilqr}
输入标称轨迹 $(\bar{\mathbf{x}}_{0:N},\bar{\mathbf{u}}_{0:N-1})$；输出改进轨迹。
\begin{enumerate}
  \item 沿标称轨迹线性化动力学（$f_{\mathbf{x}},f_{\mathbf{u}}$）、二次化代价（$\ell_{\mathbf{x}},\ell_{\mathbf{u}},\ell_{\mathbf{x}\mathbf{x}},\dots$），DDP 另算 $f_{\mathbf{x}\mathbf{x}}$ 等张量；
  \item \textbf{后向}：从 $k=N$ 起按 \cref{eq:ctrl-ilqr-Q,eq:ctrl-ilqr-gain,eq:ctrl-ilqr-V} 递推得各步 $(\mathbf{k}_k,\mathbf{K}_k)$；若 $Q_{\mathbf{u}\mathbf{u}}\not\succ0$ 加正则 $Q_{\mathbf{u}\mathbf{u}}+\mu\mathbf{I}$（信赖域，呼应 \cref{ch:nlopt} 的 LM）；
  \item \textbf{前向}：以 $\mathbf{x}_0$ 起、步长 $\alpha$ 线搜索非线性滚动，得新轨迹并比较代价；
  \item 接受/缩小 $\alpha$，更新标称轨迹；重复至梯度/代价收敛。
\end{enumerate}
\end{algo}

\begin{note}
\textbf{iLQR/DDP 与 Pontryagin、Riccati、MPC 的关系.}\;
(i) 后向传播的 $V_k'$ 在收敛时\emph{就是}协态 $\boldsymbol{\lambda}_k$（呼应 \cref{sec:ctrl-lqr-hjb} 的 Pontryagin 视角），前向滚动即状态积分——故 iLQR 是 Pontryagin 必要条件的一种数值解法。
(ii) 若动力学本就线性、代价二次，则 $\mathbf{k}=\mathbf{0}$、\cref{eq:ctrl-ilqr-V} 退化为 LQR 的 Riccati（\cref{eq:ctrl-dlqr-P}），\emph{一次迭代即收敛}——iLQR 是 LQR 的严格非线性推广。
(iii) 把 iLQR 放进\emph{滚动时域}（每步只执行首个 $\mathbf{u}_0$、重规划）即\textbf{非线性 MPC（NMPC）}的一种主流实现；约束可经控制限 DDP（control-limited DDP）或增广拉格朗日处理。这是 \cref{sec:ctrl-mpc} 在非线性、带约束情形的落地路径之一。
\end{note}

\begin{pitfall}[数值陷阱：iLQR 后向中 $Q_{\mathbf{u}\mathbf{u}}$ 非正定]
直接对 \cref{eq:ctrl-ilqr-gain} 求逆，若 $Q_{\mathbf{u}\mathbf{u}}$ 非正定（远离最优、或 DDP 张量项作怪），增益爆掉、前向滚动发散。现象：代价不降反升、轨迹甩飞。根本原因：二次近似在当前点不是凸的下界。正确做法：加正则 $Q_{\mathbf{u}\mathbf{u}}+\mu\mathbf{I}$（信赖域，$\mu$ 像 LM 阻尼一样自适应增减，见 \cref{ch:nlopt}），并对前向滚动做\emph{线搜索}（缩小 $\alpha$ 直到代价下降）——二者缺一不可。这也是 iLQR 比 DDP 鲁棒的原因之一（iLQR 丢掉张量项，$Q_{\mathbf{u}\mathbf{u}}$ 天然更接近正定）。
\end{pitfall}

% ════════════════════════════════════════════════════════════════
\section{模型预测控制 MPC}
\label{sec:ctrl-mpc}

\paragraph{动机：处理约束的最优控制。}
LQR 给出全局最优反馈，但它\emph{无法处理约束}——真实机器人有关节力矩上限、关节角限位、摩擦锥、避障等硬约束，LQR 的 $\mathbf{u}=-\mathbf{K}\mathbf{x}$ 会无视它们。\textbf{模型预测控制（MPC）}的思路：在每个时刻，以当前状态为初值，在一个\emph{有限}未来时域上解一个\emph{带约束}的最优控制问题，只施加第一步控制，下一时刻用新状态重新求解。这就是\emph{滚动时域}。MPC 是当代机器人（机械臂、足式、自动驾驶）处理约束的主力框架；其无约束、无限时域、线性二次特例\emph{就是} LQR。本节给问题表述、凝聚为 QP、持久可行性与稳定性的完整证明。

\subsection{滚动时域问题与三要素}
\label{sec:ctrl-mpc-formulation}

时刻 $t$ 以当前状态 $\mathbf{x}_0=\mathbf{x}(t)$ 为初值，解有限时域最优控制问题\cite{borrelli2017predictive,rawlings2020mpc}：
\begin{equation}
  \min_{\mathbf{u}_0,\dots,\mathbf{u}_{N-1}}\ J=\sum_{k=0}^{N-1}\ell(\mathbf{x}_k,\mathbf{u}_k)+V_f(\mathbf{x}_N)
  \label{eq:ctrl-mpc-cost}
\end{equation}
\begin{equation}
  \text{s.t.}\quad\mathbf{x}_{k+1}=f(\mathbf{x}_k,\mathbf{u}_k),\ \ \mathbf{x}_k\in\mathcal{X},\ \ \mathbf{u}_k\in\mathcal{U},\ \ \mathbf{x}_N\in\mathcal{X}_f,\ \ \mathbf{x}_0=\mathbf{x}(t),
  \label{eq:ctrl-mpc-constr}
\end{equation}
$N$ 预测时域，$\ell$ 阶段代价，$V_f$ 终端代价，$\mathcal{X},\mathcal{U}$ 状态/输入约束集，$\mathcal{X}_f$ 终端约束集。\textbf{滚动时域原理}：(1) 在 $t$ 解上述问题得 $\mathbf{u}^*=(\mathbf{u}_0^*,\dots)$；(2) \emph{只施加首步} $\kappa_N(\mathbf{x})=\mathbf{u}_0^*$；(3) 采样新状态；(4) 回到 (1)。闭环 $\mathbf{x}^+=f(\mathbf{x},\kappa_N(\mathbf{x}))$。

\textbf{三要素}\cite{borrelli2017predictive}：内部动力学模型、滚动时域代价 $J$、在约束下极小化 $J$ 的优化算法（线性 MPC 为 QP）。\textbf{MPC vs LQR}：LQR 在整个时域上算一次单一解全程使用；MPC 在滚动窗口上反复求解，可处理约束与非线性。\emph{LQR 在滚动时域下反复运行即成为一种 MPC}。线性 MPC 的问题是凸的（QP），非线性 MPC（NMPC）一般非凸（用单/多重打靶或配点法）。

\begin{example}[滚动时域的本质：闭环 $\neq$ 开环预测]
\label{ex:ctrl-mpc-double-int}
\rebuilt{（双积分器，据 \cite{borrelli2017predictive}）} 双积分器，$N=3$，$\mathbf{P}=\mathbf{Q}=\mathbf{I}$，$\mathbf{R}=10$，输入约束 $-0.5\le u\le0.5$，状态约束 $|x_i|\le5$。其关联 QP 的 Hessian
\[
  \mathbf{H}=\begin{bmatrix}13.5&-10&-0.5\\-10&22&-10\\-0.5&-10&31.5\end{bmatrix}.
\]
从 $\mathbf{x}(0)=[-4.5,2]$ 出发闭环收敛且满足约束；但从 $\mathbf{x}(0)=[-4.5,3]$ 出发，在 $\mathbf{x}(2)=[1,2]$ 处问题\emph{不可行}而停止——说明初始可行不保证未来可行，闭环轨迹与开环预测不同（滚动时域本质）。这引出下面的\emph{持久可行性}问题。
\end{example}

\subsection{线性 MPC 凝聚为 QP}
\label{sec:ctrl-mpc-qp}

线性时变（LTV）情形 $\mathbf{x}_{i+1}=\mathbf{A}_i\mathbf{x}_i+\mathbf{B}_i\mathbf{u}_i$，二次代价。把动力学递推展开，状态轨迹 $\mathbf{X}=[\mathbf{x}_1;\dots;\mathbf{x}_k]$ 表为初值与控制序列 $\mathbf{U}=[\mathbf{u}_0;\dots;\mathbf{u}_{k-1}]$ 的仿射函数\cite{dicarlo2018cheetah}：
\begin{equation}
  \mathbf{X}=\mathbf{A}_{qp}\mathbf{x}_0+\mathbf{B}_{qp}\mathbf{U},\quad
  \mathbf{A}_{qp}=\begin{bmatrix}\mathbf{A}\\\mathbf{A}^2\\\vdots\\\mathbf{A}^k\end{bmatrix},\
  \mathbf{B}_{qp}=\begin{bmatrix}\mathbf{B}&&\\\mathbf{A}\mathbf{B}&\mathbf{B}&\\\vdots&&\ddots\\\mathbf{A}^{k-1}\mathbf{B}&\cdots&\mathbf{B}\end{bmatrix}
  \label{eq:ctrl-condense-pred}
\end{equation}
（定常 $\mathbf{A},\mathbf{B}$ 示例）。代价 $\|\mathbf{X}-\mathbf{X}_{\rm ref}\|_{\mathbf{L}}^2+\|\mathbf{U}\|_{\mathbf{K}}^2$（$\mathbf{L}=\mathrm{blkdiag}(\mathbf{Q}_i),\ \mathbf{K}=\mathrm{blkdiag}(\mathbf{R}_i)$）。代入 \cref{eq:ctrl-condense-pred}、记 $\mathbf{y}=\mathbf{X}_{\rm ref}$，展开关于 $\mathbf{U}$ 的二次项 $\mathbf{U}^\top\mathbf{B}_{qp}^\top\mathbf{L}\mathbf{B}_{qp}\mathbf{U}$ 与一次项 $2\mathbf{U}^\top\mathbf{B}_{qp}^\top\mathbf{L}(\mathbf{A}_{qp}\mathbf{x}_0-\mathbf{y})$，加 $\mathbf{U}^\top\mathbf{K}\mathbf{U}$，配成标准 QP\cite{dicarlo2018cheetah}：
\begin{equation}
  \min_{\mathbf{U}}\ \tfrac12\mathbf{U}^\top\mathbf{H}\mathbf{U}+\mathbf{U}^\top\mathbf{g}\quad\text{s.t.}\ \underline{\mathbf{c}}\le\mathbf{C}\mathbf{U}\le\overline{\mathbf{c}},
  \label{eq:ctrl-mpc-qpform}
\end{equation}
\begin{equation}
  \mathbf{H}=2(\mathbf{B}_{qp}^\top\mathbf{L}\mathbf{B}_{qp}+\mathbf{K}),\qquad
  \mathbf{g}=2\mathbf{B}_{qp}^\top\mathbf{L}(\mathbf{A}_{qp}\mathbf{x}_0-\mathbf{y}).
  \label{eq:ctrl-mpc-Hg}
\end{equation}
\textbf{凝聚（condensing）}把状态轨迹变量消去，决策变量只剩 $\mathbf{U}$，$\mathbf{H}$ 维度只依赖输入维与时域、\emph{与状态数无关}——这是把动力学约束移入代价的收益（其数值结构见 \cref{sec:ctrl-mpc-numerical}）。

\subsection{持久可行性}
\label{sec:ctrl-mpc-feasibility}

\paragraph{问题。}
\cref{ex:ctrl-mpc-double-int} 显示初始可行的 $\mathbf{x}(0)$ 不保证未来恒可行。“未来所有时刻都可行”称\textbf{持久可行性}\cite{borrelli2017predictive}。记 $C_\infty$ 最大控制不变集，$O_\infty$ 闭环最大正不变集，$X_0$ 初始可行集，恒有 $O_\infty\subseteq X_0\subseteq C_\infty$。

\begin{definition}[（控制）不变集]
\label{def:ctrl-invariant}
集 $X$ 对 $\mathbf{x}^+=f(\mathbf{x})$ \textbf{正不变}，若 $\mathbf{x}\in X\Rightarrow f(\mathbf{x})\in X$。集 $X$ 对 $\mathbf{x}^+=f(\mathbf{x},\mathbf{u}),\ \mathbf{u}\in\mathcal{U}$ \textbf{控制不变}，若 $\forall\mathbf{x}\in X,\ \exists\mathbf{u}\in\mathcal{U}(\mathbf{x})$ 使 $f(\mathbf{x},\mathbf{u})\in X$。
\end{definition}

\begin{lemma}[持久可行性充要条件]
\label{lem:ctrl-pf1}
RHC 问题持久可行 $\iff X_0=O_\infty$\cite{borrelli2017predictive}。
\end{lemma}
\begin{proof}（据 \cite{borrelli2017predictive}）
持久可行要求 $X_0$ 对闭环正不变。已知 $O_\infty\subseteq X_0$；而正不变集 $X_0$ 不能大于\emph{最大}正不变集 $O_\infty$，故 $X_0=O_\infty$。
\end{proof}

\begin{lemma}[持久可行性充分条件（终端集控制不变）]
\label{lem:ctrl-pf2}
若终端约束集 $\mathcal{X}_f$ 控制不变，则 RHC 持久可行，且 $O_\infty=X_0$ 与 $\mathbf{P},\mathbf{Q},\mathbf{R}$ 无关\cite{borrelli2017predictive}。
\end{lemma}
\begin{proof}（据 \cite{borrelli2017predictive}）
$\mathcal{X}_f$ 控制不变 $\Rightarrow X_1$（一步可达 $\mathcal{X}_f$ 的可行集）控制不变，归纳得 $X_{N-1},\dots,X_1$ 均控制不变。取任意 $\mathbf{x}\in X_0$ 及其可行控制 $\mathbf{u}$，则 $\mathbf{x}^+=f(\mathbf{x},\mathbf{u})\in X_1\subseteq X_0$，故 $X_0$ 对所有可行 $\mathbf{u}$ 正不变；由 \cref{lem:ctrl-pf1} $X_0=O_\infty$，且因对所有可行 $\mathbf{u}$ 成立，$O_\infty$ 不依赖代价参数 $\mathbf{P},\mathbf{Q},\mathbf{R}$。
\end{proof}

\begin{insight}
\cref{lem:ctrl-pf2} 是 MPC 工程设计的金科玉律：\textbf{要保证持久可行，给问题加一个控制不变的终端约束集 $\mathcal{X}_f$}。直觉：$\mathcal{X}_f$ 控制不变意味着“一旦进入 $\mathcal{X}_f$，就总有办法留在里面”，于是把状态在 $N$ 步内驱入 $\mathcal{X}_f$ 的承诺可以\emph{递归地}延续下去，永不丢失可行性。这也是把有限时域“接续”成无限时域的第一步——第二步是用终端代价保证稳定（下节）。
\end{insight}

\subsection{MPC 稳定性主定理与完整证明}
\label{sec:ctrl-mpc-stability}

\paragraph{核心思想：值函数即闭环 Lyapunov 函数。}
回忆 \cref{ex:ctrl-kalman}：有限时域最优未必稳定。MPC 的修复办法是精心选\emph{终端代价 $V_f$ + 终端集 $\mathcal{X}_f$}，使最优值函数 $J_N^0(\cdot)$ 成为闭环 Lyapunov 函数\cite{rawlings2020mpc,borrelli2017predictive}。下面给出完整证明链。

\begin{definition}[MPC 基本稳定性假设]
\label{def:ctrl-mpc-asm}
$\ell,V_f,\mathcal{X}_f$ 满足\cite{rawlings2020mpc,borrelli2017predictive}：
\begin{itemize}
  \item[(A0)] 阶段代价 $\ell$ 与终端代价 $V_f$ 连续且正定（$\ell(\mathbf{0},\mathbf{0})=0,\ V_f(\mathbf{0})=0$）；
  \item[(A1)] $\mathcal{X},\mathcal{X}_f,\mathcal{U}$ 闭且内部含原点；
  \item[(A2)] $\mathcal{X}_f$ 控制不变，$\mathcal{X}_f\subseteq\mathcal{X}$；
  \item[(A3)（终端代价下降 / CLF 条件）] 对所有 $\mathbf{x}\in\mathcal{X}_f$ 存在 $\mathbf{u}$（使 $(\mathbf{x},\mathbf{u})$ 满足约束、$f(\mathbf{x},\mathbf{u})\in\mathcal{X}_f$）满足
  \begin{equation}
    V_f(f(\mathbf{x},\mathbf{u}))-V_f(\mathbf{x})\le-\ell(\mathbf{x},\mathbf{u}).
    \label{eq:ctrl-mpc-A3}
  \end{equation}
\end{itemize}
\end{definition}

\begin{insight}
(A3) 是全部的关键：它说在终端集内存在一个“局部控制器 $\kappa_f$”使终端代价 $V_f$ 像\emph{无限时域 cost-to-go} 一样沿轨迹\emph{下降至少 $\ell$}——即 $V_f$ 是终端集内闭环的\textbf{控制 Lyapunov 函数（CLF）}。这正是把有限时域“伪造”成无限时域、修复 \cref{ex:ctrl-kalman} 那种不稳定的机理。\textbf{标准落地}：取 $V_f(\mathbf{x})=\mathbf{x}^\top\mathbf{P}\mathbf{x}$（$\mathbf{P}$ 为局部 LQR 的 DARE 解）、$\kappa_f=-\mathbf{K}_{\rm LQR}\mathbf{x}$、$\mathcal{X}_f$ 取该 LQR 的约束容许不变椭球，则 (A3) 取等号，\emph{MPC 在终端集内退化为 LQR}。或终端等式约束 $\mathcal{X}_f=\{\mathbf{0}\},V_f=0$（最简）。这就是连接 MPC 与 LQR 的标准桥。
\end{insight}

下面用动态规划值函数的\emph{单调性}导出下降性。记 $V_j^0(\mathbf{x})$ 为时域 $j$ 的最优值函数，DP 递推 $V_j^0(\mathbf{x})=\min_{\mathbf{u}}\{\ell(\mathbf{x},\mathbf{u})+V_{j-1}^0(f(\mathbf{x},\mathbf{u}))\mid f(\mathbf{x},\mathbf{u})\in X_{j-1}\}$，初值 $V_0^0:=V_f,\ X_0:=\mathcal{X}_f$。

\begin{proposition}[值函数单调性]
\label{prop:ctrl-mono}
在 \cref{def:ctrl-mpc-asm} 下，$V_j^0(\mathbf{x})\le V_{j-1}^0(\mathbf{x}),\ \forall\mathbf{x}\in X_{j-1}$；且 $V_j^0(\mathbf{x})\le V_f(\mathbf{x}),\ \forall\mathbf{x}\in\mathcal{X}_f$\cite{rawlings2020mpc}。
\end{proposition}
\begin{proof}（归纳，据 \cite{rawlings2020mpc}）
\emph{基例}：由 DP，$V_1^0(\mathbf{x})=\min_{\mathbf{u}}\{\ell(\mathbf{x},\mathbf{u})+V_f(f(\mathbf{x},\mathbf{u}))\mid f(\mathbf{x},\mathbf{u})\in\mathcal{X}_f\}$。由 (A3)，存在可行 $\mathbf{u}$ 使 $\ell+V_f(f)\le V_f(\mathbf{x})$，故 $V_1^0(\mathbf{x})\le V_f(\mathbf{x})=V_0^0(\mathbf{x})$，$\forall\mathbf{x}\in X_1=\mathcal{X}_f$。
\emph{归纳步}：设 $V_{j-1}^0(\mathbf{x})\le V_{j-2}^0(\mathbf{x})$ 于 $X_{j-2}$。对 $\mathbf{x}\in X_{j-1}$，设 $\kappa_{j-1}(\mathbf{x})$ 为时域 $j-1$ 问题的最优控制，它对时域 $j$ 问题可行（给上界）：
\[
  V_j^0(\mathbf{x})\le\ell(\mathbf{x},\kappa_{j-1}(\mathbf{x}))+V_{j-1}^0(f(\mathbf{x},\kappa_{j-1}(\mathbf{x}))).
\]
而 $V_{j-1}^0(\mathbf{x})=\ell(\mathbf{x},\kappa_{j-1}(\mathbf{x}))+V_{j-2}^0(f(\mathbf{x},\kappa_{j-1}(\mathbf{x})))$。因 $\mathbf{x}\in X_{j-1}\Rightarrow f(\mathbf{x},\kappa_{j-1})\in X_{j-2}$，由归纳假设 $V_{j-1}^0(f)\le V_{j-2}^0(f)$。两式相减得 $V_j^0(\mathbf{x})\le V_{j-1}^0(\mathbf{x})$。$\{X_j\}$ 嵌套，逐次得 $V_j^0\le V_f$ 于 $\mathcal{X}_f$。
\end{proof}

\begin{theorem}[MPC 闭环渐近稳定性]
\label{thm:ctrl-mpc-main}
在 \cref{def:ctrl-mpc-asm} 下，闭环 $\mathbf{x}^+=f(\mathbf{x},\kappa_N(\mathbf{x}))$ 的原点在 $X_N$ 中\textbf{渐近稳定}，吸引域为 $X_N$\cite{rawlings2020mpc,borrelli2017predictive}。
\end{theorem}
\begin{proof}（据 \cite{rawlings2020mpc,borrelli2017predictive}）
证 $V_N^0$ 是闭环 Lyapunov 函数（\cref{thm:ctrl-lyap-direct} 离散版）。
\emph{正不变与持久可行}：(A2) + \cref{lem:ctrl-pf2} $\Rightarrow X_N$ 对闭环正不变。
\emph{下降性}：由 DP，$V_N^0(\mathbf{x})=\ell(\mathbf{x},\kappa_N(\mathbf{x}))+V_{N-1}^0(f(\mathbf{x},\kappa_N(\mathbf{x})))$。由单调性（\cref{prop:ctrl-mono}）$V_N^0(f)\le V_{N-1}^0(f)$，故
\begin{equation}
  V_N^0(f(\mathbf{x},\kappa_N(\mathbf{x})))-V_N^0(\mathbf{x})\le-\ell(\mathbf{x},\kappa_N(\mathbf{x}))\le-\alpha_1(\|\mathbf{x}\|),\quad\forall\mathbf{x}\in X_N,
  \label{eq:ctrl-mpc-descent}
\end{equation}
末步用 (A0)（$\ell$ 正定，$\ell\ge\alpha_1(\|\mathbf{x}\|)$ 某 $\mathcal{K}_\infty$ 函数 $\alpha_1$）。
\emph{下界}：$V_N^0(\mathbf{x})\ge\ell(\mathbf{x},\kappa_N(\mathbf{x}))\ge\alpha_1(\|\mathbf{x}\|)$。
\emph{上界（原点连续）}：由 \cref{prop:ctrl-mono}，$V_N^0(\mathbf{x})\le V_f(\mathbf{x})\le\alpha_f(\|\mathbf{x}\|)$ 于 $\mathcal{X}_f$（$\alpha_f\in\mathcal{K}_\infty$，由 (A0) $V_f$ 连续正定）；因 $V_N$ 在紧集上连续，此界扩展到 $X_N$ 上的某 $\alpha_2\in\mathcal{K}_\infty$：$V_N^0(\mathbf{x})\le\alpha_2(\|\mathbf{x}\|)$。
三条 $\alpha_1(\|\mathbf{x}\|)\le V_N^0(\mathbf{x})\le\alpha_2(\|\mathbf{x}\|)$ 与下降 \cref{eq:ctrl-mpc-descent} 即 Lyapunov 函数三条件，故原点在 $X_N$ 中渐近稳定。
\end{proof}

\begin{insight}
证明骨架可一句话记住：\emph{“移位 + 终端控制”造一个可行候选解，它的代价比当前最优只多了被消掉的首阶段 $\ell$、少了被 (A3) 补偿掉的终端项，于是值函数严格下降 $\ge\ell$}。下降 + 上下界 = Lyapunov 函数 = 渐近稳定。若 $\ell,V_f$ 取二次型，三个 $\mathcal{K}_\infty$ 函数可取幂律 $\alpha_i(s)=c_is^2$，则结论加强到\emph{指数稳定}——这正是线性二次 MPC 的情形。
\end{insight}

\begin{pitfall}[概念误区：以为时域 $N$ 越长 MPC 就越稳定/越好]
新手常把 $N$ 当“越大越好”的旋钮。事实：(i) 没有终端代价/终端约束时，单纯增大 $N$ \emph{不保证}稳定（\cref{ex:ctrl-kalman} 的精神）；(ii) 有了恰当的 $V_f,\mathcal{X}_f$（\cref{def:ctrl-mpc-asm}），即便 $N$ 很短也能稳定；(iii) $N$ 越大计算越慢，实时性越差。正确理解：稳定性靠\emph{终端配料}（CLF 终端代价 + 控制不变终端集），不靠堆时域；$N$ 主要影响性能/吸引域大小与算力。\textbf{Remark}：(A0) 的 $\ell$ 正定可松弛——对 2-范数代价允许 $\mathbf{Q}\succeq0$ 但需 $(\mathbf{A},\mathbf{Q}^{1/2})$ 可观。
\end{pitfall}

\subsection{MPC 与 LQR 同源：最小可见证据}
\label{sec:ctrl-mpc-lqr}

\begin{example}[约束不激活时 MPC = LQR]
\label{ex:ctrl-mpc-eq-lqr}
\rebuilt{（据 \cite{rawlings2020mpc}）} 标量系统 $x^+=x+u$，$\ell=\tfrac12(x^2+u^2),\ V_f=\tfrac12x^2$，约束 $u\in[-1,1]$，$N=2$。凝聚为 $\mathbf{u}=(u_0,u_1)$ 的 QP，Hessian $\mathbf{H}=\big[\begin{smallmatrix}3&1\\1&2\end{smallmatrix}\big]$。对每个 $x$ 解之得隐式 MPC 律
\[
  \kappa_N(x)=\operatorname{sat}_{[-1,1]}(3x/5).
\]
当 $|x|$ 小（约束不激活）时 $\kappa_N(x)=\tfrac35x$ 是\emph{线性反馈}——恰是该系统的（有限时域）LQR 增益。去掉约束、$N\to\infty$，$\kappa_N$ 即无约束 LQR 律。这是“MPC 内部即 LQR、约束不激活时退化为 LQR”的最小可见证据。
\end{example}

\subsection{MPC 的数值结构：稀疏 KKT、Riccati、Condensing}
\label{sec:ctrl-mpc-numerical}

MPC 在线解的 QP 有特殊块带结构，两条主线\cite{rawlings2020mpc}：
\begin{itemize}
  \item \textbf{稀疏（simultaneous）}：决策变量含全部 $(\mathbf{x},\mathbf{u})$，动力学作等式约束。KKT 矩阵是对称\emph{带状}（带宽 $\sim2n+m$，总尺寸 $\sim N(2n+m)$），带状 $LDL^\top$ 分解代价\emph{随时域 $N$ 线性}、随 $(2n+m)$ 三次。
  \item \textbf{Riccati 递推}：对稀疏 KKT 做结构化因子分解，等价于时变 Riccati（一矩阵 + 两向量递推），后向算增益、前向跑状态（与 \cref{thm:ctrl-lqr-dp} 同构），约为带状 $LDL^\top$ 的 $1/3$ 算力。良定条件 $\mathbf{R}_i+\mathbf{B}_i^\top\mathbf{P}_{i+1}\mathbf{B}_i\succ0$。
  \item \textbf{Condensing（凝聚）}：用 \cref{eq:ctrl-condense-pred} 消去状态轨迹，决策变量只剩 $\mathbf{U}$（尺寸 $Nm$），Hessian 稠密。\emph{短 $N$、$n\gg m$} 时更优；长时域则 Riccati（稀疏，随 $N$ 线性）更优。二者可组合（partial condensing）。
\end{itemize}
常用求解器：OSQP、qpOASES、ECOS、acados/HPIPM（嵌入式实时）。

\begin{algo}[在线 MPC（每个采样周期）]
\label{alg:ctrl-mpc}
输入当前状态 $\mathbf{x}(t)$；输出施加的控制 $\mathbf{u}_0^*$。
\begin{enumerate}
  \item 用 $\mathbf{x}(t)$ 更新 QP 的线性项：$\tilde{\mathbf{g}}=\mathbf{g}(\mathbf{x}(t))$、约束右端 $\tilde{\mathbf{c}}=\mathbf{c}(\mathbf{x}(t))$（见 \cref{eq:ctrl-mpc-Hg}）；
  \item 解 QP \cref{eq:ctrl-mpc-qpform} 得 $\mathbf{U}^*$；若不可行则报警/退回安全策略；
  \item 取首元 $\mathbf{u}_0^*$ 施加到执行器；
  \item 推进到下一采样，回到 1（滚动时域）。
\end{enumerate}
\end{algo}

\begin{practice}
\begin{enumerate}
  \item （推导，在草稿纸上完成）对 \cref{ex:ctrl-mpc-eq-lqr}，把 $x_1=x+u_0,\ x_2=x_1+u_1$ 代入 $\tfrac12(x^2+x_1^2+x_2^2+u_0^2+u_1^2)$，展开关于 $(u_0,u_1)$ 的二次型，验证 $\mathbf{H}=\big[\begin{smallmatrix}3&1\\1&2\end{smallmatrix}\big]$ 并解无约束最优、叠加饱和得 $\kappa_N(x)=\operatorname{sat}(3x/5)$。
  \item （证明）补全 \cref{thm:ctrl-mpc-main} 中“移位 + 终端控制”候选解的构造：给定 $t$ 时刻最优序列 $\{u_0^*,\dots,u_{N-1}^*\}$，写出 $t+1$ 时刻的可行候选 $\{u_1^*,\dots,u_{N-1}^*,\kappa_f(x_N)\}$，并用 (A2)(A3) 论证其可行且代价满足 \cref{eq:ctrl-mpc-descent}。
  \item （开放）若取终端等式约束 $\mathcal{X}_f=\{\mathbf{0}\}$，说明它如何自动满足 (A2)(A3)（$\kappa_f=\mathbf{0}$、$V_f=0$），并讨论它相比“LQR 终端代价 + 椭球终端集”在吸引域大小上的代价。
\end{enumerate}
\end{practice}

% ════════════════════════════════════════════════════════════════
\section{面向机器人的控制}
\label{sec:ctrl-robot}

\paragraph{动机：把前面工具用到机器人。}
前面的 PID、LQR、MPC 大多以线性系统为对象，但机器人动力学是\emph{强非线性、强耦合}的。本节把控制工具落到机器人：先给动力学方程及其关键结构性质，再讲两大范式——\emph{计算力矩}（用模型抵消非线性，化为线性误差动态）与 \emph{PD+重力补偿}（用无源性 + LaSalle 证稳定，不需精确模型），然后是操作空间控制与全身控制 QP。这些是机械臂伺服、足式机器人下层控制的基础。

\subsection{机器人动力学方程与结构性质}
\label{sec:ctrl-robot-dyn}

$n$ 自由度刚体机械臂的关节空间动力学（由 Euler--Lagrange $\frac{\mathrm d}{\mathrm dt}\frac{\partial L}{\partial\dot{\mathbf{q}}}-\frac{\partial L}{\partial\mathbf{q}}=\boldsymbol{\tau}$，$L=$ 动能减势能，推出）\cite{tedrake_underactuated}：
\begin{equation}
  \boxed{\mathbf{M}(\mathbf{q})\ddot{\mathbf{q}}+\mathbf{C}(\mathbf{q},\dot{\mathbf{q}})\dot{\mathbf{q}}+\mathbf{g}(\mathbf{q})=\boldsymbol{\tau}\,(+\textstyle\sum\mathbf{J}^\top\mathbf{F}),}
  \label{eq:ctrl-manip}
\end{equation}
$\mathbf{M}(\mathbf{q})\succ0$ 惯性（质量）矩阵，$\mathbf{C}(\mathbf{q},\dot{\mathbf{q}})\dot{\mathbf{q}}$ 科氏/离心力，$\mathbf{g}(\mathbf{q})$ 重力力矩，$\boldsymbol{\tau}$ 关节力矩，$\sum\mathbf{J}^\top\mathbf{F}$ 外接触力的广义力。科氏阵由质量阵的第一类 Christoffel 符号定义：$\bar C_{i;j,k}=\tfrac12(\partial M_{ij}/\partial q_k+\partial M_{ik}/\partial q_j-\partial M_{jk}/\partial q_i)$。

\begin{proposition}[两条承重结构性质]
\label{prop:ctrl-robot-props}
\cite{tedrake_underactuated}
\begin{enumerate}
  \item \textbf{$\mathbf{M}$ 对称正定}且一致有界（特征值有正上下界），故处处可逆；
  \item \textbf{斜对称（无源性）}：对适当选取的 $\mathbf{C}$，矩阵 $\dot{\mathbf{M}}(\mathbf{q})-2\mathbf{C}(\mathbf{q},\dot{\mathbf{q}})$ 斜对称，即 $\mathbf{z}^\top(\dot{\mathbf{M}}-2\mathbf{C})\mathbf{z}=0,\ \forall\mathbf{z}$。
\end{enumerate}
\end{proposition}
\begin{proof}\rebuilt{（性质 2 的能量证明，据 \cite{tedrake_underactuated} 与 Spong 等 \cite{spong2006robot}）}
动能 $K=\tfrac12\dot{\mathbf{q}}^\top\mathbf{M}\dot{\mathbf{q}}$，求时间导 $\dot K=\dot{\mathbf{q}}^\top\mathbf{M}\ddot{\mathbf{q}}+\tfrac12\dot{\mathbf{q}}^\top\dot{\mathbf{M}}\dot{\mathbf{q}}$。由 \cref{eq:ctrl-manip}（无外力）$\mathbf{M}\ddot{\mathbf{q}}=\boldsymbol{\tau}-\mathbf{C}\dot{\mathbf{q}}-\mathbf{g}$，代入：
\[
  \dot K=\dot{\mathbf{q}}^\top\boldsymbol{\tau}-\dot{\mathbf{q}}^\top\mathbf{g}+\tfrac12\dot{\mathbf{q}}^\top(\dot{\mathbf{M}}-2\mathbf{C})\dot{\mathbf{q}}.
\]
物理上“内力（惯性+科氏/离心）既不产生也不耗散能量”要求末项为零，对任意 $\dot{\mathbf{q}}$ 成立 $\Rightarrow\dot{\mathbf{M}}-2\mathbf{C}$ 斜对称。
\end{proof}

\begin{insight}
斜对称性质 2 是机器人控制稳定性证明的“瑞士军刀”：在 Lyapunov 求导中，科氏项 $\dot{\mathbf{q}}^\top\mathbf{C}\dot{\mathbf{q}}$ 总会与 $\tfrac12\dot{\mathbf{q}}^\top\dot{\mathbf{M}}\dot{\mathbf{q}}$ 配对消掉（\cref{sec:ctrl-robot-pd} 会用到），于是\emph{不需要知道 $\mathbf{C}$ 的具体形式}就能证稳定。这是“无源性控制”威力的根源。性质 3（动力学对惯性参数线性可参数化 $\mathbf{M}\ddot{\mathbf{q}}+\mathbf{C}\dot{\mathbf{q}}+\mathbf{g}=\mathbf{Y}(\mathbf{q},\dot{\mathbf{q}},\ddot{\mathbf{q}})\boldsymbol{\Theta}$）是自适应控制的基础。
\end{insight}

\subsection{计算力矩控制（反馈线性化）}
\label{sec:ctrl-robot-ctc}

\paragraph{控制律与误差动态。}
\emph{逆动力学}给定期望 $(\mathbf{q}_d,\dot{\mathbf{q}}_d,\ddot{\mathbf{q}}_d)$ 算所需力矩。\textbf{计算力矩控制}\cite{cttc_wiki}（误差 $\mathbf{e}=\mathbf{q}_d-\mathbf{q}$，$\hat{(\cdot)}$ 为模型估计）：
\begin{equation}
  \boxed{\boldsymbol{\tau}=\hat{\mathbf{M}}(\mathbf{q})\big(\ddot{\mathbf{q}}_d+\mathbf{K}_d\dot{\mathbf{e}}+\mathbf{K}_p\mathbf{e}+\mathbf{K}_i\!\textstyle\int_0^t\!\mathbf{e}\,\mathrm dt'\big)+\hat{\mathbf{C}}(\mathbf{q},\dot{\mathbf{q}})\dot{\mathbf{q}}+\hat{\mathbf{g}}(\mathbf{q}).}
  \label{eq:ctrl-ctc}
\end{equation}

\begin{theorem}[计算力矩的线性解耦误差动态]
\label{thm:ctrl-ctc}
若模型精确（$\hat{\mathbf{M}}=\mathbf{M},\hat{\mathbf{C}}=\mathbf{C},\hat{\mathbf{g}}=\mathbf{g}$），\cref{eq:ctrl-ctc}（取 $\mathbf{K}_i=\mathbf{0}$）使闭环误差服从线性解耦二阶系统
\begin{equation}
  \boxed{\ddot{\mathbf{e}}+\mathbf{K}_d\dot{\mathbf{e}}+\mathbf{K}_p\mathbf{e}=\mathbf{0}.}
  \label{eq:ctrl-ctc-err}
\end{equation}
\end{theorem}
\begin{proof}\rebuilt{（不跳步，据 \cite{cttc_wiki}）}
把 \cref{eq:ctrl-ctc}（$\mathbf{K}_i=\mathbf{0}$、模型精确）代入动力学 $\mathbf{M}\ddot{\mathbf{q}}+\mathbf{C}\dot{\mathbf{q}}+\mathbf{g}=\boldsymbol{\tau}$：
\[
  \mathbf{M}\ddot{\mathbf{q}}+\mathbf{C}\dot{\mathbf{q}}+\mathbf{g}=\mathbf{M}\big(\ddot{\mathbf{q}}_d+\mathbf{K}_d\dot{\mathbf{e}}+\mathbf{K}_p\mathbf{e}\big)+\mathbf{C}\dot{\mathbf{q}}+\mathbf{g}.
\]
两边消去 $\mathbf{C}\dot{\mathbf{q}}+\mathbf{g}$，左乘 $\mathbf{M}^{-1}$（$\mathbf{M}\succ0$ 可逆）：$\ddot{\mathbf{q}}=\ddot{\mathbf{q}}_d+\mathbf{K}_d\dot{\mathbf{e}}+\mathbf{K}_p\mathbf{e}$。用 $\mathbf{e}=\mathbf{q}_d-\mathbf{q}\Rightarrow\ddot{\mathbf{e}}=\ddot{\mathbf{q}}_d-\ddot{\mathbf{q}}$，移项即得 $\ddot{\mathbf{e}}+\mathbf{K}_d\dot{\mathbf{e}}+\mathbf{K}_p\mathbf{e}=\mathbf{0}$。
\end{proof}

\begin{insight}
计算力矩是机器人版的\textbf{精确反馈线性化}：前馈逆动力学项 $\hat{\mathbf{C}}\dot{\mathbf{q}}+\hat{\mathbf{g}}$ 与惯性整形 $\hat{\mathbf{M}}(\cdot)$ \emph{抵消}了非线性，剩下解耦的线性误差动态——结构是“\emph{内环}（逆动力学抵消）+ \emph{外环}（线性 PD/PID）”。$\mathbf{K}_p,\mathbf{K}_d$ 对角正定时，\cref{eq:ctrl-ctc-err} 分解为 $n$ 个独立二阶环 $\ddot e_j+k_{d,j}\dot e_j+k_{p,j}e_j=0$，取临界阻尼 $k_{d,j}=2\sqrt{k_{p,j}}$ 即无超调指数收敛。\textbf{衔接 LQR/MPC}：反馈线性化把机器人化为“双积分器”$\ddot{\mathbf{e}}=\mathbf{v}$（$\mathbf{v}$ 为外环虚拟输入），其上可叠 LQR（\cref{sec:ctrl-lqr}）或加约束的 MPC（\cref{sec:ctrl-mpc}）——这是把前几节用于机器人的标准桥。
\end{insight}

\paragraph{反馈线性化的一般理论（简述）。}
\rebuilt{SISO 输入-输出反馈线性化（据 \cite{khalil2002nonlinear} / Slotine）}：系统 $\dot x=f(x)+g(x)u,\ y=h(x)$，用 Lie 导数 $L_fh=\frac{\partial h}{\partial x}f$。\emph{相对阶} $r$ = 输出微分多少次输入才显式出现；当 $r=n$，坐标变换 $z=[h,L_fh,\dots,L_f^{n-1}h]^\top$ 把系统化为积分器链，控制律 $u=\frac{1}{L_gL_f^{n-1}h}(-L_f^nh+v)$ 得线性 $\dot z=\mathbf{A}z+\mathbf{b}v$。机器人操纵臂方程因 $\mathbf{M}$ 可逆恰好满足全状态线性化条件（解耦矩阵即 $\mathbf{M}$），计算力矩正是其特例：取虚拟输入 $\mathbf{v}=\ddot{\mathbf{q}}_d+\mathbf{K}_d\dot{\mathbf{e}}+\mathbf{K}_p\mathbf{e}$、$\mathbf{u}=\boldsymbol{\tau}$。

\subsection{PD + 重力补偿：无源性证明}
\label{sec:ctrl-robot-pd}

\paragraph{反面：计算力矩要精确模型。}
计算力矩需\emph{完整且精确}的 $\mathbf{M},\mathbf{C},\mathbf{g}$，模型误差会破坏线性化。对\emph{点到点定位}（$\mathbf{q}_d$ 常值）任务，有一个惊人简单且鲁棒的替代：PD + 重力补偿，只需重力项 $\mathbf{g}$。

\paragraph{控制律。}
误差 $\tilde{\mathbf{q}}=\mathbf{q}_d-\mathbf{q}$（$\mathbf{q}_d$ 常值），$\mathbf{K}_p,\mathbf{K}_d\succ0$\cite{takegaki1981feedback}：
\begin{equation}
  \boldsymbol{\tau}=\mathbf{g}(\mathbf{q})+\mathbf{K}_p\tilde{\mathbf{q}}-\mathbf{K}_d\dot{\mathbf{q}}.
  \label{eq:ctrl-pd-grav}
\end{equation}
Takegaki 与 Arimoto（1981）证明：尽管机器人动力学高度非线性，\emph{简单 PD + 重力补偿即可全局解决全驱动机械臂的点到点定位}\cite{takegaki1981feedback}。

\begin{theorem}[PD + 重力补偿的渐近稳定性]
\label{thm:ctrl-pd-grav}
\cref{eq:ctrl-pd-grav} 使闭环平衡点 $(\tilde{\mathbf{q}},\dot{\mathbf{q}})=(\mathbf{0},\mathbf{0})$ 渐近稳定（$V$ 径向无界时全局）。
\end{theorem}
\begin{proof}\rebuilt{（据 Spong 等 \cite{spong2006robot} 标准证）}
把 \cref{eq:ctrl-pd-grav} 代入 \cref{eq:ctrl-manip}，重力项消去，闭环
\[
  \mathbf{M}(\mathbf{q})\ddot{\mathbf{q}}+\mathbf{C}(\mathbf{q},\dot{\mathbf{q}})\dot{\mathbf{q}}=\mathbf{K}_p\tilde{\mathbf{q}}-\mathbf{K}_d\dot{\mathbf{q}}.
\]
取 Lyapunov 候选（动能 + 势能整形）
\[
  V=\tfrac12\dot{\mathbf{q}}^\top\mathbf{M}(\mathbf{q})\dot{\mathbf{q}}+\tfrac12\tilde{\mathbf{q}}^\top\mathbf{K}_p\tilde{\mathbf{q}}.
\]
$\mathbf{M}\succ0,\mathbf{K}_p\succ0\Rightarrow V$ 正定。求导（用 $\dot{\tilde{\mathbf{q}}}=-\dot{\mathbf{q}}$，$\mathbf{q}_d$ 常值）：
\[
  \dot V=\dot{\mathbf{q}}^\top\mathbf{M}\ddot{\mathbf{q}}+\tfrac12\dot{\mathbf{q}}^\top\dot{\mathbf{M}}\dot{\mathbf{q}}-\tilde{\mathbf{q}}^\top\mathbf{K}_p\dot{\mathbf{q}}.
\]
代入 $\mathbf{M}\ddot{\mathbf{q}}=\mathbf{K}_p\tilde{\mathbf{q}}-\mathbf{K}_d\dot{\mathbf{q}}-\mathbf{C}\dot{\mathbf{q}}$：
\[
  \dot V=\dot{\mathbf{q}}^\top(\mathbf{K}_p\tilde{\mathbf{q}}-\mathbf{K}_d\dot{\mathbf{q}}-\mathbf{C}\dot{\mathbf{q}})+\tfrac12\dot{\mathbf{q}}^\top\dot{\mathbf{M}}\dot{\mathbf{q}}-\tilde{\mathbf{q}}^\top\mathbf{K}_p\dot{\mathbf{q}}.
\]
$\dot{\mathbf{q}}^\top\mathbf{K}_p\tilde{\mathbf{q}}$ 与 $-\tilde{\mathbf{q}}^\top\mathbf{K}_p\dot{\mathbf{q}}$ 互消（标量 + $\mathbf{K}_p$ 对称），整理：
\[
  \dot V=-\dot{\mathbf{q}}^\top\mathbf{K}_d\dot{\mathbf{q}}+\underbrace{\tfrac12\dot{\mathbf{q}}^\top(\dot{\mathbf{M}}-2\mathbf{C})\dot{\mathbf{q}}}_{=0\;(\text{性质 2 斜对称})}=-\dot{\mathbf{q}}^\top\mathbf{K}_d\dot{\mathbf{q}}\le0,
\]
\emph{半负定}（$\dot V=0\iff\dot{\mathbf{q}}=\mathbf{0}$）。由 LaSalle（\cref{thm:ctrl-lasalle}）：令 $\dot V=0\Rightarrow\dot{\mathbf{q}}=\mathbf{0}\Rightarrow\ddot{\mathbf{q}}=\mathbf{0}$，代入闭环得 $\mathbf{K}_p\tilde{\mathbf{q}}=\mathbf{0}\Rightarrow\tilde{\mathbf{q}}=\mathbf{0}$（$\mathbf{K}_p\succ0$）。故 $\{\dot V=0\}$ 中最大不变集仅 $\{(\tilde{\mathbf{q}},\dot{\mathbf{q}})=(\mathbf{0},\mathbf{0})\}$，渐近稳定。
\end{proof}

\begin{insight}
\cref{thm:ctrl-pd-grav} 是 LaSalle（\cref{sec:ctrl-lasalle}）与斜对称性质（\cref{prop:ctrl-robot-props}）的完美合奏：$\dot V$ 只半负定（阻尼仅在速度非零时耗能），直接法不够，LaSalle 补足；科氏项靠斜对称恰好消失，故\emph{无需精确知道 $\mathbf{M},\mathbf{C}$}就能证稳定——这是无源性控制对模型误差鲁棒的根源。这正是它与计算力矩的分工：计算力矩\emph{需精确模型}但给\emph{轨迹跟踪}的线性误差动态；PD+重力补偿\emph{只需重力}、对 $\mathbf{M},\mathbf{C}$ 误差鲁棒，但只解\emph{点到点定位}。
\end{insight}

\begin{table}[H]\centering\small
\caption{计算力矩 vs PD+重力补偿}
\label{tab:ctrl-ctc-vs-pd}
\begin{tabular}{p{0.16\textwidth} p{0.36\textwidth} p{0.36\textwidth}}
\toprule
维度 & 计算力矩（\cref{sec:ctrl-robot-ctc}） & PD+重力补偿（\cref{sec:ctrl-robot-pd}） \\
\midrule
任务 & 轨迹跟踪（$\mathbf{q}_d(t)$ 时变） & 点到点定位（$\mathbf{q}_d$ 常值） \\
模型需求 & 完整且精确 $\mathbf{M},\mathbf{C},\mathbf{g}$ & 仅重力 $\mathbf{g}$（+ $\mathbf{M}\succ0$、斜对称） \\
闭环 & 精确线性解耦 $\ddot{\mathbf{e}}+\mathbf{K}_d\dot{\mathbf{e}}+\mathbf{K}_p\mathbf{e}=\mathbf{0}$ & 非线性，Lyapunov+LaSalle 证渐近稳定 \\
计算量 & 大（每周期算逆动力学，RNEA $O(n)$） & 小（仅重力补偿） \\
鲁棒性 & 对模型误差敏感 & 对 $\mathbf{M},\mathbf{C}$ 误差鲁棒（无源性） \\
\bottomrule
\end{tabular}
\end{table}

\subsection{操作空间控制}
\label{sec:ctrl-robot-osc}

\rebuilt{Khatib 操作空间公式（据 \cite{khatib1987operational}）}：末端任务坐标 $\mathbf{x}$、任务雅可比 $\mathbf{J}$（$\dot{\mathbf{x}}=\mathbf{J}\dot{\mathbf{q}}$）。把关节动力学投影到任务空间：
\begin{equation}
  \boldsymbol{\Lambda}(\mathbf{q})\ddot{\mathbf{x}}+\boldsymbol{\mu}(\mathbf{q},\dot{\mathbf{q}})+\mathbf{p}(\mathbf{q})=\mathbf{F},\qquad
  \boxed{\boldsymbol{\Lambda}=(\mathbf{J}\mathbf{M}^{-1}\mathbf{J}^\top)^{-1}}
  \label{eq:ctrl-osc}
\end{equation}
$\boldsymbol{\Lambda}$ \textbf{操作空间惯性矩阵}，$\boldsymbol{\mu}$ 任务空间科氏/离心，$\mathbf{p}$ 任务空间重力，$\mathbf{F}$ 末端广义力。关节力矩由对偶映射 $\boldsymbol{\tau}=\mathbf{J}^\top\mathbf{F}$ 给出。任务控制力 $\mathbf{F}=\hat{\boldsymbol{\Lambda}}\mathbf{F}^*+\hat{\boldsymbol{\mu}}+\hat{\mathbf{p}}$（$\mathbf{F}^*$ 为力级 PD 反馈律 $\mathbf{F}^*=\ddot{\mathbf{x}}_d+\mathbf{K}_d(\dot{\mathbf{x}}_d-\dot{\mathbf{x}})+\mathbf{K}_p(\mathbf{x}_d-\mathbf{x})$）则给出解耦行为 $\ddot{\mathbf{x}}=\mathbf{F}^*$。

\paragraph{任务空间姿态误差（衔接本书李群约定）。}
当任务含\emph{姿态}时，姿态误差不是简单相减，而应用本书 \cref{ch:lie} 的对数映射：以右扰动表达姿态误差 $\mathbf{e}_R=\mathrm{Log}(\mathbf{R}^\top\mathbf{R}_d)\in\mathbb{R}^3$（机体系下的旋转偏差），其上做 PD：$\boldsymbol{\omega}^*=\mathbf{K}_R\mathbf{e}_R+\mathbf{K}_\omega(\boldsymbol{\omega}_d-\boldsymbol{\omega})$。这样操作空间姿态控制与本书 $\SO(3)$ 主线相容。

\subsection{力控与阻抗控制}
\label{sec:ctrl-robot-impedance}

\paragraph{动机：与环境接触时纯位置控制会“打架”。}
机器人插轴、打磨、开门、与人协作时要\emph{接触}环境。此时若仍用刚性位置控制（计算力矩/操作空间 PD），微小的位置误差会被高刚度放大成巨大接触力，导致卡死或损坏。Hogan（1985）的洞见\cite{hogan1985impedance}：\textbf{接触任务不应单独控制位置或力，而应控制末端的\emph{动态关系}——即机械阻抗}（力对运动的响应）。

\begin{definition}[机械阻抗与导纳]
\label{def:ctrl-impedance}
末端\textbf{机械阻抗}是“运动 $\to$ 力”的算子（给定运动，输出反作用力）；\textbf{导纳}是其逆（给定力，输出运动）\cite{hogan1985impedance}。Hogan 论证：与含惯性的环境交互时，\emph{两个互连子系统不能都是导纳}（否则接触瞬间力/运动无法相容），故\textbf{机械手应表现为阻抗、环境为导纳}。
\end{definition}

\paragraph{目标阻抗与控制律。}
按 \cref{def:ctrl-impedance}，机械手应呈现一个阻抗。设期望末端表现为一个二阶\emph{质量--弹簧--阻尼}（目标阻抗），误差 $\mathbf{e}_x=\mathbf{x}_d-\mathbf{x}$，外力 $\mathcal{F}_{\rm ext}$\cite{hogan1985impedance}：
\begin{equation}
  \boxed{\mathbf{M}_d\ddot{\mathbf{e}}_x+\mathbf{D}_d\dot{\mathbf{e}}_x+\mathbf{K}_d^x\mathbf{e}_x=-\mathcal{F}_{\rm ext},}
  \label{eq:ctrl-impedance}
\end{equation}
$\mathbf{M}_d,\mathbf{D}_d,\mathbf{K}_d^x\succ0$ 为期望惯性/阻尼/刚度。要让真实机器人\emph{呈现}该行为，把操作空间动力学 \cref{eq:ctrl-osc} 的任务力取为
\begin{equation}
  \mathbf{F}=\hat{\boldsymbol{\Lambda}}\mathbf{M}_d^{-1}\big(\mathbf{M}_d\ddot{\mathbf{x}}_d+\mathbf{D}_d\dot{\mathbf{e}}_x+\mathbf{K}_d^x\mathbf{e}_x\big)+\hat{\boldsymbol{\mu}}+\hat{\mathbf{p}}-\big(\hat{\boldsymbol{\Lambda}}\mathbf{M}_d^{-1}-\mathbf{I}\big)\mathcal{F}_{\rm ext},
  \label{eq:ctrl-impedance-law}
\end{equation}
$\boldsymbol{\tau}=\mathbf{J}^\top\mathbf{F}$。特例 $\mathbf{M}_d=\boldsymbol{\Lambda}$（不整形惯性）时\emph{无需力传感}，律退化为 $\mathbf{F}=\hat{\boldsymbol{\Lambda}}\ddot{\mathbf{x}}_d+\mathbf{D}_d\dot{\mathbf{e}}_x+\mathbf{K}_d^x\mathbf{e}_x+\hat{\boldsymbol{\mu}}+\hat{\mathbf{p}}$，即“带前馈的任务空间 PD + 自然惯性”——这正是 \cref{eq:ctrl-osc} 末解耦律在“弹簧--阻尼”读法下的本来面目。

\begin{insight}
\textbf{刚度 $\mathbf{K}_d^x$ 是“柔顺旋钮”}：调小 $\mathbf{K}_d^x$ 让末端“软”（顺从外力、安全交互），调大让它“硬”（精确定位、抗扰）——同一控制器经一个参数从位置控制连续过渡到柔顺控制。\textbf{阻抗 vs 导纳实现}：\emph{阻抗控制}测运动、输出力矩（无需力传感器，\cref{eq:ctrl-impedance-law} 的免力传感特例），适合高带宽；\emph{导纳控制}测力、输出位置参考给内层位置环，适合刚性机器人 + 力传感器但带宽受限。\textbf{与无源性的联系}：取目标阻抗为无源（$\mathbf{D}_d\succ0$ 耗能）即保证交互稳定，这与 \cref{sec:ctrl-robot-pd} 的无源性证明同根——能量观点贯穿机器人交互控制。
\end{insight}

\begin{pitfall}[概念误区：用高增益位置控制器去“顶”接触力]
新手常想“接触力太大就把位置控制器调软一点”，却仍用纯位置环。问题：纯位置控制的等效刚度由反馈增益隐式决定、且常很高，接触时位置误差 $\to$ 巨大力；调低增益又牺牲自由空间的跟踪精度。现象：要么硬碰硬（力尖峰、跳变），要么全程松软（定位不准）。根本原因：位置与力在接触方向上\emph{不能同时独立指定}（环境约束把二者耦合）。正确做法：用阻抗控制 \cref{eq:ctrl-impedance} \emph{显式}设定力--运动关系（刚度 $\mathbf{K}_d^x$），或在已知接触方向上做\emph{混合力/位置控制}（受约束方向控力、自由方向控位）。
\end{pitfall}

\subsection{全身控制 WBC 与分层 QP}
\label{sec:ctrl-robot-wbc}

\paragraph{浮动基动力学。}
\rebuilt{足式/人形机器人（据 Wensing 等 \cite{wensing2024legged}）}：位形 $\mathbf{q}=(\mathbf{q}_b,\mathbf{q}_j)$（基座 6 DoF + 关节），广义速度 $\boldsymbol{\nu}$，动力学
\begin{equation}
  \mathbf{M}(\mathbf{q})\dot{\boldsymbol{\nu}}+\mathbf{C}(\mathbf{q},\boldsymbol{\nu})\boldsymbol{\nu}+\boldsymbol{\tau}_g(\mathbf{q})=\mathbf{S}^\top\boldsymbol{\tau}+\mathbf{J}_c(\mathbf{q})^\top\boldsymbol{\lambda},
  \label{eq:ctrl-floatbase}
\end{equation}
$\mathbf{S}=[\mathbf{0}\ \mathbf{I}_{n_j}]$ 选择矩阵（基座\emph{不可直接驱动}，前 6 行欠驱动），$\mathbf{J}_c$ 接触雅可比，$\boldsymbol{\lambda}$ 接触力。刚性接触约束：接触点零加速度 $\mathbf{J}_c\dot{\boldsymbol{\nu}}+\dot{\mathbf{J}}_c\boldsymbol{\nu}=\mathbf{0}$。

\paragraph{WBC 为一个 QP。}
任务 $i$ 的期望加速度由 PD 给 $\ddot{\mathbf{x}}_i^d=\ddot{\mathbf{x}}_i^{\rm ref}+\mathbf{K}_{p,i}(\mathbf{x}_i^{\rm ref}-\mathbf{x}_i)+\mathbf{K}_{d,i}(\dot{\mathbf{x}}_i^{\rm ref}-\dot{\mathbf{x}}_i)$。关键观察：\emph{动力学 \cref{eq:ctrl-floatbase} 对决策变量 $(\dot{\boldsymbol{\nu}},\boldsymbol{\tau},\boldsymbol{\lambda})$ 都是线性的}，故在加速度级写正则化时整个问题保持线性约束，可写成 QP\cite{wensing2024legged}：
\begin{align}
  \min_{\dot{\boldsymbol{\nu}},\boldsymbol{\tau},\boldsymbol{\lambda}}\ &\sum_i w_i\big\|\mathbf{J}_i\dot{\boldsymbol{\nu}}+\dot{\mathbf{J}}_i\boldsymbol{\nu}-\ddot{\mathbf{x}}_i^d\big\|^2+w_\tau\|\boldsymbol{\tau}\|^2\notag\\
  \text{s.t.}\ &\mathbf{M}\dot{\boldsymbol{\nu}}+\mathbf{C}\boldsymbol{\nu}+\boldsymbol{\tau}_g=\mathbf{S}^\top\boldsymbol{\tau}+\mathbf{J}_c^\top\boldsymbol{\lambda}\quad(\text{全身动力学}),\notag\\
  &\mathbf{J}_c\dot{\boldsymbol{\nu}}=-\dot{\mathbf{J}}_c\boldsymbol{\nu}\quad(\text{接触不滑动}),\quad\boldsymbol{\lambda}\in\mathcal{FC}\ (\text{摩擦锥}),\quad\boldsymbol{\tau}_{\min}\le\boldsymbol{\tau}\le\boldsymbol{\tau}_{\max}.
  \label{eq:ctrl-wbc-qp}
\end{align}
即在全身动力学 + 接触 + 摩擦锥 + 力矩限约束下，同时（加权）满足质心、足端、末端执行器等多任务的单一 QP。摩擦锥常线性化为多边形（金字塔）使之成 QP。优点：依赖成熟 QP 解器（qpOASES/OSQP）；局限：软权重难整定、低优先级无法处理不等式。

\paragraph{分层 / 严格优先级 QP（HQP）。}
任务冲突（QP 无解）时，两法处理\cite{wensing2024legged}：\emph{软优先级}（加权各违背度）或\emph{严格优先级}（定义层级，高优先级永不为低优先级牺牲）。后者称\textbf{词典最小二乘 / 分层 QP（HQP）}：优先级 $1,\dots,p$，第 $k$ 级在不破坏更高级最优残差的零空间内最小化本级残差：
\begin{equation}
  \min_{\mathbf{x}}\ \|\mathbf{A}_k\mathbf{x}-\mathbf{b}_k\|^2\quad\text{s.t.}\ \mathbf{A}_j\mathbf{x}-\mathbf{b}_j=\mathbf{w}_j^*\ (j<k),
  \label{eq:ctrl-hqp}
\end{equation}
等价于在更高级雅可比的零空间内逐级求解（零空间投影 $\mathbf{x}_k=\mathbf{x}_{k-1}+\mathbf{N}_{k-1}\mathbf{z}_k$）。HQP 是软优先级 QP 的严格超集（权重比趋于无穷的极限）。

\begin{pitfall}[概念误区：以为 WBC 是“一个大 PID”]
WBC QP \cref{eq:ctrl-wbc-qp} 表面像“给每个任务配 PD”，但本质不同：它在\emph{一个优化问题}里同时协调多个相互冲突的任务，并\emph{硬性满足}动力学一致性、接触、摩擦锥、力矩限等约束——这是单独的 PD 环做不到的。现象：若把各任务 PD 独立施加再相加，会破坏动力学一致性、违反摩擦锥导致打滑或腾空。根本原因：浮动基机器人欠驱动 + 接触约束耦合，必须联合求解。正确理解：WBC 是“瞬时（1 步）的、带约束的逆动力学 QP”；与之配套的上层 MPC 是“多步预测 QP”（见 \cref{sec:ctrl-legged}）。
\end{pitfall}

\begin{practice}
\begin{enumerate}
  \item （推导）从 \cref{eq:ctrl-osc} 验证 $\boldsymbol{\tau}=\mathbf{J}^\top\mathbf{F}$ 的对偶性：由虚功原理 $\boldsymbol{\tau}^\top\delta\mathbf{q}=\mathbf{F}^\top\delta\mathbf{x}$ 与 $\delta\mathbf{x}=\mathbf{J}\delta\mathbf{q}$ 推出。
  \item （证明）补全 \cref{thm:ctrl-ctc} 中含积分项 $\mathbf{K}_i$ 的情形：证误差满足 $\dddot{\mathbf{e}}+\mathbf{K}_d\ddot{\mathbf{e}}+\mathbf{K}_p\dot{\mathbf{e}}+\mathbf{K}_i\mathbf{e}=\mathbf{0}$（对 \cref{eq:ctrl-ctc-err} 求导并代入），讨论 $\mathbf{K}_i$ 的稳定性上界（用 \cref{sec:ctrl-pid-routh} 的 Routh 三阶条件）。
  \item （跨章综合）用 \cref{ch:lie} 的右扰动与对数映射，写出一个机械臂末端\emph{位姿}（$\SE(3)$）跟踪的操作空间 PD 律，说明姿态误差为何不能用旋转矩阵相减。
\end{enumerate}
\end{practice}

% ════════════════════════════════════════════════════════════════
\section{足式机器人控制概览}
\label{sec:ctrl-legged}

\paragraph{动机：把全章工具汇于一个旗舰场景。}
足式机器人（四足、人形）是控制理论的试金石：与环境\emph{间歇单向接触}（非光滑、刚性动力学）、\emph{高自由度}、强实时性。本节以分层架构组织：上层用\emph{简化模型 MPC}（前几节的 LQR/MPC + 凸 QP）规划质心轨迹与接触力，下层用\emph{全身控制 WBC}（\cref{sec:ctrl-robot-wbc}）分解为关节力矩。以 MIT Cheetah 3 凸 MPC 为旗舰例，完整推导其单刚体模型如何线性化、离散化为凸 QP。本节为\emph{前沿/选读}。

\subsection{分层架构与简化模型谱系}
\label{sec:ctrl-legged-arch}

足式最优控制问题（OCP）的一般在线求解仍不可行，故社区把大问题拆成小块\cite{wensing2024legged}（自上而下）：

\begin{figure}[ht]
\centering
\begin{tikzpicture}[>=Stealth, font=\small, node distance=3mm,
  lay/.style={draw, rounded corners, align=center, minimum width=58mm, minimum height=10mm},
  ann/.style={font=\scriptsize, gray!50!black, align=left}]
  \node[lay, fill=second!8, draw=second!60!black] (cp) {接触规划器（运动学模型）\\长时域、低速率 $\to$ 接触序列};
  \node[lay, fill=third!8, draw=third!60!black, below=of cp] (mpc) {简化模型 MPC（质心/单刚体 SRBD）\\中时域 $\sim$1\,s、中速率 $\to$ 质心轨迹 + 接触力};
  \node[lay, fill=main!8, draw=main!60!black, below=of mpc] (wbc) {全身控制 WBC（完整模型）\\短时域、高速率 $\sim$1\,kHz $\to$ 关节力矩};
  \node[lay, fill=gray!8, draw=gray!60!black, below=of wbc] (robot) {机器人 + 执行器};
  \draw[->, thick] (cp)--(mpc); \draw[->, thick] (mpc)--(wbc); \draw[->, thick] (wbc)--(robot);
  \node[ann, right=4mm of cp] {本节 \cref{sec:ctrl-legged-cheetah}};
  \node[ann, right=4mm of mpc] {\cref{sec:ctrl-mpc} + \cref{sec:ctrl-legged-cheetah}};
  \node[ann, right=4mm of wbc] {\cref{sec:ctrl-robot-wbc}};
\end{tikzpicture}
\caption{足式机器人控制的典型分层架构：自上而下模型越来越完整、时域越来越短、速率越来越高。（据 Wensing 等 \cite{wensing2024legged} 重绘）}
\label{fig:ctrl-legged-arch}
\end{figure}

\textbf{简化模型谱系}（自简到繁）\cite{wensing2024legged}：点质量 / 线性倒立摆（LIP）$\to$ 单刚体（SRBD，Cheetah 用）$\to$ 质心动力学 $\to$ 完整刚体。模型越简，状态越小、越线性，越利于在线快算；但简化假设（如\emph{恒定质心高度}）或被忽略的约束（关节限位/力矩限）会限制可生成的运动。

\paragraph{LIP 与质心动力学。}
LIP 动力学 $\ddot c_{x,y}=\omega^2(c_{x,y}-p_{x,y})$（$c$ 质心水平位置，$p$ 压力中心，$\omega=\sqrt{g/h}$，$h$ 恒定质心高度）\cite{wensing2024legged}。质心动力学（Newton--Euler 关于质心）：$\dot{\mathbf{h}}_G=[\dot{\mathbf{k}}_G;\dot{\mathbf{l}}_G]=[\mathbf{0};-M\mathbf{a}_g]+\sum_i[(\mathbf{p}_i-\mathbf{p}_{\rm CoM})\times\boldsymbol{\lambda}_i;\ \boldsymbol{\lambda}_i]$（角动量 + 线动量随接触力变化）。

\subsection{接触约束与摩擦锥}
\label{sec:ctrl-legged-contact}

接触是单向的（不能“拉”地）：法向力 $\lambda_n\ge0$，$\lambda_n<0$ 时分离。Coulomb 摩擦要求切向力落在\textbf{摩擦锥} $\|\boldsymbol{\lambda}_t\|\le\mu\lambda_n$ 内。为进 QP，摩擦锥常\emph{线性化为方形金字塔}\cite{dicarlo2018cheetah}：对每个站立足
\begin{equation}
  -\mu f_z\le f_x\le\mu f_z,\quad -\mu f_z\le f_y\le\mu f_z,\quad f_{\min}\le f_z\le f_{\max}.
  \label{eq:ctrl-friction-pyramid}
\end{equation}

\begin{insight}
摩擦锥金字塔化是“为了能用 QP 而做的工程近似”：真实摩擦锥是\emph{二阶锥}（圆锥），约束 $\|\boldsymbol{\lambda}_t\|\le\mu\lambda_n$ 让问题成 SOCP（二阶锥规划）；用方形金字塔（内接/外接多边形）线性化后退化为\emph{线性不等式}，问题成 QP，求解更快、库更多。代价是金字塔比圆锥略保守或略宽松（取决于内接还是外接）。这正是 \cref{sec:ctrl-robot-wbc} 中“摩擦锥 $\in\mathcal{FC}$”写成线性约束的来历。
\end{insight}

\subsection{MIT Cheetah 凸 MPC：单刚体模型完整推导}
\label{sec:ctrl-legged-cheetah}

\rebuilt{以下逐式据 Di Carlo 等 2018\cite{dicarlo2018cheetah}。} 把机器人近似为\textbf{单刚体（SRBD）}：忽略腿质量，只建质心位姿与地面反力。

\paragraph{世界系刚体动力学。}
对每足地面反力 $\mathbf{f}_i\in\mathbb{R}^3$、质心到力作用点向量 $\mathbf{r}_i$\cite{dicarlo2018cheetah}：
\begin{align}
  \ddot{\mathbf{p}}&=\frac{\sum_{i=1}^n\mathbf{f}_i}{m}-\mathbf{g},\label{eq:ctrl-srbd-trans}\\
  \frac{\mathrm d}{\mathrm dt}(\mathbf{I}\boldsymbol{\omega})&=\sum_{i=1}^n\mathbf{r}_i\times\mathbf{f}_i,\label{eq:ctrl-srbd-rot}\\
  \dot{\mathbf{R}}&=\boldsymbol{\omega}^\wedge\mathbf{R},\label{eq:ctrl-srbd-R}
\end{align}
$\mathbf{p}$ 质心位置，$m$ 质量，$\mathbf{I}$ 世界系惯性张量，$\boldsymbol{\omega}$ 角速度，$\mathbf{R}$ 体到世界旋转（本书 \cref{ch:lie} 记号，$\boldsymbol{\omega}^\wedge$ 即叉乘算子）。

\paragraph{姿态用 Z-Y-X Euler 角线性化。}
为便利凸化，姿态用 $\boldsymbol{\Theta}=[\phi\ \theta\ \psi]^\top$（roll-pitch-yaw），$\mathbf{R}=\mathbf{R}_z(\psi)\mathbf{R}_y(\theta)\mathbf{R}_x(\phi)$\cite{dicarlo2018cheetah}。\textbf{小 roll/pitch 近似}下 Euler 角速率
\begin{equation}
  \begin{bmatrix}\dot\phi\\\dot\theta\\\dot\psi\end{bmatrix}\approx\mathbf{R}_z(\psi)^\top\boldsymbol{\omega}.
  \label{eq:ctrl-euler-rate}
\end{equation}
角动量导数近似 $\frac{\mathrm d}{\mathrm dt}(\mathbf{I}\boldsymbol{\omega})=\mathbf{I}\dot{\boldsymbol{\omega}}+\boldsymbol{\omega}\times(\mathbf{I}\boldsymbol{\omega})\approx\mathbf{I}\dot{\boldsymbol{\omega}}$（舍去进动项，对小角速度体很小），世界系惯性近似 $\hat{\mathbf{I}}=\mathbf{R}_z(\psi)\,{}_B\mathbf{I}\,\mathbf{R}_z(\psi)^\top$。

\begin{note}
\textbf{足式 MPC 用 Euler 角与本书李群主线的关系.}\;
这里用 ZYX Euler 角\emph{仅为线性化便利}（让动力学只依赖 yaw 与落足位置，从而成为线性时变系统，适合凸 QP）。它与本书以 $\mathbf{R}\in\SO(3)$/李代数为主线不冲突——Euler 角是局部参数化，小 roll/pitch 下足够好。若要全姿态范围或避奇异，应回到 \cref{ch:lie} 的指数映射/右扰动（这也是更现代的“流形上 MPC”做法）。
\end{note}

\paragraph{线性状态空间（含重力增广态）。}
合并近似的姿态与平动动力学，状态 $\mathbf{x}=[\boldsymbol{\Theta};\mathbf{p};\boldsymbol{\omega};\dot{\mathbf{p}};g]\in\mathbb{R}^{13}$（3 姿态 + 3 位置 + 3 角速度 + 3 线速度 + 1 重力增广）\cite{dicarlo2018cheetah}：
\begin{equation}
  \dot{\mathbf{x}}(t)=\mathbf{A}_c(\psi)\mathbf{x}(t)+\mathbf{B}_c(\mathbf{r}_1,\dots,\mathbf{r}_n,\psi)\mathbf{u}(t),\quad\mathbf{u}=[\mathbf{f}_1;\dots;\mathbf{f}_n],
  \label{eq:ctrl-cheetah-ss}
\end{equation}
其中 $\mathbf{B}_c$ 的力到角加速度块为 $\hat{\mathbf{I}}^{-1}\mathbf{r}_i^\wedge$、力到线加速度块为 $\mathbf{I}_3/m$。\textbf{关键}：\cref{eq:ctrl-cheetah-ss} 只依赖 yaw $\psi$ 与落足位置 $\mathbf{r}_i$——若二者可提前算出，动力学成为\emph{线性时变（LTV）}，从而 MPC 问题\emph{凸}（QP）。用增广态把重力 $\mathbf{g}$ 这一仿射项并入状态（即 \cref{sec:ctrl-lqr-dp} 的仿射增广技巧）。零阶保持离散化得 $\mathbf{x}[n+1]=\hat{\mathbf{A}}\mathbf{x}[n]+\hat{\mathbf{B}}[n]\mathbf{u}[n]$。

\paragraph{凸 MPC 的 QP。}
代入 \cref{sec:ctrl-mpc-qp} 的凝聚式 \cref{eq:ctrl-mpc-qpform,eq:ctrl-mpc-Hg}，加力约束 \cref{eq:ctrl-friction-pyramid} 与摆动腿力置零的等式约束 $\mathbf{D}_i\mathbf{u}_i=\mathbf{0}$，得凸 QP\cite{dicarlo2018cheetah}。Hessian 维度 $3nk\times3nk$（$n$ 足数、$k$ 时域），\emph{与状态数无关}（凝聚收益）。Cheetah 用 ECOS/qpOASES 在 $<1$\,ms 解出（20--30\,Hz），时域 0.5\,s（一个步态周期，10--16 步）。

\begin{table}[H]\centering\small
\caption{MIT Cheetah 3 凸 MPC 关键参数（据 \cite{dicarlo2018cheetah}）}
\label{tab:ctrl-cheetah}
\begin{tabular}{llll}
\toprule
参数 & 值 & 参数 & 值 \\
\midrule
质量 $m$ & 43\,kg & $I_{xx}$ & 0.41\,kg$\cdot$m$^2$ \\
$I_{yy}=I_{zz}$ & 2.1\,kg$\cdot$m$^2$ & 摩擦系数 $\mu$ & 0.6 \\
最大力矩 $\tau_{\max}$ & 250\,N$\cdot$m & 法向力 $f_{\min}/f_{\max}$ & 10\,/\,666\,N \\
力权重 $\alpha$ & $10^{-6}$ & 求解频率 & 20--30\,Hz \\
\bottomrule
\end{tabular}
\end{table}

\paragraph{力到力矩的下层映射。}
MPC 算出地面反力 $\mathbf{f}_i$ 后，站立腿关节力矩由 $\boldsymbol{\tau}_i=\mathbf{J}_i^\top\mathbf{R}^\top\mathbf{f}_i$ 给出（$\mathbf{R}$ 体到世界旋转，$\mathbf{J}_i$ 足雅可比）；摆动腿用操作空间 PD（\cref{sec:ctrl-robot-osc}）跟踪落足轨迹。状态估计、摆动规划、腿阻抗在 1\,kHz 运行，MPC 在 20--50\,Hz——正是 \cref{fig:ctrl-legged-arch} 的分层。

\begin{insight}
Cheetah 凸 MPC 是本章多条主线的汇流：(i) 用\emph{单刚体简化模型}（谱系最简层之一）换来线性时变动力学；(ii) 用\emph{重力增广态}（\cref{sec:ctrl-lqr-dp} 仿射技巧）；(iii) 用\emph{凝聚}（\cref{sec:ctrl-mpc-qp}）把 QP 维度压到与状态无关；(iv) 用\emph{摩擦锥金字塔}把接触约束线性化；(v) 上层 MPC 给力、下层 $\boldsymbol{\tau}=\mathbf{J}^\top\mathbf{R}^\top\mathbf{f}$ 与 WBC 分解。理解了前六节，这篇旗舰论文的每一步都不再神秘。
\end{insight}

\begin{practice}
\begin{enumerate}
  \item （推导）由 \cref{eq:ctrl-srbd-rot} 与 $\boldsymbol{\omega}\times(\mathbf{I}\boldsymbol{\omega})\approx\mathbf{0}$ 写出 $\dot{\boldsymbol{\omega}}=\hat{\mathbf{I}}^{-1}\sum_i\mathbf{r}_i^\wedge\mathbf{f}_i$，并说明它为何对决策变量 $\mathbf{f}_i$ 线性（从而可进凸 QP）。
  \item （分析）摩擦锥金字塔 \cref{eq:ctrl-friction-pyramid} 用方形（4 边）近似圆锥。讨论用更多边（6/8 边）的利弊；它是 SOCP 与 QP 之间的哪种取舍？
  \item （跨章综合）把 Cheetah 的姿态动力学 \cref{eq:ctrl-srbd-R} 改用 \cref{ch:lie} 的右扰动 $\mathbf{R}\,\mathrm{Exp}(\delta\boldsymbol{\phi})$ 线性化，对比它与 Euler 角小角近似 \cref{eq:ctrl-euler-rate} 在大姿态下的差异。
\end{enumerate}
\end{practice}

% ════════════════════════════════════════════════════════════════
\section*{本章常见误解汇总}
\begin{table}[H]\centering\small
\begin{tabular}{p{0.46\textwidth} p{0.46\textwidth}}
\toprule
\textbf{常见误解} & \textbf{正确理解} \\
\midrule
积分项消除一切误差 & 只消\emph{稳态}误差，且须\emph{存在}稳态（\cref{thm:ctrl-integral}）；瞬态仍需整定 \\
限幅输出就等于抗饱和 & 须让\emph{积分器本身}感知饱和（回算 \cref{eq:ctrl-antiwindup}），否则照样 windup \\
渐近稳定 = 最终回到平衡 & 还须先“Lyapunov 稳定”（不中途跑远），二者缺一不可（\cref{def:ctrl-asymp-stable}） \\
$\mathbf{A}$ 稳定就有 $\mathbf{A}+\mathbf{A}^\top\prec0$ & 错；须对给定 $\mathbf{Q}$ 解 Lyapunov 方程得 $\mathbf{M}$（\cref{eq:ctrl-lyap-int}） \\
LaSalle 可用于任何系统 & 仅\emph{自治}系统；非自治用 Barbalat（\cref{lem:ctrl-barbalat}） \\
可控 = 能任意快地驱动 & 只保证有限时间到达；速度/代价是另一回事（受饱和约束） \\
最优控制律必使闭环稳定 & 错；有限时域最优可发散（Kalman 反例 \cref{ex:ctrl-kalman}），须无限时域 \\
LQR 的 $\mathbf{P}$ 就是协方差 & 是\emph{值函数权重}（cost-to-go Hessian），与 Kalman 协方差对偶但语义相反 \\
LQR 有好裕度 $\Rightarrow$ LQG 也有 & 错；状态换成 Kalman 估计后裕度保证失效（Doyle），\cref{sec:ctrl-margins} \\
观测器极点越快越好 & 太快放大测量噪声；Kalman 用噪声协方差最优折中（\cref{eq:ctrl-obs-error}） \\
iLQR 与 DDP 是两套不同算法 & 只差动力学二阶项：DDP 留、iLQR 丢（\cref{eq:ctrl-ilqr-Q}），后者更鲁棒 \\
MPC 时域 $N$ 越大越稳定 & 稳定靠终端代价+终端集（\cref{def:ctrl-mpc-asm}），不靠堆 $N$ \\
计算力矩与 PD+重力补偿二选一 & 分工不同：前者跟踪需精确模型，后者定位只需重力且鲁棒 \\
接触力大就把位置环调软 & 应显式设阻抗（刚度 $\mathbf{K}_d^x$，\cref{eq:ctrl-impedance}）或混合力/位控 \\
WBC 是“一个大 PID” & 是带约束的联合优化 QP，硬性满足动力学/接触/摩擦锥 \\
摩擦锥就是线性约束 & 真实是二阶锥（圆锥）；金字塔化才成线性（为进 QP） \\
\bottomrule
\end{tabular}
\end{table}

% ════════════════════════════════════════════════════════════════
\section*{本章小结}

\begin{table}[H]\centering\small
\caption{符号表}
\label{tab:ctrl-symbols}
\begin{tabular}{lll}
\toprule
符号 & 含义 & 首次出现 \\
\midrule
$\mathbf{A},\mathbf{B},\mathbf{C},\mathbf{D}$ & 状态空间矩阵（$\mathbf{A}$=转移，$\mathbf{C}$=观测） & \cref{sec:ctrl-statespace} \\
$\mathbf{Q},\mathbf{R}$ & LQR 代价权重（状态/输入，\emph{非}噪声协方差） & \cref{eq:ctrl-lqr-cost} \\
$\mathbf{P}$ & Riccati 解 / 值函数权重（\emph{非}协方差） & \cref{eq:ctrl-dare} \\
$\mathbf{K}$ & 状态反馈增益，$\mathbf{u}=-\mathbf{K}\mathbf{x}$ & \cref{eq:ctrl-clqr-K} \\
$\mathbf{L}$ & 观测器（Luenberger）增益 & \cref{eq:ctrl-luenberger} \\
$S,T,\mathrm{GM},\mathrm{PM}$ & 灵敏度/补灵敏度/增益裕度/相位裕度 & \cref{eq:ctrl-sens,def:ctrl-margins} \\
$M_p,t_s,\zeta,\omega_n$ & 超调/调节时间/阻尼比/自然频率 & \cref{eq:ctrl-2nd-order-spec} \\
$k_p,k_i,k_d$ & PID 比例/积分/微分增益 & \cref{eq:ctrl-pid-two} \\
$T_i,T_d,T_f$ & 积分/微分/滤波时间常数 & \cref{eq:ctrl-pid-two,eq:ctrl-filt-deriv} \\
$V(\mathbf{x})$ & Lyapunov 函数 & \cref{thm:ctrl-lyap-direct} \\
$\mathcal{C},\mathcal{O}$ & 可控性/可观性矩阵 & \cref{eq:ctrl-cmat,eq:ctrl-omat} \\
$Q_{\mathbf{u}\mathbf{u}},\mathbf{k},\mathbf{K}$ & iLQR/DDP 的 Q-Hessian/前馈/反馈增益 & \cref{eq:ctrl-ilqr-gain} \\
$\mathcal{X},\mathcal{U},\mathcal{X}_f$ & MPC 状态/输入/终端约束集 & \cref{eq:ctrl-mpc-constr} \\
$V_f,\ell,\kappa_N$ & MPC 终端代价/阶段代价/滚动律 & \cref{eq:ctrl-mpc-cost} \\
$\mathbf{M},\mathbf{C},\mathbf{g}$ & 机器人惯性/科氏/重力 & \cref{eq:ctrl-manip} \\
$\boldsymbol{\Lambda}$ & 操作空间惯性矩阵 $(\mathbf{J}\mathbf{M}^{-1}\mathbf{J}^\top)^{-1}$ & \cref{eq:ctrl-osc} \\
$\mathbf{M}_d,\mathbf{D}_d,\mathbf{K}_d^x$ & 阻抗控制期望惯性/阻尼/刚度 & \cref{eq:ctrl-impedance} \\
$\boldsymbol{\nu},\boldsymbol{\lambda},\mathbf{S},\mathbf{J}_c$ & 广义速度/接触力/选择矩阵/接触雅可比 & \cref{eq:ctrl-floatbase} \\
\bottomrule
\end{tabular}
\end{table}

\begin{table}[H]\centering\small
\caption{定理速查表}
\label{tab:ctrl-thm}
\begin{tabular}{lll}
\toprule
定理 / 公式 & 一句话说明（能做什么） & 对应 \\
\midrule
积分作用消差 & 只要有稳态，积分项保证稳态误差为零（不需模型） & \cref{thm:ctrl-integral} \\
Lyapunov 直接法 & 找“能量”函数判稳定，不必解微分方程 & \cref{thm:ctrl-lyap-direct} \\
LaSalle 原理 & $\dot V$ 只半负定时也能判渐近稳定 & \cref{thm:ctrl-lasalle} \\
Lyapunov 方程 & 给 $\mathbf{Q}$ 解 $\mathbf{M}$，判线性系统稳定 & \cref{eq:ctrl-lyap-eq} \\
Kalman 秩判据 & 判系统能否被驱动到任意状态 & \cref{eq:ctrl-cmat} \\
Luenberger 观测器 & 模型 + 残差纠错重构状态，误差极点可配 & \cref{eq:ctrl-obs-error} \\
分离原理 & 控制器与观测器极点可独立设计（LQG 基础） & \cref{prop:ctrl-separation} \\
离散 LQR Riccati & 后向递推得最优反馈增益 $\mathbf{K}_t$ & \cref{thm:ctrl-lqr-dp} \\
CARE / DARE & 解出无限时域最优 $\mathbf{P}$，给 $\mathbf{K}=\mathbf{R}^{-1}\mathbf{B}^\top\mathbf{P}$ & \cref{eq:ctrl-care,eq:ctrl-dare} \\
Kalman 反例 & 警示“有限时域最优 $\ne$ 稳定” & \cref{ex:ctrl-kalman} \\
iLQR / DDP & 沿轨迹迭代解 LQR 子问题，非线性轨迹优化 & \cref{eq:ctrl-ilqr-gain} \\
MPC 稳定性主定理 & 终端代价+终端集 $\Rightarrow$ 值函数=闭环 Lyapunov & \cref{thm:ctrl-mpc-main} \\
计算力矩 & 抵消非线性，得线性解耦误差动态 & \cref{thm:ctrl-ctc} \\
PD+重力补偿 & 只需重力即可全局定位（无源性+LaSalle） & \cref{thm:ctrl-pd-grav} \\
阻抗控制 & 设末端力--运动关系（刚度可调的柔顺） & \cref{eq:ctrl-impedance} \\
Cheetah 凸 MPC & 单刚体线性化 $\Rightarrow$ 凸 QP 规划接触力 & \cref{eq:ctrl-cheetah-ss} \\
\bottomrule
\end{tabular}
\end{table}

\begin{table}[H]\centering\small
\caption{知识点总表}
\label{tab:ctrl-points}
\begin{tabular}{cllcl}
\toprule
\# & 知识点 & 核心要点 & 难度 & 脉络层 \\
\midrule
1 & 反馈 + 稳定性 & 开环/闭环、性能指标、裕度/根轨迹、Lyapunov/LaSalle & \(\bigstar\bigstar\) & 骨干 \\
2 & PID & 三项作用、滤波微分、抗饱和、Routh & \(\bigstar\bigstar\) & 骨干 \\
3 & 状态空间 & 可控/可观、PBH、对偶、观测器、分离原理 & \(\bigstar\bigstar\bigstar\) & 骨干 \\
4 & LQR / Riccati & DP/HJB 两推导、CARE/DARE、Kalman 反例、iLQR/DDP & \(\bigstar\bigstar\bigstar\) & 骨干 \\
5 & MPC & 滚动时域、凝聚 QP、终端配料保稳定 & \(\bigstar\bigstar\bigstar\bigstar\) & 次优 \\
6 & 机器人控制 & 计算力矩、PD+重力补偿、操作空间、阻抗、WBC & \(\bigstar\bigstar\bigstar\) & 次优 \\
7 & 足式控制 & 分层架构、单刚体凸 MPC、摩擦锥 & \(\bigstar\bigstar\bigstar\bigstar\) & 前沿 \\
\bottomrule
\end{tabular}
\end{table}

% ════════════════════════════════════════════════════════════════
\section*{延伸阅读}
\begin{itemize}
  \item \textbf{经典控制与 PID}：Åström \& Murray《Feedback Systems》\cite{astrom2008feedback}——反馈、灵敏度、PID 整定与抗饱和的权威且易读的来源（本章第 1--2 节主要出处）。
  \item \textbf{非线性系统与 Lyapunov}：Khalil《Nonlinear Systems》\cite{khalil2002nonlinear}——Lyapunov 稳定性、LaSalle、Barbalat、反馈线性化的标准研究生教材。
  \item \textbf{最优控制与 LQR}：Tedrake《Underactuated Robotics》\cite{tedrake_underactuated}（在线，含 HJB$\to$Riccati、LQR 裕度与轨迹 LQR/iLQR）；Bansal LQR 讲义\cite{bansal2017lqr}（离散 DP 完整归纳证明）；Anderson \& Moore《Optimal Control: Linear Quadratic Methods》\cite{anderson2007optimal}（LQR/LQG、保证裕度、Riccati 的经典权威）。
  \item \textbf{iLQR / DDP}：Tassa, Mansard \& Todorov《Control-limited differential dynamic programming》\cite{tassa2014ddp}（控制限 DDP，机器人轨迹优化标准参考；含 iLQR/DDP 后向-前向与正则）。
  \item \textbf{MPC}：Rawlings, Mayne \& Diehl《Model Predictive Control》\cite{rawlings2020mpc}（稳定性理论、Lemma 1.3、数值结构的权威主源）；Borrelli, Bemporad \& Morari《Predictive Control for Linear and Hybrid Systems》\cite{borrelli2017predictive}（持久可行性、Theorem 12.2）；Mayne 等 2000\cite{mayne2000constrained}（MPC 稳定性奠基综述）。
  \item \textbf{机器人控制与交互}：Spong, Hutchinson \& Vidyasagar《Robot Modeling and Control》\cite{spong2006robot}（计算力矩、PD+重力补偿、无源性证）；Khatib 1987\cite{khatib1987operational}（操作空间公式）；Takegaki \& Arimoto 1981\cite{takegaki1981feedback}（PD+重力补偿里程碑）；Hogan 1985\cite{hogan1985impedance}（阻抗控制三部曲，力/柔顺控制奠基）。
  \item \textbf{足式控制}：Di Carlo 等 2018\cite{dicarlo2018cheetah}（MIT Cheetah 3 凸 MPC，单刚体模型旗舰论文）；Wensing 等 2024\cite{wensing2024legged}（基于优化的足式控制综述：WBC、分层 QP、质心动力学）。
\end{itemize}

\section*{本章与后续章节的关系}
\begin{table}[H]\centering\small
\begin{tabular}{lll}
\toprule
关联章节 & 关系 & 本章铺垫 \\
\midrule
\cref{ch:eskf} 卡尔曼/ESKF & 估计--控制对偶：Kalman Riccati $\leftrightarrow$ LQR Riccati & \cref{sec:ctrl-lqe-dual} \\
\cref{ch:nlopt} 非线性优化 & MPC 的 QP/LQR 的二次最小化同源 & \cref{sec:ctrl-mpc-qp} \\
\cref{ch:lie} 李群与李代数 & 任务空间姿态误差、单刚体姿态动力学 & \cref{sec:ctrl-robot-osc,sec:ctrl-legged-cheetah} \\
\cref{ch:planning} 规划导论 & 规划给参考轨迹，控制（MPC/WBC）跟踪之 & \cref{sec:ctrl-mpc,sec:ctrl-legged} \\
\bottomrule
\end{tabular}
\end{table}

\section*{故障排查手册}
\begin{table}[H]\centering\small
\begin{tabular}{p{0.22\textwidth} p{0.26\textwidth} p{0.40\textwidth}}
\toprule
症状 & 可能原因 & 排查步骤 \\
\midrule
改设定点后控制量猛冲、执行器抖 & 微分踢（对参考微分） & 改为只对测量微分（\cref{eq:ctrl-setpoint-weight} 取 $\gamma=0$）+ 滤波微分 \cref{eq:ctrl-filt-deriv} \\
大扰动后长时间大超调 & 积分饱和（windup） & 加回算抗饱和 \cref{eq:ctrl-antiwindup} 或改增量式；查执行器是否常饱和 \\
LQR 闭环仿真发散 & 用了有限时域首增益或未验稳定 & 解 DARE/CARE 至收敛，验 $\operatorname{eig}(\mathbf{A}-\mathbf{B}\mathbf{K})$（\cref{ex:ctrl-kalman}） \\
LQG 实机裕度差/易振 & 把 LQR 的保证裕度误当 LQG 也有 & LQG 裕度无保证（\cref{sec:ctrl-margins} note）；用 LQG/LTR 或调观测器带宽 \\
观测器估计跳变/噪声大 & 观测器极点配得过快（$\mathbf{L}$ 过大） & 放慢误差极点（\cref{eq:ctrl-obs-error}），或改用 Kalman 按噪声协方差折中 \\
iLQR/DDP 不收敛、代价上升 & $Q_{\mathbf{u}\mathbf{u}}$ 非正定 / 步长过大 & 加正则 $Q_{\mathbf{u}\mathbf{u}}+\mu\mathbf{I}$ + 前向线搜索缩 $\alpha$（\cref{alg:ctrl-ilqr}） \\
MPC 输出 NaN / 求解失败 & 初始不可行 / 无终端集 & 查 $\mathbf{x}(0)$ 可行性；加控制不变终端集 $\mathcal{X}_f$（\cref{lem:ctrl-pf2}）；降时域 \\
MPC 跑着跑着不可行 & 持久可行性缺失 & 用 \cref{lem:ctrl-pf2} 加控制不变 $\mathcal{X}_f$；检查约束是否过紧 \\
计算力矩跟踪误差大 & 模型 $\mathbf{M},\mathbf{C},\mathbf{g}$ 不准 & 标定动力学参数；或退回 PD+重力补偿（对模型误差鲁棒，\cref{tab:ctrl-ctc-vs-pd}） \\
接触时力尖峰/卡死 & 用纯刚性位置控制顶接触 & 改阻抗控制设刚度 $\mathbf{K}_d^x$（\cref{eq:ctrl-impedance}）或混合力/位控 \\
足式 WBC 打滑/腾空 & 接触力出摩擦锥 / 动力学不一致 & 在 QP 加摩擦锥 \cref{eq:ctrl-friction-pyramid} 与动力学等式约束 \cref{eq:ctrl-wbc-qp} \\
\bottomrule
\end{tabular}
\end{table}

\section*{研究实践建议}
\begin{itemize}
  \item \textbf{新手}：从 PID + 双积分器（\cref{sec:ctrl-statespace} 练习）入手，先吃透“稳态误差—积分—饱和”这条线；再用 SciPy \texttt{solve\_continuous\_are} 跑一个简单 LQR，亲手验闭环极点。
  \item \textbf{有经验者}：把 LQR 当 MPC 的特例理解（\cref{ex:ctrl-mpc-eq-lqr}），用 OSQP 实现一个带约束的线性 MPC，体会终端集对持久可行性的作用；机器人方向可复现 \cref{sec:ctrl-legged-cheetah} 的单刚体凸 MPC。
  \item \textbf{统一视角}：始终记住三个“同源”——LQR=无约束无限时域 MPC、Kalman Riccati 对偶 LQR Riccati、计算力矩把机器人化为双积分器后接 LQR/MPC。这些连接比记住每个公式更重要。
\end{itemize}

% ════════════════════════════════════════════════════════════════
\section{本章推导附录}
\label{sec:ctrl-appendix}

本章主要证明（积分消差、Lyapunov 直接法/LaSalle、Lyapunov 方程、可控 Gramian 双向、分离原理、LQR 的 DP/HJB 推导、iLQR/DDP 后向递推、MPC 单调性与主定理、计算力矩与 PD+重力补偿）均已在正文相应定理处给出。此处补全 Joseph 形式与紧凑形式的等价代数细节。

\begin{derivation}[离散 Riccati：Joseph 形式 $\equiv$ 紧凑形式]
\cref{thm:ctrl-lqr-dp} 中 Joseph 形式 \cref{eq:ctrl-dlqr-joseph}
$\mathbf{P}_{k-1}=\mathbf{Q}+\mathbf{K}_{k-1}^\top\mathbf{R}\mathbf{K}_{k-1}+(\mathbf{A}-\mathbf{B}\mathbf{K}_{k-1})^\top\mathbf{P}_k(\mathbf{A}-\mathbf{B}\mathbf{K}_{k-1})$ 与紧凑形式 \cref{eq:ctrl-dlqr-P} 代数恒等。记 $\mathbf{G}=\mathbf{R}+\mathbf{B}^\top\mathbf{P}_k\mathbf{B}$，$\mathbf{K}_{k-1}=\mathbf{G}^{-1}\mathbf{B}^\top\mathbf{P}_k\mathbf{A}$。展开 Joseph 形式：
\begin{align*}
  \mathbf{P}_{k-1}&=\mathbf{Q}+\mathbf{A}^\top\mathbf{P}_k\mathbf{A}-\mathbf{A}^\top\mathbf{P}_k\mathbf{B}\mathbf{K}_{k-1}-\mathbf{K}_{k-1}^\top\mathbf{B}^\top\mathbf{P}_k\mathbf{A}
   +\mathbf{K}_{k-1}^\top\underbrace{(\mathbf{R}+\mathbf{B}^\top\mathbf{P}_k\mathbf{B})}_{=\mathbf{G}}\mathbf{K}_{k-1}.
\end{align*}
代入 $\mathbf{K}_{k-1}=\mathbf{G}^{-1}\mathbf{B}^\top\mathbf{P}_k\mathbf{A}$，则 $\mathbf{K}_{k-1}^\top\mathbf{G}\mathbf{K}_{k-1}=\mathbf{A}^\top\mathbf{P}_k\mathbf{B}\mathbf{G}^{-1}\mathbf{B}^\top\mathbf{P}_k\mathbf{A}=\mathbf{K}_{k-1}^\top\mathbf{B}^\top\mathbf{P}_k\mathbf{A}$，恰与第三项相消，且第二项 $\mathbf{A}^\top\mathbf{P}_k\mathbf{B}\mathbf{K}_{k-1}=\mathbf{A}^\top\mathbf{P}_k\mathbf{B}\mathbf{G}^{-1}\mathbf{B}^\top\mathbf{P}_k\mathbf{A}$。合并得
\[
  \mathbf{P}_{k-1}=\mathbf{Q}+\mathbf{A}^\top\mathbf{P}_k\mathbf{A}-\mathbf{A}^\top\mathbf{P}_k\mathbf{B}(\mathbf{R}+\mathbf{B}^\top\mathbf{P}_k\mathbf{B})^{-1}\mathbf{B}^\top\mathbf{P}_k\mathbf{A},
\]
即紧凑形式 \cref{eq:ctrl-dlqr-P}。Joseph 形式数值上更稳（自动保对称半正定），紧凑式更省算——这正是 \cref{sec:ctrl-mpc-numerical} 选择 Riccati 因子分解的依据。
\end{derivation}
